\documentclass[11pt]{article}
\usepackage[margin=2.1cm]{geometry}
\usepackage{amsmath,amssymb,amsthm,mathrsfs}
\usepackage{longtable,array}
\usepackage{url}
\usepackage[hidelinks]{hyperref}
\numberwithin{equation}{section}
\theoremstyle{plain}
\newtheorem{theorem}{Theorem}[section]
\newtheorem{proposition}[theorem]{Proposition}
\newtheorem{lemma}[theorem]{Lemma}
\newtheorem{corollary}[theorem]{Corollary}
\theoremstyle{definition}
\newtheorem{definition}[theorem]{Definition}
\newtheorem{example}[theorem]{Example}
\newtheorem{remark}[theorem]{Remark}
\newtheorem{hypothesis}[theorem]{Hypothesis}
\allowdisplaybreaks
\emergencystretch=3em
\newcommand{\Ap}{A|_{\perp}}
\newcommand{\AV}{A|_V}
\newcommand{\Ob}{\mathrm{Ob}}
\newcommand{\ICT}{I^{CT}[\varepsilon]}
\newcommand{\lc}{\mathrm{lc}}
\newcommand{\dslash}{{\not d}}

\title{The action-level one-loop equivariant quantum master equation\\ along every generator of $\mathfrak{psl}(4|4)$}
\author{Step 2, strict Costello route, action side (theory)}
\date{2026-10-30}

\begin{document}
\maketitle

\begin{abstract}
For $\mathcal N=4$ super Yang--Mills theory in the formalism of [C07] on Euclidean $\mathbb R^4$, with the heat-kernel cutoff $\varepsilon$ on every edge and the transfer endpoint $T$, we set up and analyse the one-loop equivariant quantum master equation of Definition c8-quant at the level of the action, arity by arity, along every generator: the dilatation $\Delta$, the supercharges $\xi\in S_+\otimes W$ and $\xi\in S_-\otimes W^*$, the special supercharges $\eta\in\Sigma^*$, the special conformal fields $K_b$, and, as regressions, the translations $P_b$, the rotations $\mathfrak{so}(V)$ and the R-symmetry $\mathfrak{so}(W^\perp_{\mathbb R})$. The obstruction $\Ob$ of Definition c8-ob is written component by component as the $\varepsilon\to0$ limit of the defect graphs. We prove: the one-loop counterterm cochain is $\tfrac94$ units times $\varepsilon^{-1}\int(A,A)$ plus $\log\varepsilon$ times an $s$-closed cochain (no pole in arity one or two); every component $\Ob^{(1)}(x)$ and $\Ob^{(2)}(x,y)$ is the sum of the $\varepsilon^0$ coefficients of the irreducible defect cycles with one special edge, the vertex $S^{(1)}(x)$ (and $S^{(1)}(y)$ or $S^{(2)}(x,y)$) placed on the cycle in the listed ways, all other classes cancelling or converging to trees; a second, independent evaluation applies $Q$ to the renormalised arity-one and arity-two amplitudes, which the conjugation of $S^{(1)}(x)$ by $e^{\hbar\partial_P}$ (terminating at second order along every generator, the cubic term of $S^{(1)}(\xi)$ included) organises into variation kernels, composite lines and endpoint terms. The census of the sectors that weight, ghost number, the centres of $\mathrm{Spin}(4)\times SU(4)$ and colour allow is given for all components. The Wess--Zumino identities are derived in every arity. The arity-one equation is solvable exactly when a residual class in $H^0(Q+\{I,-\})$ vanishes; along $\Sigma$ the two-leg cohomology is zero, so the two-leg equation is solvable with the two free directions $S^{(1)}(\xi)$ restricted to the $V$-arrows and to the $W^\perp_{\mathbb R}$-arrows. The cubic part of the arity-one equation fixes the differences of the sector coefficients exactly when the sector parts of $S^{(1)}(x)$ fail to extend to symmetries of the interacting classical theory. The arity-two residuals form a two-cocycle of $\mathfrak{psl}(4|4)$ with values in $H^{-1}(Q+\{I,-\})$; under the hypothesis that the classical symmetries of local degree $-1$ are the superconformal ones, the arity-two equations fix the uniform coefficients $t_x$ of $S^{(1)}(x)$ in $\hbar F_1^{(1)}(x)$ up to the inner direction generated by $S^{(1)}(\Delta)$ and the outer direction $G_B$; in particular the coefficient $\theta$ of Theorem c9-kappa is fixed by the two-leg part of the arity-two equation along $(\Delta,P_b)$, $2u\,\theta_d$ equal to the coefficient of the class of $S^{(1)}(P_b)_d$ in the arity-two obstruction. The one-loop corrections to $[Q,Q]$, $\{Q,S\}$, $[K,P]$, $[D,Q]$ are the components of this cocycle. Closed forms (simplex integrals with the special edge, generating functions, momentum injection for the position) and the specification of two independent engine routes with their gates complete the note; the scaling identity, $K=P^\dagger$, the twist and the value $\theta=-\tfrac13$ enter as cross-checks only.
\end{abstract}

\tableofcontents

\section{Data, names and inputs}
\label{sec:at-data}

The source of truth is the manuscript, Chapters 1--11; labels such as Definition c8-ob refer to it. The unit is
\[
u:=\frac{\hbar\operatorname{tr}_{\mathfrak g}(\mathrm{ad}_X\mathrm{ad}_Y)}{16\pi^2(X,Y)}=-\frac{g^2N}{32\pi^2}\quad(\mathfrak g=\mathfrak{su}(N),\ (X,Y)=-\tfrac{4\hbar}{g^2}\operatorname{tr}(XY),\ \text{Proposition c8-twocomp}).
\]
Edges carry $P(\varepsilon,T)=\int_\varepsilon^T(Q^{GF}\otimes1)K_t\,dt$ (Definition c8-PT); $T=\infty$ along every generator except $K_b$, and along the pairs $(K_b,P_{b'})$ and $(\eta,\eta')$ with $\Gamma(\eta,\eta')|_V\neq0$, where $T$ is finite (Proposition c8-Tsafe, Definition c8-K). No dimensional regularisation and no infrared regulator other than $T$ is used anywhere.

\begin{definition}[Generators and their data]
\label{def:at-gens}
For a homogeneous generator $x$ of $\mathfrak{psl}(4|4)$ let $\lambda_x$ be its $\mathrm{ad}_\Delta$-eigenvalue ($[\Delta,x]=\lambda_xx$, Remark c2-dilatation), $p_x$ the polynomial degree of its Killing data in the position, and $|x|$ its parity. The functional $S^{(1)}(x)$ (Definition c2-S1, \eqref{eq:at-S1list} below) has ghost degree $-1$ and Weyl weight $4-\tfrac12\lambda_x$, the position counted at weight $-1$.
\[
\begin{array}{l|c|c|c|l|l}
x & \lambda_x & p_x & |x| & \text{terms of }S^{(1)}(x) & \text{endpoint}\\
\hline
\Delta\ (v=2E) & 0 & 1 & 0 & \text{quadratic: }(v.A,A^*),\ -(v.\lambda,\lambda^*),\ -(v.c,c^*) & T=\infty\\
P_b & -2 & 0 & 0 & \text{quadratic} & T=\infty\\
\mathfrak{so}(V)\ni M_\omega\ (v=\omega x) & 0 & 1 & 0 & \text{quadratic} & T=\infty\\
\mathfrak{so}(W^\perp_{\mathbb R})\ni R_r & 0 & 0 & 0 & \text{quadratic, on }\Ap,\lambda\text{ only} & T=\infty\\
\xi\in S_+\otimes W\ \text{or}\ S_-\otimes W^* & -1 & 0 & 1 & T_1,\dots,T_4\ \text{quadratic},\ T_5\ \text{cubic} & T=\infty\\
\eta\in S_+\otimes W^*\ \text{or}\ S_-\otimes W & 1 & 1 & 1 & \textstyle\sum_\mu x^\mu S^{(1)}(\rho(e_\mu)\eta)+\int s_0 & T=\infty\\
K_b\ (v=(x,x)b-2(b,x)x) & 2 & 2 & 0 & \text{quadratic} & T<\infty
\end{array}
\]
Here, with $T_1,\dots,T_5$ the terms of Table c11-terms and $s_0=-(\rho(\Ap)\eta,\lambda^*)$ ((c11-Seta) of Chapter 11),
\begin{equation}
\begin{aligned}
&S^{(1)}(x)=\int\mathrm{dvol}\,\big((v.A,A^*)-(v.\lambda,\lambda^*)-(v.c,c^*)\big)\quad(x\ \text{even}),\\
&S^{(1)}(\xi)=T_1+T_2+T_3+T_4+T_5,\qquad T_1=-\int(\Gamma(\xi,\lambda)|_V,A^*|_V),\ T_2=-\int(\Gamma(\xi,\lambda)|_\perp,A^*|_\perp),\\
&T_3=\tfrac12\int(\rho(d\AV)\xi,\lambda^*),\ T_4=\tfrac12\int(\rho(d\Ap)\xi,\lambda^*),\ T_5=\tfrac12\int(\rho(\tfrac12[A\wedge A])\xi,\lambda^*).
\end{aligned}
\label{eq:at-S1list}
\end{equation}
The leg names used throughout: $\lambda_+=\lambda|_{S_+\otimes W}$, $\lambda_-=\lambda|_{S_-\otimes W^*}$, $\lambda^*_+=\lambda^*|_{S_+\otimes W^*}$, $\lambda^*_-=\lambda^*|_{S_-\otimes W}$, $\xi_+\in S_+\otimes W$, $\xi_-\in S_-\otimes W^*$, $\eta_+\in S_+\otimes W^*$, $\eta_-\in S_-\otimes W$.
\end{definition}

\begin{definition}[Components, sector parts, coefficients]
\label{def:at-names}
\leavevmode
\begin{itemize}
\item (i) A cochain $f$ of $C^\bullet(\mathfrak{psl}(4|4),\mathscr O_l(\mathscr E))$ has components $f^{(k)}(x_1,\dots,x_k)$ (Definition c2-cochains), normalised as in Lemma c2-components: the arity-one component of $d_{CE}f+\tfrac12\{f,f\}$ at $f=S_{\mathfrak{psl}(4|4)}$ is $\{S,S^{(1)}(x)\}$ and its arity-two component is $\{S,S^{(2)}(x,y)\}+S^{(1)}([x,y])+(-1)^{|x|}\{S^{(1)}(x),S^{(1)}(y)\}$. $\Lambda^{(0)}:=Q+\{I,-\}=\{S,-\}$ on local functionals (Theorem c8-L (iv)). $\Ob^{(k)}$ and $F_1^{(k)}$ are the components of the obstruction and of the finite counterterm $\hbar F_1$ in this normalisation.
\item (ii) Sector parts. For even $x$ the sector parts of $S^{(1)}(x)$ are
\[
S^{(1)}(x)_\perp=\int(v.\Ap,A^*|_\perp),\qquad S^{(1)}(x)_\lambda=-\int(v.\lambda,\lambda^*),\qquad S^{(1)}(x)_V=\int(v.\AV,A^*|_V)-\int(v.c,c^*),
\]
with sum $S^{(1)}(x)$; each is $Q$-closed, $v.$ commuting with $Q$ (Propositions c7-cov, c7-dil). For $\xi\in\Sigma$, $S^V_2(\xi):=T_1+T_3$ and $S^\perp_2(\xi):=T_2+T_4$ (Definition c11-terms); for $\eta\in\Sigma^*$, $S^V(\eta):=\sum_\mu\int x^\mu S^V_2(\rho(e_\mu)\eta)$ and $S^\perp(\eta):=\sum_\mu\int x^\mu S^\perp_2(\rho(e_\mu)\eta)+\int s_0$, whose sum is the quadratic part of $S^{(1)}(\eta)$. The index $d$ runs over the sectors $\{\perp,\lambda,V\}$ (even $x$) or $\{V,\perp\}$ (odd $x$).
\item (iii) The quadratic part of $\hbar F_1^{(1)}(x)$ contains $u\sum_d\theta^x_d$ times the sector parts; $\theta^\Delta_\perp,\theta^\Delta_\lambda,\theta^\Delta_V$ are $\theta,\theta_\lambda,\theta_V$ of Definition c11-free and $\theta^{P_b}_d$ are $\theta_P,\theta_{P,\lambda},\theta_{P,V}$; along $\xi$, $\theta^\xi_V,\theta^\xi_\perp$ are those of Proposition qs-freexi. When all $\theta^x_d$ are equal their common value is the uniform coefficient $t_x$, the coefficient of $u\,S^{(1)}(x)$ (completed by its cubic part).
\item (iv) $\ICT$ is the counterterm cochain $I^{CT}_{\mathfrak{psl}(4|4)}[\varepsilon]$ of Definition c8-quant; its coefficient of $\log\varepsilon$ is the cochain $\lc$, and $\lc^{(k)}$ its components.
\end{itemize}
\end{definition}

\begin{remark}[Physics notation]
\label{rem:at-phys}
$S_+$ carries undotted and $S_-$ dotted indices; $W$ and $W^*$ are the $\mathbf 4$ and $\bar{\mathbf 4}$ of $SU(4)_R$ (upper and lower index $i$), $W^\perp_{\mathbb R}$ the vector $\mathbf 6$ of $SO(6)_R$, realised by the antisymmetric $A_{ij}=A_I\Sigma^I_{ij}$ (Definition c8-twocomp); $\rho$ is Clifford multiplication, in components the ten-dimensional chiral $\Gamma^M$, in four dimensions $\sigma^m,\bar\sigma^m$, and on the Dirac column $\rho(v)=-i\gamma_S^\mu v_\mu=\gamma_W^\mu v_\mu$ (Srednicki, Weinberg). The gaugino $\lambda$ lies in $\Pi\Sigma\otimes\mathfrak g$, $\Pi$ the parity shift, $\Sigma=S_+\otimes W\oplus S_-\otimes W^*$ the ten-dimensional chiral spinor $\mathbf{16}$, $\Sigma^*$ the $\mathbf{16}'$ (lower ten-dimensional index). Functional degrees (ghost numbers of the coordinate functionals): $c$: $+1$, $A$, $\lambda$: $0$, $A^*$, $\lambda^*$: $-1$, $c^*$: $-2$; Weyl weights $0,1,\tfrac32,3,\tfrac52,4$. For $\mathfrak g=\mathfrak{su}(N)$, $A_m=-igA^a_mT^a$, $\lambda=-\tfrac{ig}{\sqrt2}\lambda^aT^a$, $\operatorname{tr}T^aT^b=\tfrac12\delta^{ab}$ (Proposition c8-twocomp).
\end{remark}

\begin{remark}[Inputs used as established]
\label{rem:at-inputs}
From the manuscript: the arity-zero results (Theorems c4-quadratic, c4-log, c4-obstruction, c4-finite, c9-ob3, c9-F3, c9-wz), the two-leg data along $\Delta$ of Theorem c9-quadratic and Remark c9-Pi as regression targets (their derivation in Chapter 9 uses the scaling identity; here they are to be recomputed through the defect graphs), the commutators of Lemma c11-lines and Proposition qs-comm, the simplex theorems c10-simplex, c11-simplex, c11-oddsimplex, and the endpoint statements of Chapter 8. From the earlier engines: their sign rules and normalisations (Section \ref{sec:at-engines}).
\end{remark}

\section{The arity-one obstruction: definition, graphs, bounds}
\label{sec:at-ob1}

\subsection{The defect vertex by arity and the Wess--Zumino identities}

Write $\{f,g\}_\varepsilon-\{f,g\}$ for the special-edge part of the scale-$\varepsilon$ bracket: the two functionals joined by one edge $K_\varepsilon-K_0=-\int_0^\varepsilon(H\otimes1)K_t\,dt$, in momentum space $e^{-\varepsilon|p|^2}-1$ times the components of $K_0$ (Proposition c4-defect1).

\begin{lemma}[Components of the defect vertex]
\label{lem:at-defect}
Let $J=I_{\mathfrak{psl}(4|4)}-\ICT-\hbar F_1$ and let $D:=(d_{CE}+Q)J+\tfrac12\{J,J\}_\varepsilon+\hbar\Delta_\varepsilon J$ be the defect vertex of Definition c8-ob. Modulo $\hbar^2$ and constants, in the normalisation of Definition \ref{def:at-names} (i), with $S(x):=S^{(1)}(x)$,
\begin{equation}
\begin{aligned}
D^{(0)}={}&\tfrac12\big(\{I,I\}_\varepsilon-\{I,I\}\big)-Q\,\ICT^{(0)}-\{I,\ICT^{(0)}\}_\varepsilon-\hbar\big(QF_1^{(0)}+\{I,F_1^{(0)}\}_\varepsilon\big),\\
D^{(1)}(x)={}&\big(\{I,S(x)\}_\varepsilon-\{I,S(x)\}\big)+\hbar\Delta_\varepsilon S(x)-Q\,\ICT^{(1)}(x)-\{I,\ICT^{(1)}(x)\}_\varepsilon-\{\ICT^{(0)},S(x)\}_\varepsilon\\
&-\hbar\big(QF_1^{(1)}(x)+\{I,F_1^{(1)}(x)\}_\varepsilon+\{F_1^{(0)},S(x)\}_\varepsilon\big),\\
D^{(2)}(x,y)={}&(-1)^{|x|}\big(\{S(x),S(y)\}_\varepsilon-\{S(x),S(y)\}\big)+\big(\{I,S^{(2)}(x,y)\}_\varepsilon-\{I,S^{(2)}(x,y)\}\big)\\
&-\Big[Q\,\ICT^{(2)}(x,y)+\{I,\ICT^{(2)}(x,y)\}_\varepsilon+\ICT^{(1)}([x,y])\\
&\qquad+(-1)^{|x|}\big(\{\ICT^{(1)}(x),S(y)\}_\varepsilon+\{S(x),\ICT^{(1)}(y)\}_\varepsilon\big)+\{\ICT^{(0)},S^{(2)}(x,y)\}_\varepsilon\Big]\\
&-\hbar\big[\text{the same bracket with }F_1\text{ in place of }\ICT\big].
\end{aligned}
\label{eq:at-defect}
\end{equation}
$\hbar\Delta_\varepsilon I=0$, $\hbar\Delta_\varepsilon S(x)$ is a constant for even $x$ and $0$ for odd $x$, and $\hbar\Delta_\varepsilon S^{(2)}(x,y)=0$.
\end{lemma}

\begin{proof}
The order-$\hbar^0$, $\varepsilon$-independent part of $D$ is $d_{CE}I_{\mathfrak{psl}(4|4)}+\tfrac12\{I_{\mathfrak{psl}(4|4)},I_{\mathfrak{psl}(4|4)}\}$, which vanishes by (c2-mc); its arity components are those of Lemma c2-components, and the $\varepsilon$-dependent remainder is the special-edge part of each bracket. The terms linear in $\ICT$ and in $F_1$ are the linearisation of the same component formulas at $S_{\mathfrak{psl}(4|4)}$ (Lemma \ref{lem:at-sformulas} below), with the bracket at scale $\varepsilon$. $\Delta_\varepsilon$ contracts a field with its own antifield inside one term: in $I$ this produces $\operatorname{tr}\mathrm{ad}=0$ (proof of Theorem c4-obstruction); in $S(x)$ for even $x$ the terms $(v.e,e^*)$ give the field-independent supertrace of $v.\,e^{-\varepsilon H}$; the quadratic terms of $S(\xi)$ and of $S(\eta)$ pair $\lambda$ with $A^*$ and $A$ with $\lambda^*$, the cubic term has legs $A,A,\lambda^*$, and $s_0$ has legs $\Ap,\lambda^*$; the terms of $S^{(2)}$ ((c2-S2)) have the legs $(\lambda^*,\lambda^*)$ and $(A,c^*)$.
\end{proof}

\begin{definition}[Arity components of the obstruction; residuals]
\label{def:at-ob}
With $F_1=0$ the one-loop obstruction $\Ob$ is the local cochain of degree $1$ of Definition c8-ob: its tree-level flow with edges $P(0,T)$ is the coefficient of $\hbar$ of the $\varepsilon\to0$ limit of the left side of (c8-qme), equivalently of the connected graphs with one vertex $D$, the others from $J$, edges $P(\varepsilon,T)$. $\Ob^{(0)}$, $\Ob^{(1)}(x)$, $\Ob^{(2)}(x,y)$ are its components in the normalisation of Definition \ref{def:at-names} (i), so that $s(\hbar F_1)=\hbar\Ob$ reads, by Lemma \ref{lem:at-sformulas},
\begin{equation}
\begin{aligned}
&\Lambda^{(0)}F_1^{(0)}=\Ob^{(0)},\qquad \Lambda^{(0)}F_1^{(1)}(x)+\{F_1^{(0)},S(x)\}=\Ob^{(1)}(x),\\
&\Lambda^{(0)}F_1^{(2)}(x,y)+F_1^{(1)}([x,y])+(-1)^{|x|}\big(\{F_1^{(1)}(x),S(y)\}+\{S(x),F_1^{(1)}(y)\}\big)+\{F_1^{(0)},S^{(2)}(x,y)\}=\Ob^{(2)}(x,y).
\end{aligned}
\label{eq:at-eqs}
\end{equation}
Given $F_1^{(0)}$ solving the first equation, the arity-one residual is $\mathcal R^{(1)}(x):=\Ob^{(1)}(x)-\{F_1^{(0)},S(x)\}$; given also $F_1^{(1)}$ solving the second, the arity-two residual is
\[
\mathcal R^{(2)}(x,y):=\Ob^{(2)}(x,y)-F_1^{(1)}([x,y])-(-1)^{|x|}\big(\{F_1^{(1)}(x),S(y)\}+\{S(x),F_1^{(1)}(y)\}\big)-\{F_1^{(0)},S^{(2)}(x,y)\}.
\]
$\Ob^{(1)}(x)$ is local of degree $0$ and Weyl weight $4-\tfrac12\lambda_x$, $\Ob^{(2)}(x,y)$ of degree $-1$ and weight $4-\tfrac12(\lambda_x+\lambda_y)$; $F_1^{(1)}(x)$ and $F_1^{(2)}(x,y)$ have degrees $-1$, $-2$ and the same weights.
\end{definition}

\begin{lemma}[Component formulas of $s$; Wess--Zumino identities]
\label{lem:at-sformulas}
Let $s=d_{CE}+\{S_{\mathfrak{psl}(4|4)},-\}$. In the normalisation of Definition \ref{def:at-names} (i):
\begin{itemize}
\item (i) for a cochain $F$ of degree $0$, $(sF)^{(0)}=\Lambda^{(0)}F^{(0)}$, $(sF)^{(1)}(x)=\Lambda^{(0)}F^{(1)}(x)+\{F^{(0)},S(x)\}$, and $(sF)^{(2)}(x,y)$ is the left side of the third equation of \eqref{eq:at-eqs};
\item (ii) for a cochain $G$ of degree $1$, $(sG)^{(1)}(x)=\Lambda^{(0)}G^{(1)}(x)+\{G^{(0)},S(x)\}$ and $(sG)^{(2)}(x,y)=\Lambda^{(0)}G^{(2)}(x,y)-G^{(1)}([x,y])+(-1)^{|x|}\{G^{(1)}(x),S(y)\}-\{S(x),G^{(1)}(y)\}+\{G^{(0)},S^{(2)}(x,y)\}$, so $sG=0$ in arities $0,1,2$ reads
\[
\Lambda^{(0)}G^{(0)}=0,\qquad \Lambda^{(0)}G^{(1)}(x)=-\{G^{(0)},S(x)\},\qquad \Lambda^{(0)}G^{(2)}(x,y)=G^{(1)}([x,y])-(-1)^{|x|}\{G^{(1)}(x),S(y)\}+\{S(x),G^{(1)}(y)\}-\{G^{(0)},S^{(2)}(x,y)\};
\]
\item (iii) (Wess--Zumino) these hold for $G=\Ob$; consequently $\Lambda^{(0)}\mathcal R^{(1)}(x)=0$ when $\Lambda^{(0)}F_1^{(0)}=\Ob^{(0)}$, and $\Lambda^{(0)}\mathcal R^{(2)}(x,y)=0$ when the equations of arity $0$ and $1$ hold.
\end{itemize}
\end{lemma}

\begin{proof}
(i) Differentiate the component formulas of Lemma c2-components along the family $S_{\mathfrak{psl}(4|4)}+tF$; $\tfrac{d}{dt}\tfrac12\{f,f\}=\{f,\dot f\}$ for even $f$. (ii) Derivation rule. By Definition c1-bracket, $\{f,g\}=\Delta_0(fg)-\Delta_0(f)g-(-1)^{|f|}f\Delta_0(g)=(-1)^{|f|}\{f,g\}_{\rm std}$, with $\{,\}_{\rm std}$ the odd bracket for which $\Delta_0(fg)=\Delta_0(f)g+(-1)^{|f|}f\Delta_0(g)+(-1)^{|f|}\{f,g\}_{\rm std}$. A derivation $\mathcal D$ of degree $1$ of $\{,\}_{\rm std}$ ($\mathcal D=Q$, $\{I,-\}$, $\Lambda^{(0)}$, or $\{F,-\}$ with $F$ of degree $0$) satisfies $\mathcal D\{a,b\}_{\rm std}=\{\mathcal Da,b\}_{\rm std}+(-1)^{|a|+1}\{a,\mathcal Db\}_{\rm std}$, hence
\[
\mathcal D\{a,b\}=-\{\mathcal Da,b\}-(-1)^{|a|}\{a,\mathcal Db\}
\]
for the bracket of Definition c1-bracket. The components of $sG$ are universal expressions $\Lambda^{(0)}G^{(2)}+\alpha G^{(1)}([x,y])+\beta\{G^{(1)}(x),S(y)\}+\gamma\{S(x),G^{(1)}(y)\}+\delta\{G^{(0)},S^{(2)}(x,y)\}$ and $\Lambda^{(0)}G^{(1)}(x)+\delta_1\{G^{(0)},S(x)\}$ with signs fixed by the Koszul rules, and $s^2=0$ forces them on the exact cochains $G=sF$. Arity one, $F=F^{(0)}$: $G^{(0)}=\Lambda^{(0)}F^{(0)}$, $G^{(1)}(x)=\{F^{(0)},S(x)\}$, and $\Lambda^{(0)}\{F^{(0)},S(x)\}=-\{\Lambda^{(0)}F^{(0)},S(x)\}$ ($\Lambda^{(0)}S(x)=0$), so $\delta_1=1$. Arity two, $F=F^{(1)}$: $G^{(1)}=\Lambda^{(0)}F^{(1)}$ and $G^{(2)}=F^{(1)}([x,y])+(-1)^{|x|}(\{F^{(1)}(x),S(y)\}+\{S(x),F^{(1)}(y)\})$; with $|S(x)|\equiv1+|x|$, $\Lambda^{(0)}G^{(2)}=\Lambda^{(0)}F^{(1)}([x,y])-(-1)^{|x|}\{\Lambda^{(0)}F^{(1)}(x),S(y)\}+\{S(x),\Lambda^{(0)}F^{(1)}(y)\}$, so $\alpha=\gamma=-1$, $\beta=(-1)^{|x|}$. Arity two, $F=F^{(0)}$: $G^{(1)}(x)=\{F^{(0)},S(x)\}$, $G^{(2)}=\{F^{(0)},S^{(2)}(x,y)\}$; with $\Lambda^{(0)}S^{(2)}(x,y)=-S([x,y])-(-1)^{|x|}\{S(x),S(y)\}$ (Lemma c2-components) and the derivation rule for $\{F^{(0)},-\}$, $\{F^{(0)},\{S(x),S(y)\}\}=-\{\{F^{(0)},S(x)\},S(y)\}+(-1)^{|x|}\{S(x),\{F^{(0)},S(y)\}\}$, every term cancels except $(\delta-1)\{\Lambda^{(0)}F^{(0)},S^{(2)}(x,y)\}$, so $\delta=1$. (The forms are checked in a toy BV algebra with the bracket of Definition c1-bracket: step2/strict/action/qa/src.nosync/q\_bv\_toy.py.) (iii) $s\Ob=0$ is Theorem c8-obstruction (i); insert $G=\Ob-sF_1$, whose lower components vanish.
\end{proof}

\subsection{Poles and logarithms of the counterterm cochain}

\begin{proposition}[Pole structure]
\label{prop:at-poles}
At one loop,
\[
\ICT=\tfrac94\,u\,\varepsilon^{-1}\int\mathrm{dvol}\,(A,A)+\log\varepsilon\cdot\lc,
\]
with $\lc$ a local cochain of degree $0$: the components of arity one and two have no pole. The finite special vertices of Definition c11-defect (iii) are therefore absent at the level of the action in arity one and two, and in arity zero they come only from the mass term.
\end{proposition}

\begin{proof}
The coefficient of $\varepsilon^{-m}$ of $\ICT^{(k)}(x_1,\dots,x_k)$ is a colour-singlet local functional of degree $-k$ and total Weyl weight $4-\tfrac12\sum\lambda_{x_i}-2m$, of degree at most $\sum p_{x_i}$ in the position (Proposition c10-bound1). A colour singlet of degree $-1$ contains an antifield and a second adjoint field: its field weight is at least $4$, or $\tfrac72$ with the legs $(A,\lambda^*)$ when the generator supplies a spinor. A colour singlet of degree $-2$ has field weight at least $5$ ($(A,c^*)$, $(\lambda^*,\lambda^*)$), or $\tfrac{11}2$ with an odd number of spinor legs. The total weight is at least the field weight minus the position degree. For $m\ge1$ the required inequality fails in each case: arity one, $\Delta$: $4-2m<3$; $P_b$: $5-2m<4$; $\mathfrak{so}(V)$: $4-2m<3$; $\mathfrak{so}(W^\perp_{\mathbb R})$: $4-2m<4$; $\Sigma$: $\tfrac92-2m<\tfrac72$; $\Sigma^*$: $\tfrac72-2m<\tfrac52$; $K_b$: $3-2m<2$; arity two, $(\Delta,P_b)$: $5-2m<4$; $(K_b,P_{b'})$: $4-2m<3$; $(\xi,\xi')$: $5-2m<5$; $(\xi,\eta)$: $4-2m<4$; $(\xi,P_b)$: $\tfrac{11}2-2m<\tfrac{11}2$; $(\eta,\eta')$: $3-2m<3$; $(\Delta,\xi)$: $\tfrac92-2m<\tfrac92$. In arity zero, weight $4-2m$ with $m\ge1$ leaves weight $2$ and the legs $(A,A)$, the term of Theorem c4-quadratic. One loop gives integer powers of $\varepsilon$ and the first power of $\log\varepsilon$ (Lemma c4-local).
\end{proof}

\begin{proposition}[The logarithmic coefficient is $s$-closed]
\label{prop:at-log}
$s\,\lc=0$. In components: $\Lambda^{(0)}\lc^{(0)}=0$; $\Lambda^{(0)}\lc^{(1)}(x)+\{\lc^{(0)},S(x)\}=0$; and the arity-two formula of Lemma \ref{lem:at-sformulas} (i) applied to $\lc$ vanishes. With $\lc^{(0)}=u\,\Lambda^{(0)}N$, $N:=2\int(A,A^*)-4\int(\lambda^*,\lambda)-\int(c^*,c)$ (Theorem c4-log), $\lc^{(1)}(x)-\{uN,S(x)\}$ is $\Lambda^{(0)}$-closed of degree $-1$.
\end{proposition}

\begin{proof}
By [C07, Thm B] and Definition c8-quant the flow $W(P(\varepsilon,T),I_{\mathfrak{psl}(4|4)}-\ICT)$ has a limit; $Q$ acts on plane-wave values by polynomial substitutions, so the left side of (c8-qme) at scale $T$ has a limit, and so has the flow of $D$ (Lemma c8-transfer). In the flow of $D$ at order $\hbar$ the coefficient of $\log\varepsilon$ comes from the trees of the counterterm vertices of \eqref{eq:at-defect}, whose $K_0$-parts give $-\log\varepsilon\,(s\,\lc)$ and whose special-edge parts carry $(e^{-\varepsilon|p|^2}-1)\log\varepsilon\to0$; the cycles with a special edge have no $\log\varepsilon$ (Lemma \ref{lem:at-special}); the cycles without special edge are absent from the flow of $D$ at order $\hbar$, and the subgraphs of the reducible graphs are of the two previous kinds. The tree-level flow is injective on local cochains (proof of Theorem c8-L (i)), so $s\,\lc=0$. The last statement: by the derivation rule of the proof of Lemma \ref{lem:at-sformulas}, $\Lambda^{(0)}\{uN,S(x)\}=-\{\Lambda^{(0)}uN,S(x)\}=-\{\lc^{(0)},S(x)\}=\Lambda^{(0)}\lc^{(1)}(x)$, since $\Lambda^{(0)}S(x)=0$.
\end{proof}

\begin{remark}[Exactness of the logarithmic coefficient]
\label{rem:at-logexact}
Along $\Delta$, $\lc^{(1)}(\Delta)$ is quadratically $Q\big(u\int(\iota_x\AV,c^*)-\tfrac u2\int(\rho(x)\lambda^*,\lambda^*)\big)$ and $\{N,S(\Delta)\}=0$ (Theorem c9-quadratic (i)). Whether $\lc=s\Phi$ for a local cochain $\Phi$ of degree $-1$ is a computed statement: by Proposition \ref{prop:at-log} and Theorem \ref{thm:at-H} below, $\lc^{(1)}(x)-\{uN,S(x)\}\equiv l_x\,u\,S(x)$ modulo $\operatorname{Im}\Lambda^{(0)}$ under Hypothesis \ref{hyp:at-H}, with $l_{[x,y]}=l_x+l_y$; the solutions are $l_x=c\lambda_x+aB_x$, $B_x$ the eigenvalue of the outer derivation $B$ (Definition c8-B). The inner part is $s$-exact; $a$ is the coefficient of $\log\varepsilon$ along the class $[G_B]$ (gate G-logex of the engine specification).
\end{remark}

\subsection{Placements, the special-edge lemma and the graph classes}

\begin{lemma}[Placements of $S^{(1)}(x)$ on a cycle]
\label{lem:at-place}
The propagator has the components $A\otimes A$, $c\otimes A^*|_V$, $A^*|_V\otimes c$, $\lambda\otimes\lambda$ only (Lemma o-legs), and a special edge joins a field to its antifield. On a cycle with exactly one special edge, the vertex $S(x)$ occurs in one of the following placements, and in no other:
\begin{itemize}
\item (line) both cycle legs on edges $P$: the term $(v.\AV,A^*|_V)$ for even $x$ (an $\AV$--$c$ line), the term $T_1$ for $\xi$ and $x^\mu T_1(\rho(e_\mu)\eta)$ for $\eta$ (a $\lambda$--$c$ line);
\item (at the special edge) one cycle leg on the special edge and the other on an edge $P$: for even $x$, ($A^*$ on the special edge, $A$ on $P$), ($\AV$ on the special edge, $A^*|_V$ on $P$), ($\lambda^*$ on the special edge, $\lambda$ on $P$), ($c^*$ on the special edge, $c$ on $P$); for $\xi$, $T_1$ with either leg on the special edge, $T_2$ with $A^*|_\perp$ on it, $T_3$, $T_4$ with $\lambda^*$ on it; for $\eta$ the same at $\rho(x)\eta$ and $s_0$ with $\lambda^*$ on it;
\item (cubic) for $\xi$ and $\eta$: $T_5$ with two legs $A$ on the cycle (both on edges $P$, the special edge elsewhere, or one on the special edge) and the third leg external, or with $\lambda^*$ on the special edge, one $A$ on $P$ and one $A$ external.
\end{itemize}
$S^{(2)}(x,y)$ has no line placement; its legs $\lambda^*$, $c^*$ lie on the special edge or are external.
\end{lemma}

\begin{proof}
A vertex on a cycle uses two of its legs on cycle edges. $\lambda^*$, $c^*$ and $A^*|_\perp$ have no component of $P$, so they lie on the special edge or are external; $A^*|_V$ lies on $c\otimes A^*|_V$; the remaining assignments are read off \eqref{eq:at-S1list} and Table c11-terms. For $S^{(2)}$, Lemma o-legs.
\end{proof}

\begin{lemma}[Special-edge lemma for action cycles]
\label{lem:at-special}
Let $\mathcal C$ be a cycle all of whose vertices lie on it, with external legs at generic momenta, $n_t$ edges, one special edge $e_\sigma$ (Schwinger time in $[0,\varepsilon]$) and the other edges $P(\varepsilon,T)$, possibly one composite line, with vertices from $I$, from $S(x)$ and $S(y)$ or $S^{(2)}(x,y)$; along $\eta$, $\Delta$, $\mathfrak{so}(V)$ and $K_b$ the position factors at the vertex $S(x)$ are realised by derivatives in a momentum $r$ injected there, at $r=0$. Let $D$ be its superficial degree, the number of derivatives of the local functionals in its sector (Definition \ref{def:at-ob}). Then at $T=\infty$
\[
\mathcal C(\varepsilon)=\sum_{m\ge1}\varepsilon^{-m}a_m+a_0+o(1),\qquad 16\pi^2a_m=\frac1m\int_{\alpha_\sigma\ \text{smallest}}\Big[\big(\min_{f\ne\sigma}\alpha_f\big)^m-\alpha_\sigma^m\Big]F_m\,d\alpha,\qquad 16\pi^2a_0=\int_{\alpha_\sigma\ \text{smallest}}F_0\log\frac{\min_{f\ne\sigma}\alpha_f}{\alpha_\sigma}\,d\alpha,
\]
with $F_m$ the $m$-th simplex integrand of Definition c10-cycle; there is no term in $\log\varepsilon$, $a_m$ is a polynomial of degree $D-2m$ in the external momenta and $r$, and $a_0$ is a polynomial of degree $D$, zero for $D<0$. At finite $T$ the expansion has the same coefficients and $T\partial_T\mathcal C(\varepsilon)=O(\varepsilon)$ at fixed $T$.
\end{lemma}

\begin{proof}
As in Theorems c11-simplex and c11-oddsimplex: the Schwinger domain $\{t_\sigma\le\varepsilon\le t_f\}$ is the difference of the domains with the cutoff shapes $1/\min_{f\ne\sigma}\alpha_f$ and $1/\min_f\alpha_f$, which differ on the region where $\alpha_\sigma$ is the smallest coordinate; Proposition c10-shape applies to each, the coefficients of $\log\varepsilon$ cancel, and the finite parts differ by the stated integral. The region of large total Schwinger time converges at $T=\infty$ because the external momenta are generic and the composite line is of order $|p|^{-3}$ (Proposition c9-composite (iv)); the injected momentum enters polynomially. At finite $T$, $T\partial_T$ replaces one edge by $(Q^{GF}\otimes1)K_T$, which confines the loop momentum to $|l|^2\lesssim T^{-1}$, where the special edge is $O(\varepsilon|l|^2)$; the position factors grow at most as $T^{p_x/2}$ and are compensated by the decay of the other edges, as in Remark c8-K.
\end{proof}

\begin{theorem}[Graph classes and the local formula for $\Ob^{(1)}(x)$]
\label{thm:at-classes}
Fix a generator $x$ and $n\ge2$. At order $\hbar$, with $F_1=0$, the part with $n$ legs of the arity-one component along $x$ of the flow of $D$ is the sum of the following classes of connected graphs with edges $P(\varepsilon,T)$, one vertex from \eqref{eq:at-defect} and one vertex carrying $x$:
\begin{itemize}
\item (C) irreducible defect cycles: every vertex on the cycle, exactly one special edge, exactly one vertex from $S(x)$ in a placement of Lemma \ref{lem:at-place}, the other vertices cubic or quartic vertices of $I$ with their external legs; the special edge joins $S(x)$ to $I$ (placement at the special edge) or two vertices of $I$ (line and cubic placements);
\item (Rd) reducible graphs: a subgraph of type (C) with fewer legs, or a one-loop subgraph of the arity-zero defect, joined by edges $P(\varepsilon,T)$ to trees of vertices of $I$ and of $S(x)$;
\item (Br) graphs whose special edge is a bridge;
\item (Ct) trees with one counterterm vertex of \eqref{eq:at-defect} of order $\hbar$, with the $K_0$-part and the special-edge part of each bracket;
\item (K) the constant $\hbar\Delta_\varepsilon S(x)$.
\end{itemize}
As $\varepsilon\to0$: (Br) together with the special-edge parts of (Ct) tends to $0$; (Rd) tends to the trees of the components of $\Ob$ with fewer legs (arity one joined to $I$, arity zero joined to $S(x)$ and $I$); the $\log\varepsilon$ terms of (Ct) cancel identically (Proposition \ref{prop:at-log}) and (C) has none; the poles satisfy
\begin{equation}
\sum_{\mathcal C\in(C),\,n=2}a_1(\mathcal C)=\tfrac94\,u\,\Big\{\int(A,A),S(x)\Big\},\qquad \sum_{\mathcal C\in(C),\,n}a_m(\mathcal C)=0\quad(m\ge2\ \text{or}\ n\ge3),
\label{eq:at-pole}
\end{equation}
and
\begin{equation}
\hbar\,\Ob^{(1)}(x)\big|_{n\ \text{legs}}=\sum_{\mathcal C\in(C),\,n\ \text{legs}}a_0(\mathcal C),
\label{eq:at-localformula}
\end{equation}
a local functional, independent of $T$, with $F_1=0$. With an arity-zero finite counterterm included the same graphs give $\mathcal R^{(1)}(x)$ (Definition \ref{def:at-ob}), the trees of $-\{F_1^{(0)},S(x)\}_\varepsilon$ having the limit $-\{F_1^{(0)},S(x)\}$.
\end{theorem}

\begin{proof}
Order $\hbar$ forces one cycle with all vertices of order $\hbar^0$ or a tree with one vertex of order $\hbar$ (Proposition c3-graphs); arity one places $x$ on one vertex. The order-$\hbar^0$ marked vertices of $D^{(1)}(x)$ and of $D^{(0)}$ (with a separate $S(x)$) are special-edge parts, so the special edge lies on the cycle, on a bridge, or on a tree attached to the cycle; in the last case it is a bridge. A one-loop graph with the special edge on the cycle is a cycle with trees attached; if a tree is attached through an edge, the graph is reducible, otherwise it is in (C). (Br): removing the bridge separates the cubic or $S(x)$ side from a one-loop subgraph of the flow at generic momentum $p$; summed with the counterterm trees of the same subgraph and the special-edge part of the counterterm bracket across the same edge, the edge carries $e^{-\varepsilon|p|^2}-1$ times the renormalised subgraph, which has a limit (Definition c8-quant), so the sum tends to $0$; along $\eta$ and $K_b$ the derivatives in the injected momentum act on $(e^{-\varepsilon|p+r|^2}-1)$ and on a function of class $C^2$ at $p\ne0$. (Rd): the subgraph equals the corresponding component of the obstruction at its momenta plus $o(1)$ after its counterterm trees are included (induction on the number of legs; the arity-zero two-leg subgraphs are Theorem c4-obstruction), and the edges $P(\varepsilon,T)$ at generic momenta converge to $P(0,T)$. (Ct): by Proposition \ref{prop:at-poles} the only pole is $-\varepsilon^{-1}\tfrac94u\{\int(A,A),S(x)\}$ in $-\{\ICT^{(0)},S(x)\}_\varepsilon$, which has two legs ($\{\int(A,A),T_5\}=0$, $T_5$ having no $A^*$); its special-edge part is part of (Br). The limit exists (proof of Proposition \ref{prop:at-log}), so the poles of (C) cancel this term, which is \eqref{eq:at-pole}; Lemma \ref{lem:at-special} gives the finite parts. The tree-level flow of the arity-one part of $\Ob$ with $n$ legs is $\Ob^{(1)}(x)|_n$ plus trees with fewer-leg components, which are the limits of (Rd), so \eqref{eq:at-localformula} follows; $T$-independence is Definition c8-ob and Lemma \ref{lem:at-special}.
\end{proof}

\begin{remark}[Counting]
\label{rem:at-counting}
Each graph of (C) is weighted by $1/|\mathrm{Aut}|$ and the order factors of $\tfrac12\{I,I\}_\varepsilon$, which cancel against the two orders of the vertices joined by the special edge (table caption of Section 9.1); the colour factor of a cycle with $n$ external adjoint legs is a trace of $\mathrm{ad}$'s, reduced by Lemma c9-legs for $n=3$; for $n=4$ it contains, besides products of structure constants, the symmetrised trace $\operatorname{str}_{\mathfrak g}(\mathrm{ad}_{X_1}\mathrm{ad}_{X_2}\mathrm{ad}_{X_3}\mathrm{ad}_{X_4})$, which the four-leg sectors keep as a separate colour structure (realisations $\mathfrak{su}(2)$, $\mathfrak{su}(3)$ and $\mathfrak{gl}(3)$ in the engines).
\end{remark}

\subsection{Composite lines, variation kernels and the second evaluation}

\begin{proposition}[Composite lines from the commutators with the regulator]
\label{prop:at-composite}
Let $P\star=Q^{GF}g(H)$, $g(H)=\int_\varepsilon^Te^{-tH}dt$, $E(H)=e^{-\varepsilon H}-e^{-TH}$. The line placement of $S(x)$ is the composite kernel $2(P\star\,x.\,P\star\otimes1)K_0$ (Corollary qs-edges (ii)), and
\[
P\star\,x.\,P\star=x.\,(P\star)^2-[x.,Q^{GF}]\,g(H)\,P\star-Q^{GF}[x.,g(H)]\,P\star .
\]
\begin{itemize}
\item $P_b$, $\mathfrak{so}(V)$, $\mathfrak{so}(W^\perp_{\mathbb R})$: the commutators vanish; the kernel is $2(x.(P\star)^2\otimes1)K_0$, zero along $\mathfrak{so}(W^\perp_{\mathbb R})$ ($v.=0$ on the slots of $c$ and $\AV$); $[Q,P\star x.P\star]=0$.
\item $\Delta$: $[v.,Q^{GF}]=-4Q^{GF}$, $[v.,g(H)]=4Hg'(H)$; $2P\star v.P\star=v.(P\star)^2+(P\star)^2v.$ and $[Q,P\star v.P\star]=0$ (Proposition c9-composite).
\item $\xi$: $[\xi.,Q^{GF}]=C_\xi$, $[\xi.,g(H)]=0$; $P\star\xi.P\star=C_\xi Q^{GF}g^2-\xi.(Q^{GF})^2g^2$ and $[Q,P\star\xi.P\star]=E(H)C_\xi g(H)\ne0$ (Lemma c11-lines (iii)).
\item $\eta$: $[\eta.,Q^{GF}]=\sum_\mu x^\mu C_{\xi_\mu}+C^0_\eta$, $[\eta.,H]=2\sum_\mu\xi_\mu.\partial_\mu$, $\xi_\mu=\rho(e_\mu)\eta$ (Proposition qs-comm (ii)).
\item $K_b$: $[v.,Q^{GF}]=4(b,x)Q^{GF}+C_K$, $[v.,H]=4(b,x)H-4(w-1)(b,\partial)+4b_\nu\Sigma^{\rho\nu}\partial_\rho$ (Lemma c11-lines (v)); $[Q,P\star v.P\star]=E(H)v.P\star-P\star v.E(H)\ne0$.
\end{itemize}
The variation kernel $\mathcal V_x=(x.\otimes1+1\otimes x.)P(\varepsilon,T)$ is $0$ along $P_b$, $\mathfrak{so}(V)$, $\mathfrak{so}(W^\perp_{\mathbb R})$, $-4(\varepsilon\partial_\varepsilon+T\partial_T)P$ along $\Delta$, $(C_\xi g(H)\otimes1)K_0$ along $\xi$, and the kernels of Corollary qs-edges (i) along $\eta$ and $K_b$.
\end{proposition}

\begin{proof}
$P\star x.P\star=x.P\star P\star+[P\star,x.]P\star$ and $[P\star,x.]=-[x.,Q^{GF}]g-Q^{GF}[x.,g]$ (graded). The listed commutators are those of the cited statements; $[Q,P\star v.P\star]=E\,v.P\star-P\star v.E$ for even $v.$ with $[v.,Q]=0$ follows from $QP\star+P\star Q=E$.
\end{proof}

\begin{theorem}[Second evaluation: $Q$ of the renormalised arity-one amplitudes]
\label{thm:at-second}
Let $W^{(1)}(x)$ be the arity-one component along $x$ of $W\big(P(\varepsilon,T),I_{\mathfrak{psl}(4|4)}-\ICT\big)$, the linearised flow of $S(x)-\ICT^{(1)}(x)$ at $I-\ICT^{(0)}$, and $W^{(1)}_{\rm ren}(x)$ the $\varepsilon\to0$ limit of its coefficient of $\hbar$ at generic momenta.
\begin{itemize}
\item (i) At $T=\infty$: $\hbar\big(\text{tree flow of }\Ob\big)^{(1)}(x)=Q\,W^{(1)}_{\rm ren}(x)$, the action of $Q$ on plane-wave values commuting with the limit; at finite $T$ add $\{W^{(0)},W^{(1)}(x)\}_T+\hbar\Delta_TW^{(1)}(x)$ (smooth kernels). The local functional $\Ob^{(1)}(x)|_n$ is the part of $Q\,W^{(1)}_{\rm ren}(x)$ with $n$ legs minus the trees of the components with fewer legs.
\item (ii) Conjugation. With $D_x$ the derivation induced by $P\star x.$, $D^{(3)}_x$ the first-order operator given by the cubic term of $S(x)$ with one leg $A$ joined by $P$ to the argument, $\partial^{(3)}_x$ the second-order operator with both legs $A$ joined by $P$, and $\mathcal K_x=2(P\star x.P\star\otimes1)K_0$,
\begin{equation}
e^{\hbar\partial_P}\,S(x)\,e^{-\hbar\partial_P}=S(x)+\hbar\big(D_x+D^{(3)}_x\big)+\tfrac{\hbar^2}2\big(\partial_{\mathcal K_x}+\partial^{(3)}_x\big),
\label{eq:at-conj}
\end{equation}
and the series terminates. Hence, with $W=W_0+\hbar W_1$ and $B_{\mathcal K}(f,g):=\partial_{\mathcal K}(fg)-(\partial_{\mathcal K}f)g-f\,\partial_{\mathcal K}g$ (Koszul signs understood), the linearised flow of $S(x)$ at $J$ is $S(x)+(D_x+D^{(3)}_x)W+\tfrac\hbar2(\partial_{\mathcal K_x}+\partial^{(3)}_x)W+\tfrac14(B_{\mathcal K_x}+B^{(3)}_x)(W,W)$, whose part of order $\hbar$ is
\[
(D_x+D^{(3)}_x)W_1+\tfrac12(\partial_{\mathcal K_x}+\partial^{(3)}_x)W_0+\tfrac12(B_{\mathcal K_x}+B^{(3)}_x)(W_0,W_1);
\]
the composite cycles $\tfrac12\partial_{\mathcal K_x}W_0$ and the cubic cycles $\tfrac12\partial^{(3)}_xW_0$ carry the singular parts $\ICT^{(1)}(x)$.
\item (iii) $\mathfrak{so}(W^\perp_{\mathbb R})$: $\ICT^{(1)}(R_r)=0$ and $\Ob^{(1)}(R_r)=0$ exactly.
\item (iv) $P_b$, $\mathfrak{so}(V)$: $\Ob^{(1)}(x)$ is $Q$ of the renormalised composite cycles minus $\Lambda^{(0)}\ICT^{(1)}(x)$ (the composite part).
\item (v) $\Delta$: $\Ob^{(1)}(\Delta)=4\,\lc^{(0)}+(\text{composite part})$ in every sector in which the arity-zero obstruction is removed (Proposition c9-structure (ii)), a theorem of the construction used as cross-check X-scal.
\item (vi) $\xi$: $\Ob^{(1)}(\xi)$ is $Q$ of the renormalised composite and cubic cycles, minus $\Lambda^{(0)}\ICT^{(1)}(\xi)$, plus the variation-kernel graphs: the one-loop graphs of $W$ with one edge replaced by $\mathcal V_\xi$ (Proposition qs-ward), on the cycle.
\end{itemize}
\end{theorem}

\begin{proof}
(i) Lemma c8-transfer at arity one; $d_{CE}$ vanishes on arity zero. (ii) $\mathrm{ad}_{\partial_P}$ applied to the multiplication by the quadratic part gives $D_x$ (Lemma c9-SDelta (iv), its proof applying to every quadratic $S(x)$ whose legs have no common component of $P$); applied twice it gives $\partial_{\mathcal K_x}$ by (c9-commutator), a constant-coefficient operator commuting with $\partial_P$. For the cubic term, $\partial_PT_5=0$ (the contraction of its two legs $A$ gives $[X_a,X^a]=0$), $\mathrm{ad}_{\partial_P}T_5=D^{(3)}_x$, $\mathrm{ad}^2_{\partial_P}T_5=\partial^{(3)}_x$, whose coefficient $\lambda^*$ has no component of $P$, so $\mathrm{ad}^3_{\partial_P}T_5=0$. The expansion of the linearised flow is the product rule of the second-order operators (proof of Proposition c9-conj). (iii) $v.$ commutes with $Q$, $Q^{GF}$, $H$, preserves the pairing, annihilates $I$ and the $\mathfrak{so}(W^\perp_{\mathbb R})$-invariant counterterms, and $\mathcal K_x=0$; the one-loop part of the linearised flow is $D_xW_1$, whose singular part is that of $W_1$ with $P\star v.$ on a leg, already removed by $\ICT^{(0)}$; so $\ICT^{(1)}(R_r)=0$. Then $QW^{(1)}=[Q,D_x]W_1+D_xQW_1$, $[Q,D_x]$ is the derivation by $E(H)v.$, whose action on $W_1$ tends to $v.W_{1,{\rm ren}}=0$, and $D_xQW_1$ is $S(x)$ joined to a leg of the flow of the arity-zero defect, a tree of $\Ob^{(0)}$ with $S(x)$. (iv) The same with $\mathcal K_x\ne0$ and $[Q,\mathcal K_x]=0$; $v.$ annihilates the translation- and rotation-invariant $I$ and $\ICT^{(0)}$. (v) Proposition c9-structure (ii), whose proof uses (c9-D1b) and the scaling identity (c9-ward), both identities of the flow. (vi) Proposition qs-ward and $[Q,P\star\xi.]=E(H)\xi.$, $[Q,\mathcal K_\xi]=2(E C_\xi g\otimes1)K_0$.
\end{proof}

\subsection{Position decomposition}

\begin{proposition}[Position parts at the level of the action]
\label{prop:at-pos}
For a local functional $G$ with polynomial dependence on the position, its position part along $b$ is $-\{S(P_b),G\}$ (Lemma c11-canonical). Graph by graph, with the momentum injection of Lemma \ref{lem:at-special}:
\begin{itemize}
\item (i) the position part along $b$ of a cycle along $\Delta$ is $2$ times the same cycle along $P_b$; along $M_\omega$ the cycle along $P_{\omega b}$; along $\eta$ the cycle along $\rho(b)\eta\in\Sigma$; along $K_b$ the cycle along the generator with the Killing field $\partial_{b'}v$, which is $2(b,b')\cdot\tfrac12\Delta$ plus the rotation $x\mapsto2(b,x)b'-2(b',x)b$ (Remark c7-KDelta), and the second position part along $(b',b'')$ is the cycle along the translation $P_{\partial_{b'}\partial_{b''}v}$;
\item (ii) hence
\[
-\{S(P_b),\Ob^{(1)}(\Delta)\}=2\Ob^{(1)}(P_b),\qquad -\{S(P_b),\Ob^{(1)}(\eta)\}=\Ob^{(1)}(\rho(b)\eta),\qquad -\{S(P_b),\Ob^{(1)}(M_\omega)\}=\Ob^{(1)}(P_{\omega b}),
\]
and for $K_b$ the analogous identities; the same holds for $\ICT^{(1)}$ and $\lc^{(1)}$;
\item (iii) the intrinsic part of $\Ob^{(1)}(x)$ is $\Ob^{(1)}(x)$ minus a fixed lift of its position parts, $\int x^\mu o_\mu$ with $o_\mu$ chosen densities of the position parts; it is translation invariant and changes by a translation-invariant functional when the densities $o_\mu$ change by total derivatives.
\end{itemize}
\end{proposition}

\begin{proof}
(i) As Proposition c11-position (i): the factor $x^\mu$ at the vertex $S(x)$ is $-i\partial/\partial r_\mu$ of $e^{i(r,x)}$; its action on the momentum delta function is the position part, $-i$ times the graph with the vertex differentiated along $e_\mu$; the derivative of the Killing data of $\Delta$ along $e_\mu$ is $2e_\mu$, of $\eta$ is $\rho(e_\mu)\eta$, of $M_\omega$ is $\omega e_\mu$, and of $K_b$ is $\partial_\mu v$. The special edge, the edges $P$ and the vertices of $I$ are translation invariant. (ii) The obstruction is the sum \eqref{eq:at-localformula}. (iii) $-\{S(P_b),\int x^\mu o_\mu\}=\int b^\mu o_\mu$.
\end{proof}

\subsection{The sectors that can be nonzero}

\begin{theorem}[Sector bound]
\label{thm:at-bound}
$\Ob^{(1)}(x)$ is a colour-singlet local functional of degree $0$, Weyl weight $4-\tfrac12\lambda_x$, degree at most $p_x$ in the position, linear in the Killing data of $x$, of fermion parity $|x|$, equivariant for $\mathrm{Spin}(4)\times\mathrm{Spin}(6)$ and the adjoint action of $\mathfrak g$, and invariant under the reflection of one direction of $V$ together with one of $W^\perp_{\mathbb R}$ (Proposition c9-basis). A leg set with $n$ legs, field weight $w_f$ and $D$ derivatives occurs at position degree $p$ only if $D=4-\tfrac12\lambda_x+p-w_f\ge0$, the numbers $n_\pm$ of $S_\pm$-indices and $n_V$ of $V$-indices (legs $A|_V$, $A^*|_V$, derivatives, positions, the vectors $b$, $\omega$ of the data) satisfy $n_++n_V\equiv n_-+n_V\equiv0\pmod 2$ (centre of $\mathrm{Spin}(4)$), and the $SU(4)$-charges ($W$: $+1$, $W^*$: $-1$, $W^\perp_{\mathbb R}$: $2$) add to $0$ modulo $4$. The same rules with degree $-1$ and weight $4-\tfrac12(\lambda_x+\lambda_y)$ hold for $\Ob^{(2)}(x,y)$, and with degrees $-1$, $-2$ for $F_1^{(1)}$, $F_1^{(2)}$. The leg sets allowed by these rules are listed in Tables \ref{tab:at-census1} and \ref{tab:at-census2}; in particular, at position degree $0$ the number of legs is at most $4$ along $\Delta$, $\mathfrak{so}(V)$, $\mathfrak{so}(W^\perp_{\mathbb R})$ and $\Sigma$, at most $5$ along $P_b$ (the leg set $(\AV)^5$ without derivatives), at most $3$ along $\Sigma^*$, and the intrinsic part along $K_b$ has two legs only (its cubic leg sets $(\AV,\AV,\AV)_0$, $(\AV,\Ap,\Ap)_0$ pair a symmetric kinematic factor with the antisymmetric colour factor, or are parity odd, and vanish); the parts of higher position degree are fixed by Proposition \ref{prop:at-pos}. The arity-two obstruction has at most three legs at position degree $0$, four along $(\xi,P_b)$, and none at position degree $0$ along $(\eta,\eta')$.
\end{theorem}

\begin{proof}
Weight and degree: Definition \ref{def:at-ob}. A local functional with $D$ derivatives has the weight of its legs plus $D$ minus the position degree. The centres of $SU(2)\times SU(2)$ and of $SU(4)$ act on each leg and datum by the stated signs and charges, and an invariant exists only if they act trivially. A one-loop graph with external legs in the adjoint has at least two of them. Equivariance: Proposition c7-cov and Proposition c7-select. The census is the output of the enumeration of all multisets of at most six legs.
\end{proof}

{\footnotesize
\begin{longtable}{>{\raggedright}p{2.3cm}|p{0.6cm}|p{0.25cm}|p{0.25cm}|>{\raggedright\arraybackslash}p{11.6cm}}
cochain & weight & $p$ & $n$ & leg sets $(\cdots)_D$\\
\hline
\endhead
$\Ob^{(1)}(\Delta)$ & $4$ & 0 & 2 & $(c,A^*_V)_{1}$, $(A_V,A_V)_{2}$, $(A_\perp,A_\perp)_{2}$, $(\lambda_+,\lambda_-)_{1}$ \\
 &  & 0 & 3 & $(c,c,c^*)_{0}$, $(c,A_V,A^*_V)_{0}$, $(c,A_\perp,A^*_\perp)_{0}$, $(c,\lambda_+,\lambda^*_+)_{0}$, $(c,\lambda_-,\lambda^*_-)_{0}$, $(A_V,A_V,A_V)_{1}$, $(A_V,A_\perp,A_\perp)_{1}$, $(A_V,\lambda_+,\lambda_-)_{0}$, $(A_\perp,\lambda_+,\lambda_+)_{0}$, $(A_\perp,\lambda_-,\lambda_-)_{0}$ \\
 &  & 0 & 4 & $(A_V,A_V,A_V,A_V)_{0}$, $(A_V,A_V,A_\perp,A_\perp)_{0}$, $(A_\perp,A_\perp,A_\perp,A_\perp)_{0}$ \\
 &  & 1 & 2 & $(c,A^*_V)_{2}$, $(A_V,A_V)_{3}$, $(A_\perp,A_\perp)_{3}$, $(\lambda_+,\lambda_-)_{2}$ \\
 &  & 1 & 3 & $(c,c,c^*)_{1}$, $(c,A_V,A^*_V)_{1}$, $(c,A_\perp,A^*_\perp)_{1}$, $(c,\lambda_+,\lambda^*_+)_{1}$, $(c,\lambda_-,\lambda^*_-)_{1}$, $(A_V,A_V,A_V)_{2}$, $(A_V,A_\perp,A_\perp)_{2}$, $(A_V,\lambda_+,\lambda_-)_{1}$, $(A_\perp,\lambda_+,\lambda_+)_{1}$, $(A_\perp,\lambda_-,\lambda_-)_{1}$ \\
 &  & 1 & 4 & $(c,c,A_V,c^*)_{0}$, $(c,c,\lambda^*_+,\lambda^*_-)_{0}$, $(c,A_V,A_V,A^*_V)_{0}$, $(c,A_V,A_\perp,A^*_\perp)_{0}$, $(c,A_V,\lambda_+,\lambda^*_+)_{0}$, $(c,A_V,\lambda_-,\lambda^*_-)_{0}$, $(c,A_\perp,A_\perp,A^*_V)_{0}$, $(c,A_\perp,\lambda_+,\lambda^*_-)_{0}$, $(c,A_\perp,\lambda_-,\lambda^*_+)_{0}$, $(A_V,A_V,A_V,A_V)_{1}$, $(A_V,A_V,A_\perp,A_\perp)_{1}$, $(A_V,A_V,\lambda_+,\lambda_-)_{0}$, $(A_V,A_\perp,\lambda_+,\lambda_+)_{0}$, $(A_V,A_\perp,\lambda_-,\lambda_-)_{0}$, $(A_\perp,A_\perp,A_\perp,A_\perp)_{1}$, $(A_\perp,A_\perp,\lambda_+,\lambda_-)_{0}$ \\
 &  & 1 & 5 & $(A_V,A_V,A_V,A_V,A_V)_{0}$, $(A_V,A_V,A_V,A_\perp,A_\perp)_{0}$, $(A_V,A_\perp,A_\perp,A_\perp,A_\perp)_{0}$ \\
\hline
$F_1^{(1)}(\Delta)$ & $4$ & 0 & 2 & $(c,c^*)_{0}$, $(A_V,A^*_V)_{0}$, $(A_\perp,A^*_\perp)_{0}$, $(\lambda_+,\lambda^*_+)_{0}$, $(\lambda_-,\lambda^*_-)_{0}$ \\
 &  & 1 & 2 & $(c,c^*)_{1}$, $(A_V,A^*_V)_{1}$, $(A_\perp,A^*_\perp)_{1}$, $(\lambda_+,\lambda^*_+)_{1}$, $(\lambda_-,\lambda^*_-)_{1}$ \\
 &  & 1 & 3 & $(c,A_V,c^*)_{0}$, $(c,\lambda^*_+,\lambda^*_-)_{0}$, $(A_V,A_V,A^*_V)_{0}$, $(A_V,A_\perp,A^*_\perp)_{0}$, $(A_V,\lambda_+,\lambda^*_+)_{0}$, $(A_V,\lambda_-,\lambda^*_-)_{0}$, $(A_\perp,A_\perp,A^*_V)_{0}$, $(A_\perp,\lambda_+,\lambda^*_-)_{0}$, $(A_\perp,\lambda_-,\lambda^*_+)_{0}$ \\
\hline
$\Ob^{(1)}(P_b)$ & $5$ & 0 & 2 & $(c,A^*_V)_{2}$, $(A_V,A_V)_{3}$, $(A_\perp,A_\perp)_{3}$, $(\lambda_+,\lambda_-)_{2}$ \\
 &  & 0 & 3 & $(c,c,c^*)_{1}$, $(c,A_V,A^*_V)_{1}$, $(c,A_\perp,A^*_\perp)_{1}$, $(c,\lambda_+,\lambda^*_+)_{1}$, $(c,\lambda_-,\lambda^*_-)_{1}$, $(A_V,A_V,A_V)_{2}$, $(A_V,A_\perp,A_\perp)_{2}$, $(A_V,\lambda_+,\lambda_-)_{1}$, $(A_\perp,\lambda_+,\lambda_+)_{1}$, $(A_\perp,\lambda_-,\lambda_-)_{1}$ \\
 &  & 0 & 4 & $(c,c,A_V,c^*)_{0}$, $(c,c,\lambda^*_+,\lambda^*_-)_{0}$, $(c,A_V,A_V,A^*_V)_{0}$, $(c,A_V,A_\perp,A^*_\perp)_{0}$, $(c,A_V,\lambda_+,\lambda^*_+)_{0}$, $(c,A_V,\lambda_-,\lambda^*_-)_{0}$, $(c,A_\perp,A_\perp,A^*_V)_{0}$, $(c,A_\perp,\lambda_+,\lambda^*_-)_{0}$, $(c,A_\perp,\lambda_-,\lambda^*_+)_{0}$, $(A_V,A_V,A_V,A_V)_{1}$, $(A_V,A_V,A_\perp,A_\perp)_{1}$, $(A_V,A_V,\lambda_+,\lambda_-)_{0}$, $(A_V,A_\perp,\lambda_+,\lambda_+)_{0}$, $(A_V,A_\perp,\lambda_-,\lambda_-)_{0}$, $(A_\perp,A_\perp,A_\perp,A_\perp)_{1}$, $(A_\perp,A_\perp,\lambda_+,\lambda_-)_{0}$ \\
 &  & 0 & 5 & $(A_V,A_V,A_V,A_V,A_V)_{0}$, $(A_V,A_V,A_V,A_\perp,A_\perp)_{0}$, $(A_V,A_\perp,A_\perp,A_\perp,A_\perp)_{0}$ \\
\hline
$F_1^{(1)}(P_b)$ & $5$ & 0 & 2 & $(c,c^*)_{1}$, $(A_V,A^*_V)_{1}$, $(A_\perp,A^*_\perp)_{1}$, $(\lambda_+,\lambda^*_+)_{1}$, $(\lambda_-,\lambda^*_-)_{1}$ \\
 &  & 0 & 3 & $(c,A_V,c^*)_{0}$, $(c,\lambda^*_+,\lambda^*_-)_{0}$, $(A_V,A_V,A^*_V)_{0}$, $(A_V,A_\perp,A^*_\perp)_{0}$, $(A_V,\lambda_+,\lambda^*_+)_{0}$, $(A_V,\lambda_-,\lambda^*_-)_{0}$, $(A_\perp,A_\perp,A^*_V)_{0}$, $(A_\perp,\lambda_+,\lambda^*_-)_{0}$, $(A_\perp,\lambda_-,\lambda^*_+)_{0}$ \\
\hline
$\Ob^{(1)}(M_\omega)$ & $4$ & 0 & 2 & $(c,A^*_V)_{1}$, $(A_V,A_V)_{2}$, $(A_\perp,A_\perp)_{2}$, $(\lambda_+,\lambda_-)_{1}$ \\
 &  & 0 & 3 & $(c,c,c^*)_{0}$, $(c,A_V,A^*_V)_{0}$, $(c,A_\perp,A^*_\perp)_{0}$, $(c,\lambda_+,\lambda^*_+)_{0}$, $(c,\lambda_-,\lambda^*_-)_{0}$, $(A_V,A_V,A_V)_{1}$, $(A_V,A_\perp,A_\perp)_{1}$, $(A_V,\lambda_+,\lambda_-)_{0}$, $(A_\perp,\lambda_+,\lambda_+)_{0}$, $(A_\perp,\lambda_-,\lambda_-)_{0}$ \\
 &  & 0 & 4 & $(A_V,A_V,A_V,A_V)_{0}$, $(A_V,A_V,A_\perp,A_\perp)_{0}$, $(A_\perp,A_\perp,A_\perp,A_\perp)_{0}$ \\
 &  & 1 & 2 & $(c,A^*_V)_{2}$, $(A_V,A_V)_{3}$, $(A_\perp,A_\perp)_{3}$, $(\lambda_+,\lambda_-)_{2}$ \\
 &  & 1 & 3 & $(c,c,c^*)_{1}$, $(c,A_V,A^*_V)_{1}$, $(c,A_\perp,A^*_\perp)_{1}$, $(c,\lambda_+,\lambda^*_+)_{1}$, $(c,\lambda_-,\lambda^*_-)_{1}$, $(A_V,A_V,A_V)_{2}$, $(A_V,A_\perp,A_\perp)_{2}$, $(A_V,\lambda_+,\lambda_-)_{1}$, $(A_\perp,\lambda_+,\lambda_+)_{1}$, $(A_\perp,\lambda_-,\lambda_-)_{1}$ \\
 &  & 1 & 4 & $(c,c,A_V,c^*)_{0}$, $(c,c,\lambda^*_+,\lambda^*_-)_{0}$, $(c,A_V,A_V,A^*_V)_{0}$, $(c,A_V,A_\perp,A^*_\perp)_{0}$, $(c,A_V,\lambda_+,\lambda^*_+)_{0}$, $(c,A_V,\lambda_-,\lambda^*_-)_{0}$, $(c,A_\perp,A_\perp,A^*_V)_{0}$, $(c,A_\perp,\lambda_+,\lambda^*_-)_{0}$, $(c,A_\perp,\lambda_-,\lambda^*_+)_{0}$, $(A_V,A_V,A_V,A_V)_{1}$, $(A_V,A_V,A_\perp,A_\perp)_{1}$, $(A_V,A_V,\lambda_+,\lambda_-)_{0}$, $(A_V,A_\perp,\lambda_+,\lambda_+)_{0}$, $(A_V,A_\perp,\lambda_-,\lambda_-)_{0}$, $(A_\perp,A_\perp,A_\perp,A_\perp)_{1}$, $(A_\perp,A_\perp,\lambda_+,\lambda_-)_{0}$ \\
 &  & 1 & 5 & $(A_V,A_V,A_V,A_V,A_V)_{0}$, $(A_V,A_V,A_V,A_\perp,A_\perp)_{0}$, $(A_V,A_\perp,A_\perp,A_\perp,A_\perp)_{0}$ \\
\hline
$F_1^{(1)}(M_\omega)$ & $4$ & 0 & 2 & $(c,c^*)_{0}$, $(A_V,A^*_V)_{0}$, $(A_\perp,A^*_\perp)_{0}$, $(\lambda_+,\lambda^*_+)_{0}$, $(\lambda_-,\lambda^*_-)_{0}$ \\
 &  & 1 & 2 & $(c,c^*)_{1}$, $(A_V,A^*_V)_{1}$, $(A_\perp,A^*_\perp)_{1}$, $(\lambda_+,\lambda^*_+)_{1}$, $(\lambda_-,\lambda^*_-)_{1}$ \\
 &  & 1 & 3 & $(c,A_V,c^*)_{0}$, $(c,\lambda^*_+,\lambda^*_-)_{0}$, $(A_V,A_V,A^*_V)_{0}$, $(A_V,A_\perp,A^*_\perp)_{0}$, $(A_V,\lambda_+,\lambda^*_+)_{0}$, $(A_V,\lambda_-,\lambda^*_-)_{0}$, $(A_\perp,A_\perp,A^*_V)_{0}$, $(A_\perp,\lambda_+,\lambda^*_-)_{0}$, $(A_\perp,\lambda_-,\lambda^*_+)_{0}$ \\
\hline
$\Ob^{(1)}(R_r)$ & $4$ & 0 & 2 & $(c,A^*_V)_{1}$, $(A_V,A_V)_{2}$, $(A_\perp,A_\perp)_{2}$, $(\lambda_+,\lambda_-)_{1}$ \\
 &  & 0 & 3 & $(c,c,c^*)_{0}$, $(c,A_V,A^*_V)_{0}$, $(c,A_\perp,A^*_\perp)_{0}$, $(c,\lambda_+,\lambda^*_+)_{0}$, $(c,\lambda_-,\lambda^*_-)_{0}$, $(A_V,A_V,A_V)_{1}$, $(A_V,A_\perp,A_\perp)_{1}$, $(A_V,\lambda_+,\lambda_-)_{0}$, $(A_\perp,\lambda_+,\lambda_+)_{0}$, $(A_\perp,\lambda_-,\lambda_-)_{0}$ \\
 &  & 0 & 4 & $(A_V,A_V,A_V,A_V)_{0}$, $(A_V,A_V,A_\perp,A_\perp)_{0}$, $(A_\perp,A_\perp,A_\perp,A_\perp)_{0}$ \\
\hline
$F_1^{(1)}(R_r)$ & $4$ & 0 & 2 & $(c,c^*)_{0}$, $(A_V,A^*_V)_{0}$, $(A_\perp,A^*_\perp)_{0}$, $(\lambda_+,\lambda^*_+)_{0}$, $(\lambda_-,\lambda^*_-)_{0}$ \\
\hline
$\Ob^{(1)}(\xi_+)$ & $9/2$ & 0 & 2 & $(c,\lambda^*_+)_{2}$, $(A_V,\lambda_-)_{2}$, $(A_\perp,\lambda_+)_{2}$ \\
 &  & 0 & 3 & $(c,A_V,\lambda^*_+)_{1}$, $(c,A_\perp,\lambda^*_-)_{1}$, $(c,\lambda_+,A^*_\perp)_{0}$, $(c,\lambda_-,A^*_V)_{0}$, $(A_V,A_V,\lambda_-)_{1}$, $(A_V,A_\perp,\lambda_+)_{1}$, $(A_\perp,A_\perp,\lambda_-)_{1}$, $(\lambda_+,\lambda_+,\lambda_+)_{0}$, $(\lambda_+,\lambda_-,\lambda_-)_{0}$ \\
 &  & 0 & 4 & $(c,A_V,A_V,\lambda^*_+)_{0}$, $(c,A_V,A_\perp,\lambda^*_-)_{0}$, $(c,A_\perp,A_\perp,\lambda^*_+)_{0}$, $(A_V,A_V,A_V,\lambda_-)_{0}$, $(A_V,A_V,A_\perp,\lambda_+)_{0}$, $(A_V,A_\perp,A_\perp,\lambda_-)_{0}$, $(A_\perp,A_\perp,A_\perp,\lambda_+)_{0}$ \\
\hline
$F_1^{(1)}(\xi_+)$ & $9/2$ & 0 & 2 & $(A_V,\lambda^*_+)_{1}$, $(A_\perp,\lambda^*_-)_{1}$, $(\lambda_+,A^*_\perp)_{0}$, $(\lambda_-,A^*_V)_{0}$ \\
 &  & 0 & 3 & $(A_V,A_V,\lambda^*_+)_{0}$, $(A_V,A_\perp,\lambda^*_-)_{0}$, $(A_\perp,A_\perp,\lambda^*_+)_{0}$ \\
\hline
$\Ob^{(1)}(\xi_-)$ & $9/2$ & 0 & 2 & $(c,\lambda^*_-)_{2}$, $(A_V,\lambda_+)_{2}$, $(A_\perp,\lambda_-)_{2}$ \\
 &  & 0 & 3 & $(c,A_V,\lambda^*_-)_{1}$, $(c,A_\perp,\lambda^*_+)_{1}$, $(c,\lambda_+,A^*_V)_{0}$, $(c,\lambda_-,A^*_\perp)_{0}$, $(A_V,A_V,\lambda_+)_{1}$, $(A_V,A_\perp,\lambda_-)_{1}$, $(A_\perp,A_\perp,\lambda_+)_{1}$, $(\lambda_+,\lambda_+,\lambda_-)_{0}$, $(\lambda_-,\lambda_-,\lambda_-)_{0}$ \\
 &  & 0 & 4 & $(c,A_V,A_V,\lambda^*_-)_{0}$, $(c,A_V,A_\perp,\lambda^*_+)_{0}$, $(c,A_\perp,A_\perp,\lambda^*_-)_{0}$, $(A_V,A_V,A_V,\lambda_+)_{0}$, $(A_V,A_V,A_\perp,\lambda_-)_{0}$, $(A_V,A_\perp,A_\perp,\lambda_+)_{0}$, $(A_\perp,A_\perp,A_\perp,\lambda_-)_{0}$ \\
\hline
$F_1^{(1)}(\xi_-)$ & $9/2$ & 0 & 2 & $(A_V,\lambda^*_-)_{1}$, $(A_\perp,\lambda^*_+)_{1}$, $(\lambda_+,A^*_V)_{0}$, $(\lambda_-,A^*_\perp)_{0}$ \\
 &  & 0 & 3 & $(A_V,A_V,\lambda^*_-)_{0}$, $(A_V,A_\perp,\lambda^*_+)_{0}$, $(A_\perp,A_\perp,\lambda^*_-)_{0}$ \\
\hline
$\Ob^{(1)}(\eta_+)$ & $7/2$ & 0 & 2 & $(c,\lambda^*_-)_{1}$, $(A_V,\lambda_+)_{1}$, $(A_\perp,\lambda_-)_{1}$ \\
 &  & 0 & 3 & $(c,A_V,\lambda^*_-)_{0}$, $(c,A_\perp,\lambda^*_+)_{0}$, $(A_V,A_V,\lambda_+)_{0}$, $(A_V,A_\perp,\lambda_-)_{0}$, $(A_\perp,A_\perp,\lambda_+)_{0}$ \\
 &  & 1 & 2 & $(c,\lambda^*_-)_{2}$, $(A_V,\lambda_+)_{2}$, $(A_\perp,\lambda_-)_{2}$ \\
 &  & 1 & 3 & $(c,A_V,\lambda^*_-)_{1}$, $(c,A_\perp,\lambda^*_+)_{1}$, $(c,\lambda_+,A^*_V)_{0}$, $(c,\lambda_-,A^*_\perp)_{0}$, $(A_V,A_V,\lambda_+)_{1}$, $(A_V,A_\perp,\lambda_-)_{1}$, $(A_\perp,A_\perp,\lambda_+)_{1}$, $(\lambda_+,\lambda_+,\lambda_-)_{0}$, $(\lambda_-,\lambda_-,\lambda_-)_{0}$ \\
 &  & 1 & 4 & $(c,A_V,A_V,\lambda^*_-)_{0}$, $(c,A_V,A_\perp,\lambda^*_+)_{0}$, $(c,A_\perp,A_\perp,\lambda^*_-)_{0}$, $(A_V,A_V,A_V,\lambda_+)_{0}$, $(A_V,A_V,A_\perp,\lambda_-)_{0}$, $(A_V,A_\perp,A_\perp,\lambda_+)_{0}$, $(A_\perp,A_\perp,A_\perp,\lambda_-)_{0}$ \\
\hline
$F_1^{(1)}(\eta_+)$ & $7/2$ & 0 & 2 & $(A_V,\lambda^*_-)_{0}$, $(A_\perp,\lambda^*_+)_{0}$ \\
 &  & 1 & 2 & $(A_V,\lambda^*_-)_{1}$, $(A_\perp,\lambda^*_+)_{1}$, $(\lambda_+,A^*_V)_{0}$, $(\lambda_-,A^*_\perp)_{0}$ \\
 &  & 1 & 3 & $(A_V,A_V,\lambda^*_-)_{0}$, $(A_V,A_\perp,\lambda^*_+)_{0}$, $(A_\perp,A_\perp,\lambda^*_-)_{0}$ \\
\hline
$\Ob^{(1)}(\eta_-)$ & $7/2$ & 0 & 2 & $(c,\lambda^*_+)_{1}$, $(A_V,\lambda_-)_{1}$, $(A_\perp,\lambda_+)_{1}$ \\
 &  & 0 & 3 & $(c,A_V,\lambda^*_+)_{0}$, $(c,A_\perp,\lambda^*_-)_{0}$, $(A_V,A_V,\lambda_-)_{0}$, $(A_V,A_\perp,\lambda_+)_{0}$, $(A_\perp,A_\perp,\lambda_-)_{0}$ \\
 &  & 1 & 2 & $(c,\lambda^*_+)_{2}$, $(A_V,\lambda_-)_{2}$, $(A_\perp,\lambda_+)_{2}$ \\
 &  & 1 & 3 & $(c,A_V,\lambda^*_+)_{1}$, $(c,A_\perp,\lambda^*_-)_{1}$, $(c,\lambda_+,A^*_\perp)_{0}$, $(c,\lambda_-,A^*_V)_{0}$, $(A_V,A_V,\lambda_-)_{1}$, $(A_V,A_\perp,\lambda_+)_{1}$, $(A_\perp,A_\perp,\lambda_-)_{1}$, $(\lambda_+,\lambda_+,\lambda_+)_{0}$, $(\lambda_+,\lambda_-,\lambda_-)_{0}$ \\
 &  & 1 & 4 & $(c,A_V,A_V,\lambda^*_+)_{0}$, $(c,A_V,A_\perp,\lambda^*_-)_{0}$, $(c,A_\perp,A_\perp,\lambda^*_+)_{0}$, $(A_V,A_V,A_V,\lambda_-)_{0}$, $(A_V,A_V,A_\perp,\lambda_+)_{0}$, $(A_V,A_\perp,A_\perp,\lambda_-)_{0}$, $(A_\perp,A_\perp,A_\perp,\lambda_+)_{0}$ \\
\hline
$F_1^{(1)}(\eta_-)$ & $7/2$ & 0 & 2 & $(A_V,\lambda^*_+)_{0}$, $(A_\perp,\lambda^*_-)_{0}$ \\
 &  & 1 & 2 & $(A_V,\lambda^*_+)_{1}$, $(A_\perp,\lambda^*_-)_{1}$, $(\lambda_+,A^*_\perp)_{0}$, $(\lambda_-,A^*_V)_{0}$ \\
 &  & 1 & 3 & $(A_V,A_V,\lambda^*_+)_{0}$, $(A_V,A_\perp,\lambda^*_-)_{0}$, $(A_\perp,A_\perp,\lambda^*_+)_{0}$ \\
\hline
$\Ob^{(1)}(K_b)$ & $3$ & 0 & 2 & $(c,A^*_V)_{0}$, $(A_V,A_V)_{1}$, $(A_\perp,A_\perp)_{1}$, $(\lambda_+,\lambda_-)_{0}$ \\
 &  & 0 & 3 & $(A_V,A_V,A_V)_{0}$, $(A_V,A_\perp,A_\perp)_{0}$ \\
 &  & 1 & 2 & $(c,A^*_V)_{1}$, $(A_V,A_V)_{2}$, $(A_\perp,A_\perp)_{2}$, $(\lambda_+,\lambda_-)_{1}$ \\
 &  & 1 & 3 & $(c,c,c^*)_{0}$, $(c,A_V,A^*_V)_{0}$, $(c,A_\perp,A^*_\perp)_{0}$, $(c,\lambda_+,\lambda^*_+)_{0}$, $(c,\lambda_-,\lambda^*_-)_{0}$, $(A_V,A_V,A_V)_{1}$, $(A_V,A_\perp,A_\perp)_{1}$, $(A_V,\lambda_+,\lambda_-)_{0}$, $(A_\perp,\lambda_+,\lambda_+)_{0}$, $(A_\perp,\lambda_-,\lambda_-)_{0}$ \\
 &  & 1 & 4 & $(A_V,A_V,A_V,A_V)_{0}$, $(A_V,A_V,A_\perp,A_\perp)_{0}$, $(A_\perp,A_\perp,A_\perp,A_\perp)_{0}$ \\
 &  & 2 & 2 & $(c,A^*_V)_{2}$, $(A_V,A_V)_{3}$, $(A_\perp,A_\perp)_{3}$, $(\lambda_+,\lambda_-)_{2}$ \\
 &  & 2 & 3 & $(c,c,c^*)_{1}$, $(c,A_V,A^*_V)_{1}$, $(c,A_\perp,A^*_\perp)_{1}$, $(c,\lambda_+,\lambda^*_+)_{1}$, $(c,\lambda_-,\lambda^*_-)_{1}$, $(A_V,A_V,A_V)_{2}$, $(A_V,A_\perp,A_\perp)_{2}$, $(A_V,\lambda_+,\lambda_-)_{1}$, $(A_\perp,\lambda_+,\lambda_+)_{1}$, $(A_\perp,\lambda_-,\lambda_-)_{1}$ \\
 &  & 2 & 4 & $(c,c,A_V,c^*)_{0}$, $(c,c,\lambda^*_+,\lambda^*_-)_{0}$, $(c,A_V,A_V,A^*_V)_{0}$, $(c,A_V,A_\perp,A^*_\perp)_{0}$, $(c,A_V,\lambda_+,\lambda^*_+)_{0}$, $(c,A_V,\lambda_-,\lambda^*_-)_{0}$, $(c,A_\perp,A_\perp,A^*_V)_{0}$, $(c,A_\perp,\lambda_+,\lambda^*_-)_{0}$, $(c,A_\perp,\lambda_-,\lambda^*_+)_{0}$, $(A_V,A_V,A_V,A_V)_{1}$, $(A_V,A_V,A_\perp,A_\perp)_{1}$, $(A_V,A_V,\lambda_+,\lambda_-)_{0}$, $(A_V,A_\perp,\lambda_+,\lambda_+)_{0}$, $(A_V,A_\perp,\lambda_-,\lambda_-)_{0}$, $(A_\perp,A_\perp,A_\perp,A_\perp)_{1}$, $(A_\perp,A_\perp,\lambda_+,\lambda_-)_{0}$ \\
 &  & 2 & 5 & $(A_V,A_V,A_V,A_V,A_V)_{0}$, $(A_V,A_V,A_V,A_\perp,A_\perp)_{0}$, $(A_V,A_\perp,A_\perp,A_\perp,A_\perp)_{0}$ \\
\hline
$F_1^{(1)}(K_b)$ & $3$ & 1 & 2 & $(c,c^*)_{0}$, $(A_V,A^*_V)_{0}$, $(A_\perp,A^*_\perp)_{0}$, $(\lambda_+,\lambda^*_+)_{0}$, $(\lambda_-,\lambda^*_-)_{0}$ \\
 &  & 2 & 2 & $(c,c^*)_{1}$, $(A_V,A^*_V)_{1}$, $(A_\perp,A^*_\perp)_{1}$, $(\lambda_+,\lambda^*_+)_{1}$, $(\lambda_-,\lambda^*_-)_{1}$ \\
 &  & 2 & 3 & $(c,A_V,c^*)_{0}$, $(c,\lambda^*_+,\lambda^*_-)_{0}$, $(A_V,A_V,A^*_V)_{0}$, $(A_V,A_\perp,A^*_\perp)_{0}$, $(A_V,\lambda_+,\lambda^*_+)_{0}$, $(A_V,\lambda_-,\lambda^*_-)_{0}$, $(A_\perp,A_\perp,A^*_V)_{0}$, $(A_\perp,\lambda_+,\lambda^*_-)_{0}$, $(A_\perp,\lambda_-,\lambda^*_+)_{0}$ \\
\hline

\caption{Census of the leg sets of the arity-one components allowed by Theorem \ref{thm:at-bound}: $(\cdots)_D$ is a leg set with $D$ derivatives, $p$ the degree in the position; $\lambda_\pm$, $\lambda^*_\pm$, $\xi_\pm$, $\eta_\pm$ as in Definition \ref{def:at-gens}. Necessary conditions; the reflection parity and the symmetry of identical legs remove further structures inside each leg set. $\mathrm{F}^{(1)}$ stands for $F_1^{(1)}$.}
\label{tab:at-census1}
\end{longtable}

\begin{longtable}{>{\raggedright}p{2.3cm}|p{0.6cm}|p{0.25cm}|p{0.25cm}|>{\raggedright\arraybackslash}p{11.6cm}}
cochain & weight & $p$ & $n$ & leg sets $(\cdots)_D$\\
\hline
\endhead
$\Ob^{(2)}(\Delta,P_b)$ & $5$ & 0 & 2 & $(c,c^*)_{1}$, $(A_V,A^*_V)_{1}$, $(A_\perp,A^*_\perp)_{1}$, $(\lambda_+,\lambda^*_+)_{1}$, $(\lambda_-,\lambda^*_-)_{1}$ \\
 &  & 0 & 3 & $(c,A_V,c^*)_{0}$, $(c,\lambda^*_+,\lambda^*_-)_{0}$, $(A_V,A_V,A^*_V)_{0}$, $(A_V,A_\perp,A^*_\perp)_{0}$, $(A_V,\lambda_+,\lambda^*_+)_{0}$, $(A_V,\lambda_-,\lambda^*_-)_{0}$, $(A_\perp,A_\perp,A^*_V)_{0}$, $(A_\perp,\lambda_+,\lambda^*_-)_{0}$, $(A_\perp,\lambda_-,\lambda^*_+)_{0}$ \\
 &  & 1 & 2 & $(c,c^*)_{2}$, $(A_V,A^*_V)_{2}$, $(A_\perp,A^*_\perp)_{2}$, $(\lambda_+,\lambda^*_+)_{2}$, $(\lambda_-,\lambda^*_-)_{2}$ \\
 &  & 1 & 3 & $(c,A_V,c^*)_{1}$, $(c,A^*_V,A^*_V)_{0}$, $(c,A^*_\perp,A^*_\perp)_{0}$, $(c,\lambda^*_+,\lambda^*_-)_{1}$, $(A_V,A_V,A^*_V)_{1}$, $(A_V,A_\perp,A^*_\perp)_{1}$, $(A_V,\lambda_+,\lambda^*_+)_{1}$, $(A_V,\lambda_-,\lambda^*_-)_{1}$, $(A_\perp,A_\perp,A^*_V)_{1}$, $(A_\perp,\lambda_+,\lambda^*_-)_{1}$, $(A_\perp,\lambda_-,\lambda^*_+)_{1}$, $(\lambda_+,\lambda_+,A^*_\perp)_{0}$, $(\lambda_+,\lambda_-,A^*_V)_{0}$, $(\lambda_-,\lambda_-,A^*_\perp)_{0}$ \\
 &  & 1 & 4 & $(c,A_V,A_V,c^*)_{0}$, $(c,A_V,\lambda^*_+,\lambda^*_-)_{0}$, $(c,A_\perp,A_\perp,c^*)_{0}$, $(c,A_\perp,\lambda^*_+,\lambda^*_+)_{0}$, $(c,A_\perp,\lambda^*_-,\lambda^*_-)_{0}$, $(A_V,A_V,A_V,A^*_V)_{0}$, $(A_V,A_V,A_\perp,A^*_\perp)_{0}$, $(A_V,A_V,\lambda_+,\lambda^*_+)_{0}$, $(A_V,A_V,\lambda_-,\lambda^*_-)_{0}$, $(A_V,A_\perp,A_\perp,A^*_V)_{0}$, $(A_V,A_\perp,\lambda_+,\lambda^*_-)_{0}$, $(A_V,A_\perp,\lambda_-,\lambda^*_+)_{0}$, $(A_\perp,A_\perp,A_\perp,A^*_\perp)_{0}$, $(A_\perp,A_\perp,\lambda_+,\lambda^*_+)_{0}$, $(A_\perp,A_\perp,\lambda_-,\lambda^*_-)_{0}$ \\
\hline
$F_1^{(2)}(\Delta,P_b)$ & $5$ & 0 & 2 & $(A_V,c^*)_{0}$, $(\lambda^*_+,\lambda^*_-)_{0}$ \\
 &  & 1 & 2 & $(A_V,c^*)_{1}$, $(A^*_V,A^*_V)_{0}$, $(A^*_\perp,A^*_\perp)_{0}$, $(\lambda^*_+,\lambda^*_-)_{1}$ \\
 &  & 1 & 3 & $(A_V,A_V,c^*)_{0}$, $(A_V,\lambda^*_+,\lambda^*_-)_{0}$, $(A_\perp,A_\perp,c^*)_{0}$, $(A_\perp,\lambda^*_+,\lambda^*_+)_{0}$, $(A_\perp,\lambda^*_-,\lambda^*_-)_{0}$ \\
\hline
$\Ob^{(2)}(K_b,P_b)$ & $4$ & 0 & 2 & $(c,c^*)_{0}$, $(A_V,A^*_V)_{0}$, $(A_\perp,A^*_\perp)_{0}$, $(\lambda_+,\lambda^*_+)_{0}$, $(\lambda_-,\lambda^*_-)_{0}$ \\
 &  & 1 & 2 & $(c,c^*)_{1}$, $(A_V,A^*_V)_{1}$, $(A_\perp,A^*_\perp)_{1}$, $(\lambda_+,\lambda^*_+)_{1}$, $(\lambda_-,\lambda^*_-)_{1}$ \\
 &  & 1 & 3 & $(c,A_V,c^*)_{0}$, $(c,\lambda^*_+,\lambda^*_-)_{0}$, $(A_V,A_V,A^*_V)_{0}$, $(A_V,A_\perp,A^*_\perp)_{0}$, $(A_V,\lambda_+,\lambda^*_+)_{0}$, $(A_V,\lambda_-,\lambda^*_-)_{0}$, $(A_\perp,A_\perp,A^*_V)_{0}$, $(A_\perp,\lambda_+,\lambda^*_-)_{0}$, $(A_\perp,\lambda_-,\lambda^*_+)_{0}$ \\
 &  & 2 & 2 & $(c,c^*)_{2}$, $(A_V,A^*_V)_{2}$, $(A_\perp,A^*_\perp)_{2}$, $(\lambda_+,\lambda^*_+)_{2}$, $(\lambda_-,\lambda^*_-)_{2}$ \\
 &  & 2 & 3 & $(c,A_V,c^*)_{1}$, $(c,A^*_V,A^*_V)_{0}$, $(c,A^*_\perp,A^*_\perp)_{0}$, $(c,\lambda^*_+,\lambda^*_-)_{1}$, $(A_V,A_V,A^*_V)_{1}$, $(A_V,A_\perp,A^*_\perp)_{1}$, $(A_V,\lambda_+,\lambda^*_+)_{1}$, $(A_V,\lambda_-,\lambda^*_-)_{1}$, $(A_\perp,A_\perp,A^*_V)_{1}$, $(A_\perp,\lambda_+,\lambda^*_-)_{1}$, $(A_\perp,\lambda_-,\lambda^*_+)_{1}$, $(\lambda_+,\lambda_+,A^*_\perp)_{0}$, $(\lambda_+,\lambda_-,A^*_V)_{0}$, $(\lambda_-,\lambda_-,A^*_\perp)_{0}$ \\
 &  & 2 & 4 & $(c,A_V,A_V,c^*)_{0}$, $(c,A_V,\lambda^*_+,\lambda^*_-)_{0}$, $(c,A_\perp,A_\perp,c^*)_{0}$, $(c,A_\perp,\lambda^*_+,\lambda^*_+)_{0}$, $(c,A_\perp,\lambda^*_-,\lambda^*_-)_{0}$, $(A_V,A_V,A_V,A^*_V)_{0}$, $(A_V,A_V,A_\perp,A^*_\perp)_{0}$, $(A_V,A_V,\lambda_+,\lambda^*_+)_{0}$, $(A_V,A_V,\lambda_-,\lambda^*_-)_{0}$, $(A_V,A_\perp,A_\perp,A^*_V)_{0}$, $(A_V,A_\perp,\lambda_+,\lambda^*_-)_{0}$, $(A_V,A_\perp,\lambda_-,\lambda^*_+)_{0}$, $(A_\perp,A_\perp,A_\perp,A^*_\perp)_{0}$, $(A_\perp,A_\perp,\lambda_+,\lambda^*_+)_{0}$, $(A_\perp,A_\perp,\lambda_-,\lambda^*_-)_{0}$ \\
\hline
$F_1^{(2)}(K_b,P_b)$ & $4$ & 1 & 2 & $(A_V,c^*)_{0}$, $(\lambda^*_+,\lambda^*_-)_{0}$ \\
 &  & 2 & 2 & $(A_V,c^*)_{1}$, $(A^*_V,A^*_V)_{0}$, $(A^*_\perp,A^*_\perp)_{0}$, $(\lambda^*_+,\lambda^*_-)_{1}$ \\
 &  & 2 & 3 & $(A_V,A_V,c^*)_{0}$, $(A_V,\lambda^*_+,\lambda^*_-)_{0}$, $(A_\perp,A_\perp,c^*)_{0}$, $(A_\perp,\lambda^*_+,\lambda^*_+)_{0}$, $(A_\perp,\lambda^*_-,\lambda^*_-)_{0}$ \\
\hline
$\Ob^{(2)}(\xi_+,\xi_-)$ & $5$ & 0 & 2 & $(c,c^*)_{1}$, $(A_V,A^*_V)_{1}$, $(A_\perp,A^*_\perp)_{1}$, $(\lambda_+,\lambda^*_+)_{1}$, $(\lambda_-,\lambda^*_-)_{1}$ \\
 &  & 0 & 3 & $(c,A_V,c^*)_{0}$, $(c,\lambda^*_+,\lambda^*_-)_{0}$, $(A_V,A_V,A^*_V)_{0}$, $(A_V,A_\perp,A^*_\perp)_{0}$, $(A_V,\lambda_+,\lambda^*_+)_{0}$, $(A_V,\lambda_-,\lambda^*_-)_{0}$, $(A_\perp,A_\perp,A^*_V)_{0}$, $(A_\perp,\lambda_+,\lambda^*_-)_{0}$, $(A_\perp,\lambda_-,\lambda^*_+)_{0}$ \\
\hline
$F_1^{(2)}(\xi_+,\xi_-)$ & $5$ & 0 & 2 & $(A_V,c^*)_{0}$, $(\lambda^*_+,\lambda^*_-)_{0}$ \\
\hline
$\Ob^{(2)}(\xi_+,\xi_+)$ & $5$ & 0 & 2 & $(A_V,A^*_\perp)_{1}$, $(A_\perp,A^*_V)_{1}$, $(\lambda_+,\lambda^*_-)_{1}$, $(\lambda_-,\lambda^*_+)_{1}$ \\
 &  & 0 & 3 & $(c,A_\perp,c^*)_{0}$, $(c,\lambda^*_+,\lambda^*_+)_{0}$, $(c,\lambda^*_-,\lambda^*_-)_{0}$, $(A_V,A_V,A^*_\perp)_{0}$, $(A_V,A_\perp,A^*_V)_{0}$, $(A_V,\lambda_+,\lambda^*_-)_{0}$, $(A_V,\lambda_-,\lambda^*_+)_{0}$, $(A_\perp,A_\perp,A^*_\perp)_{0}$, $(A_\perp,\lambda_+,\lambda^*_+)_{0}$, $(A_\perp,\lambda_-,\lambda^*_-)_{0}$ \\
\hline
$F_1^{(2)}(\xi_+,\xi_+)$ & $5$ & 0 & 2 & $(A_\perp,c^*)_{0}$, $(\lambda^*_+,\lambda^*_+)_{0}$, $(\lambda^*_-,\lambda^*_-)_{0}$ \\
\hline
$\Ob^{(2)}(\xi_+,\eta_+)$ & $4$ & 0 & 2 & $(c,c^*)_{0}$, $(A_V,A^*_V)_{0}$, $(A_\perp,A^*_\perp)_{0}$, $(\lambda_+,\lambda^*_+)_{0}$, $(\lambda_-,\lambda^*_-)_{0}$ \\
 &  & 1 & 2 & $(c,c^*)_{1}$, $(A_V,A^*_V)_{1}$, $(A_\perp,A^*_\perp)_{1}$, $(\lambda_+,\lambda^*_+)_{1}$, $(\lambda_-,\lambda^*_-)_{1}$ \\
 &  & 1 & 3 & $(c,A_V,c^*)_{0}$, $(c,\lambda^*_+,\lambda^*_-)_{0}$, $(A_V,A_V,A^*_V)_{0}$, $(A_V,A_\perp,A^*_\perp)_{0}$, $(A_V,\lambda_+,\lambda^*_+)_{0}$, $(A_V,\lambda_-,\lambda^*_-)_{0}$, $(A_\perp,A_\perp,A^*_V)_{0}$, $(A_\perp,\lambda_+,\lambda^*_-)_{0}$, $(A_\perp,\lambda_-,\lambda^*_+)_{0}$ \\
\hline
$F_1^{(2)}(\xi_+,\eta_+)$ & $4$ & 1 & 2 & $(A_V,c^*)_{0}$, $(\lambda^*_+,\lambda^*_-)_{0}$ \\
\hline
$\Ob^{(2)}(\xi_+,\eta_-)$ & $4$ & 0 & 2 & $(A_V,A^*_\perp)_{0}$, $(A_\perp,A^*_V)_{0}$, $(\lambda_+,\lambda^*_-)_{0}$, $(\lambda_-,\lambda^*_+)_{0}$ \\
 &  & 1 & 2 & $(A_V,A^*_\perp)_{1}$, $(A_\perp,A^*_V)_{1}$, $(\lambda_+,\lambda^*_-)_{1}$, $(\lambda_-,\lambda^*_+)_{1}$ \\
 &  & 1 & 3 & $(c,A_\perp,c^*)_{0}$, $(c,\lambda^*_+,\lambda^*_+)_{0}$, $(c,\lambda^*_-,\lambda^*_-)_{0}$, $(A_V,A_V,A^*_\perp)_{0}$, $(A_V,A_\perp,A^*_V)_{0}$, $(A_V,\lambda_+,\lambda^*_-)_{0}$, $(A_V,\lambda_-,\lambda^*_+)_{0}$, $(A_\perp,A_\perp,A^*_\perp)_{0}$, $(A_\perp,\lambda_+,\lambda^*_+)_{0}$, $(A_\perp,\lambda_-,\lambda^*_-)_{0}$ \\
\hline
$F_1^{(2)}(\xi_+,\eta_-)$ & $4$ & 1 & 2 & $(A_\perp,c^*)_{0}$, $(\lambda^*_+,\lambda^*_+)_{0}$, $(\lambda^*_-,\lambda^*_-)_{0}$ \\
\hline
$\Ob^{(2)}(\xi_+,P_b)$ & $11/2$ & 0 & 2 & $(A_V,\lambda^*_+)_{2}$, $(A_\perp,\lambda^*_-)_{2}$, $(\lambda_+,A^*_\perp)_{1}$, $(\lambda_-,A^*_V)_{1}$ \\
 &  & 0 & 3 & $(c,\lambda_-,c^*)_{0}$, $(c,A^*_V,\lambda^*_+)_{0}$, $(c,A^*_\perp,\lambda^*_-)_{0}$, $(A_V,A_V,\lambda^*_+)_{1}$, $(A_V,A_\perp,\lambda^*_-)_{1}$, $(A_V,\lambda_+,A^*_\perp)_{0}$, $(A_V,\lambda_-,A^*_V)_{0}$, $(A_\perp,A_\perp,\lambda^*_+)_{1}$, $(A_\perp,\lambda_+,A^*_V)_{0}$, $(A_\perp,\lambda_-,A^*_\perp)_{0}$, $(\lambda_+,\lambda_+,\lambda^*_-)_{0}$, $(\lambda_+,\lambda_-,\lambda^*_+)_{0}$, $(\lambda_-,\lambda_-,\lambda^*_-)_{0}$ \\
 &  & 0 & 4 & $(A_V,A_V,A_V,\lambda^*_+)_{0}$, $(A_V,A_V,A_\perp,\lambda^*_-)_{0}$, $(A_V,A_\perp,A_\perp,\lambda^*_+)_{0}$, $(A_\perp,A_\perp,A_\perp,\lambda^*_-)_{0}$ \\
\hline
$F_1^{(2)}(\xi_+,P_b)$ & $11/2$ & 0 & 2 & $(\lambda_-,c^*)_{0}$, $(A^*_V,\lambda^*_+)_{0}$, $(A^*_\perp,\lambda^*_-)_{0}$ \\
\hline
$\Ob^{(2)}(\eta_+,\eta_-)$ & $3$ & 1 & 2 & $(c,c^*)_{0}$, $(A_V,A^*_V)_{0}$, $(A_\perp,A^*_\perp)_{0}$, $(\lambda_+,\lambda^*_+)_{0}$, $(\lambda_-,\lambda^*_-)_{0}$ \\
 &  & 2 & 2 & $(c,c^*)_{1}$, $(A_V,A^*_V)_{1}$, $(A_\perp,A^*_\perp)_{1}$, $(\lambda_+,\lambda^*_+)_{1}$, $(\lambda_-,\lambda^*_-)_{1}$ \\
 &  & 2 & 3 & $(c,A_V,c^*)_{0}$, $(c,\lambda^*_+,\lambda^*_-)_{0}$, $(A_V,A_V,A^*_V)_{0}$, $(A_V,A_\perp,A^*_\perp)_{0}$, $(A_V,\lambda_+,\lambda^*_+)_{0}$, $(A_V,\lambda_-,\lambda^*_-)_{0}$, $(A_\perp,A_\perp,A^*_V)_{0}$, $(A_\perp,\lambda_+,\lambda^*_-)_{0}$, $(A_\perp,\lambda_-,\lambda^*_+)_{0}$ \\
\hline
$F_1^{(2)}(\eta_+,\eta_-)$ & $3$ & 2 & 2 & $(A_V,c^*)_{0}$, $(\lambda^*_+,\lambda^*_-)_{0}$ \\
\hline
$\Ob^{(2)}(\eta_+,\eta_+)$ & $3$ & 1 & 2 & $(A_V,A^*_\perp)_{0}$, $(A_\perp,A^*_V)_{0}$, $(\lambda_+,\lambda^*_-)_{0}$, $(\lambda_-,\lambda^*_+)_{0}$ \\
 &  & 2 & 2 & $(A_V,A^*_\perp)_{1}$, $(A_\perp,A^*_V)_{1}$, $(\lambda_+,\lambda^*_-)_{1}$, $(\lambda_-,\lambda^*_+)_{1}$ \\
 &  & 2 & 3 & $(c,A_\perp,c^*)_{0}$, $(c,\lambda^*_+,\lambda^*_+)_{0}$, $(c,\lambda^*_-,\lambda^*_-)_{0}$, $(A_V,A_V,A^*_\perp)_{0}$, $(A_V,A_\perp,A^*_V)_{0}$, $(A_V,\lambda_+,\lambda^*_-)_{0}$, $(A_V,\lambda_-,\lambda^*_+)_{0}$, $(A_\perp,A_\perp,A^*_\perp)_{0}$, $(A_\perp,\lambda_+,\lambda^*_+)_{0}$, $(A_\perp,\lambda_-,\lambda^*_-)_{0}$ \\
\hline
$F_1^{(2)}(\eta_+,\eta_+)$ & $3$ & 2 & 2 & $(A_\perp,c^*)_{0}$, $(\lambda^*_+,\lambda^*_+)_{0}$, $(\lambda^*_-,\lambda^*_-)_{0}$ \\
\hline
$\Ob^{(2)}(\Delta,\xi_+)$ & $9/2$ & 0 & 2 & $(A_V,\lambda^*_+)_{1}$, $(A_\perp,\lambda^*_-)_{1}$, $(\lambda_+,A^*_\perp)_{0}$, $(\lambda_-,A^*_V)_{0}$ \\
 &  & 0 & 3 & $(A_V,A_V,\lambda^*_+)_{0}$, $(A_V,A_\perp,\lambda^*_-)_{0}$, $(A_\perp,A_\perp,\lambda^*_+)_{0}$ \\
 &  & 1 & 2 & $(A_V,\lambda^*_+)_{2}$, $(A_\perp,\lambda^*_-)_{2}$, $(\lambda_+,A^*_\perp)_{1}$, $(\lambda_-,A^*_V)_{1}$ \\
 &  & 1 & 3 & $(c,\lambda_-,c^*)_{0}$, $(c,A^*_V,\lambda^*_+)_{0}$, $(c,A^*_\perp,\lambda^*_-)_{0}$, $(A_V,A_V,\lambda^*_+)_{1}$, $(A_V,A_\perp,\lambda^*_-)_{1}$, $(A_V,\lambda_+,A^*_\perp)_{0}$, $(A_V,\lambda_-,A^*_V)_{0}$, $(A_\perp,A_\perp,\lambda^*_+)_{1}$, $(A_\perp,\lambda_+,A^*_V)_{0}$, $(A_\perp,\lambda_-,A^*_\perp)_{0}$, $(\lambda_+,\lambda_+,\lambda^*_-)_{0}$, $(\lambda_+,\lambda_-,\lambda^*_+)_{0}$, $(\lambda_-,\lambda_-,\lambda^*_-)_{0}$ \\
 &  & 1 & 4 & $(A_V,A_V,A_V,\lambda^*_+)_{0}$, $(A_V,A_V,A_\perp,\lambda^*_-)_{0}$, $(A_V,A_\perp,A_\perp,\lambda^*_+)_{0}$, $(A_\perp,A_\perp,A_\perp,\lambda^*_-)_{0}$ \\
\hline
$F_1^{(2)}(\Delta,\xi_+)$ & $9/2$ & 1 & 2 & $(\lambda_-,c^*)_{0}$, $(A^*_V,\lambda^*_+)_{0}$, $(A^*_\perp,\lambda^*_-)_{0}$ \\
\hline

\caption{Census of the leg sets of the arity-two components, as in Table \ref{tab:at-census1}. The pair $(\xi_+,\xi_-)$ has $\Gamma(\xi_+,\xi_-)|_V\neq0$, the pair $(\xi_+,\xi'_+)$ has bracket $0$; $(\xi_+,\eta_+)$ contains $\Delta$, $\mathfrak{so}(V)$ and $\mathfrak{so}(W^\perp_{\mathbb R})$ in its bracket, $(\xi_+,\eta_-)$ has bracket $0$; $(\eta_+,\eta_-)$ has bracket $K_{\Gamma(\eta_+,\eta_-)|_V}$.}
\label{tab:at-census2}
\end{longtable}
}

\begin{proposition}[Endpoint]
\label{prop:at-T}
The components $\Ob^{(1)}(x)$, $\Ob^{(2)}(x,y)$, $\lc^{(k)}$ and the solutions $F_1^{(k)}$ are independent of $T$. Along $\Sigma$, $\Sigma^*$, $\Delta$, $P_b$, $\mathfrak{so}(V)$, $\mathfrak{so}(W^\perp_{\mathbb R})$ and the pairs $(\Delta,P_b)$, $(\xi,\xi')$, $(\xi,\eta)$, $(\xi,P_b)$, $(\Delta,\xi)$ they are computed at $T=\infty$. Along $K_b$, $(K_b,P_{b'})$ and $(\eta,\eta')$ with $\Gamma(\eta,\eta')|_V\neq0$ the defect cycles are evaluated at a finite $T$; their $\varepsilon\to0$ limit is the same at every $T$ (Lemma \ref{lem:at-special}), and the coefficient $\tfrac3{8\pi^2}\hbar\log T$ of the composite $K_b$-insertions (Proposition c8-Tsafe) sits in the flows $W$ and in the trees of the observables, not in the obstruction.
\end{proposition}

\begin{proof}
Definition c8-ob and Proposition c8-Tindep give the $T$-independence of the singular parts and of the obstruction; Lemma \ref{lem:at-special} that of the finite parts of the cycles; Proposition c8-Tsafe the convergence at $T=\infty$ in the listed components; Theorem c8-K and Remark c8-K the $\log T$ terms, whose coefficient in a defect graph carries a special edge on a tree at a generic momentum and is $O(\varepsilon)$.
\end{proof}

\section{The arity-one equation: solvability and free directions}
\label{sec:at-eq1}

\begin{theorem}[The arity-one equation]
\label{thm:at-eq1}
Let $F_1^{(0)}$ solve the arity-zero equation, $\Lambda^{(0)}F_1^{(0)}=\Ob^{(0)}$, and let $x$ be a generator.
\begin{itemize}
\item (i) $\mathcal R^{(1)}(x)$ is $\Lambda^{(0)}$-closed, of degree $0$ and Weyl weight $4-\tfrac12\lambda_x$, of position degree at most $p_x$, with the covariance of Theorem \ref{thm:at-bound}.
\item (ii) The equation $\Lambda^{(0)}F_1^{(1)}(x)=\mathcal R^{(1)}(x)$ has a local solution exactly when the class $[\mathcal R^{(1)}(x)]$ vanishes in the cohomology $H^0(\Lambda^{(0)})$ of the local functionals with these quantum numbers.
\item (iii) Two solutions differ by a $\Lambda^{(0)}$-closed local functional of degree $-1$ with the same quantum numbers; modulo $\Lambda^{(0)}$-exact ones, by a class of $H^{-1}(\Lambda^{(0)})$.
\item (iv) Replacing $F_1^{(0)}$ by $F_1^{(0)}+G^{(0)}$ with $\Lambda^{(0)}G^{(0)}=0$ replaces $\mathcal R^{(1)}(x)$ by $\mathcal R^{(1)}(x)-\{G^{(0)},S(x)\}$; for $G^{(0)}=\Lambda^{(0)}\phi^{(0)}$ the change is $+\Lambda^{(0)}\{\phi^{(0)},S(x)\}$, absorbed by $F_1^{(1)}(x)\mapsto F_1^{(1)}(x)+\{\phi^{(0)},S(x)\}$ (the field redefinition generated by $\phi^{(0)}$, Proposition c8-laws (a)).
\item (v) The map $x\mapsto[\mathcal R^{(1)}(x)]$ satisfies, in $H^0(\Lambda^{(0)})$,
\[
[\mathcal R^{(1)}([x,y])]=(-1)^{|x|}[\{\mathcal R^{(1)}(x),S(y)\}]-[\{S(x),\mathcal R^{(1)}(y)\}],
\]
the Chevalley--Eilenberg one-cocycle condition of $\mathfrak{psl}(4|4)$ with values in $H^0(\Lambda^{(0)})$ for the action through the brackets with $S(y)$; the change (iv), $\mathcal R^{(1)}(x)\mapsto\mathcal R^{(1)}(x)-\{G^{(0)},S(x)\}$, satisfies the same identity with the left side included (it is the coboundary of $[G^{(0)}]$). The obstruction to solving through arity one, for some choice of the arity-zero solution, is its class in $H^1\big(\mathfrak{psl}(4|4);H^0(\Lambda^{(0)})\big)$.
\end{itemize}
\end{theorem}

\begin{proof}
(i) Lemma \ref{lem:at-sformulas} (iii). (ii), (iii): definition of the cohomology; locality of the solution is part of the statement of $H^0$ on local functionals. (iv) $\{\Lambda^{(0)}\phi^{(0)},S(x)\}=-\Lambda^{(0)}\{\phi^{(0)},S(x)\}$ by the derivation rule of the proof of Lemma \ref{lem:at-sformulas} and $\Lambda^{(0)}S(x)=0$. (v) The arity-two identity of Lemma \ref{lem:at-sformulas} (ii) for $G=\Ob-sF_1$, whose arity-zero component vanishes, reads $\Lambda^{(0)}G^{(2)}(x,y)=G^{(1)}([x,y])-(-1)^{|x|}\{G^{(1)}(x),S(y)\}+\{S(x),G^{(1)}(y)\}$; in $H^0$ the left side vanishes, which is the stated identity with $G^{(1)}=\mathcal R^{(1)}$. For $\mathcal R^{(1)}(x)=-\{G^{(0)},S(x)\}$ the identity follows from $(-1)^{|x|}\{S(x),S(y)\}=-S([x,y])-\Lambda^{(0)}S^{(2)}(x,y)$ and the derivation rule $\{G^{(0)},\{a,b\}\}=-\{\{G^{(0)},a\},b\}-(-1)^{|a|}\{a,\{G^{(0)},b\}\}$, the term $\{G^{(0)},\Lambda^{(0)}S^{(2)}\}=-\Lambda^{(0)}\{G^{(0)},S^{(2)}\}$ being exact.
\end{proof}

\subsection{Two legs}

On two legs the bracket with $I$ raises the number of legs, so the equation reads $QF_1^{(1)}(x)|_2=\mathcal R^{(1)}(x)|_2$ and the relevant cohomology is that of $Q$ on quadratic functionals.

\begin{proposition}[Two legs along $\Sigma$]
\label{prop:at-xi2}
Let $\xi\in S_+\otimes W$ (for $S_-\otimes W^*$ exchange the chiralities). Write $\Gamma_V:=\Gamma(\xi,\lambda)|_V$, $\Gamma_\perp:=\Gamma(\xi,\lambda)|_{W^\perp_{\mathbb R}}$ and $(\xi,\lambda^*)$ for the $\mathfrak g$-valued pairings, and
\[
g_1=\int(\Gamma_V,d^*d\AV),\quad g_2=\int(\Gamma_V,dd^*\AV),\quad g_3=\int(\Gamma_\perp,d^*d\Ap),\quad g_4=\int\big(d^*dc,(\xi,\lambda^*)\big),\quad T_0=\int\big(d^*\AV,(\xi,\lambda^*)\big).
\]
\begin{itemize}
\item (i) The quadratic local functionals of degree $-1$, Weyl weight $\tfrac92$, linear in $\xi$ and covariant, are spanned by $T_0,T_1,T_2,T_3,T_4$; those of degree $0$ by $g_1,g_2,g_3,g_4$ (legs $(\lambda_-,\AV)$, $(\lambda_+,\Ap)$, $(c,\lambda^*_+)$).
\item (ii) $QT_1=-QT_3=q_1g_1$, $QT_2=-QT_4=q_3g_3$, $QT_0=a\,g_2+b\,g_4$ with $q_1,q_3,a,b$ nonzero; in the normalisation of Definition \ref{def:at-names} with the left-derivation rule for $Q=\{S,-\}$, both engine routes give $q_1=q_3=-\tfrac12$, $a=b=-1$. The kernel of $Q$ on the degree $-1$ space is spanned by $S^V_2(\xi)=T_1+T_3$ and $S^\perp_2(\xi)=T_2+T_4$; no nonzero combination is $Q$-exact.
\item (iii) The $Q$-closed functionals of degree $0$ are spanned by $g_1$, $g_3$ and $a\,g_2+b\,g_4$, which is the image of $Q$: the cohomology $H^0(Q)$ of these quadratic functionals is zero.
\item (iv) Consequently the two-leg part of the arity-one equation along $\xi$ is solvable for every $F_1^{(0)}$ solving arity zero: with $\mathcal R^{(1)}(\xi)|_2=r_1g_1+r_3g_3+r_0(ag_2+bg_4)$,
\[
F_1^{(1)}(\xi)|_2=\frac{r_1}{q_1}T_1+\frac{r_3}{q_3}T_2+r_0T_0+u\big(\theta^\xi_VS^V_2(\xi)+\theta^\xi_\perp S^\perp_2(\xi)\big),
\]
with free $\theta^\xi_V,\theta^\xi_\perp$. The coefficient $r_0$ of the gauge-fixing structure $T_0$ is fed by the finite counterterm $\tfrac{27}{16}u\int(d^*\AV,d^*\AV)$ through $\{F_1^{(0)},T_1\}\propto g_2$, and by the parts $g_2$, $g_4$ of $\Ob^{(1)}(\xi)|_2$.
\item (v) Wess--Zumino on two legs: $Q\,\Ob^{(1)}(\xi)|_2=-\{\Ob^{(0)}|_2,T_1\}$ with $\Ob^{(0)}|_2=-\tfrac{27}8u\int(d^*dc,d^*\AV)$ (Theorem c4-obstruction): the components of $\Ob^{(1)}(\xi)|_2$ along $g_2$, $g_4$ outside the closed combination are fixed by the arity-zero two-leg obstruction (gate G-WZ1).
\item (vi) The $\beta$-directions of $F_1^{(0)}$ (Theorem c9-F3) are $Q$-exact on two legs, $Q$ of the field-counting functionals; through $\{Q N_f,S^{(1)}_2(\xi)\}=-Q\{N_f,S^{(1)}_2(\xi)\}$ they shift $F_1^{(1)}(\xi)|_2$ by $+\{N_f,S^{(1)}_2(\xi)\}$ (Theorem \ref{thm:at-eq1} (iv)), which multiplies each arrow $f\to f'$ of the classical transformation by the difference of the counting weights (Proposition qs-laws (a)).
\end{itemize}
\end{proposition}

\begin{proof}
(i) $\xi\otimes\lambda_-\in(S_+\otimes S_-)\otimes(W\otimes W^*)$ contains the singlet of $SU(4)$ only in $V\otimes\mathbf 1$, so every bilinear in $(\xi,\lambda_-)$ is a component of $\Gamma_V$; with $\AV$ and two derivatives, after integration by parts onto $\AV$, the contractions are $\Gamma_V\cdot d^*d\AV$ and $\Gamma_V\cdot dd^*\AV$ (the $\epsilon$-contraction with two symmetric derivatives vanishes). $\xi\otimes\lambda_+\in(\mathbf 1\oplus\Lambda^2_+)\otimes(\mathbf 6\oplus\mathbf{10})$ pairs with $\Ap\in\mathbf 6$ through $\Gamma_\perp$, and the $\Lambda^2_+$ part needs an antisymmetric pair of derivatives; $(\xi\otimes\lambda_+,\AV)$ and $(\xi\otimes\lambda_-,\Ap)$ have no $SU(4)$-singlet. $\xi\otimes\lambda^*_+$ contains the scalar $(\xi,\lambda^*)$, paired with $c$ and two derivatives; $\xi\otimes\lambda^*_-$ has no singlet. Degree $-1$: $(\lambda,A^*)$ without derivatives gives $T_1$, $T_2$; $(A,\lambda^*)$ with one derivative gives $(\xi,\lambda^*_+)\,d^*\AV$, the self-dual part of $d\AV$ against $\Lambda^2_+$ ($T_3$), and $\rho(d\Ap)\xi$ against $\lambda^*_-$ ($T_4$); the weight count of Theorem \ref{thm:at-bound} excludes the rest. (ii) $Q$ substitutes $A^*\mapsto-\tfrac12d^*dA$, $\lambda^*\mapsto\dslash\lambda$, $\AV\mapsto-dc$ (Proposition c1-Q): $T_1\mapsto\Gamma_V\cdot d^*d\AV$, $T_2\mapsto\Gamma_\perp\cdot d^*d\Ap$; in $T_0$ the substitution in $\AV$ gives $g_4$ and that in $\lambda^*$ gives $(\xi,\dslash\lambda)\,d^*\AV$ with $(\xi,\dslash\lambda)=\sum_\mu(e_\mu,\Gamma(\xi,\partial_\mu\lambda))=-d^*\Gamma_V$, hence $g_2$; $S^V_2$, $S^\perp_2$ are closed as the two sector parts of the quadratic part of $\{S,S(\xi)\}=0$ (Proposition qs-freexi). The images $g_1$, $g_3$, $ag_2+bg_4$ are independent, so the kernel is two-dimensional. A $Q$-exact functional of degree $-1$ would be $Q$ of a degree $-2$ quadratic functional of weight $\tfrac92$, which does not exist (Proposition qs-freexi (i)). (iii) $Qg_1=-\int(\Gamma_V,d^*d\,dc)=0$, $Qg_3=0$ ($\Ap$ and $\lambda$ are inert), and $Qg_2=-\int(\Gamma_V,dd^*dc)$, $Qg_4=\int(d^*dc,(\xi,\dslash\lambda))=-\int(dd^*dc,\Gamma_V)$ are proportional and nonzero; the closed space is three-dimensional and contains the three-dimensional image. (iv) Lemma \ref{lem:at-sformulas} (iii) makes $\mathcal R^{(1)}(\xi)|_2$ closed; (iii) and (ii). The bracket $\{\int(d^*\AV,d^*\AV),T_1\}$ contracts $\AV$ with $A^*|_V$ and equals $\pm2\int(d^*\Gamma_V,d^*\AV)=\pm2g_2$. (v) Lemma \ref{lem:at-sformulas} (ii) on two legs. (vi) The quadratic part of $s$ of $N_f$ is $QN_f$, and $\{QN_f,S_2\}=-Q\{N_f,S_2\}$ by the derivation rule and $QS_2=0$.
\end{proof}

\begin{proposition}[Two legs along the other generators]
\label{prop:at-two}
\leavevmode
\begin{itemize}
\item (i) $\Delta$. The $Q$-closed quadratic functionals of degree $-1$, weight $4$, position degree at most one and reflection-even are spanned modulo $Q$-exact ones by the sector parts $S(\Delta)_\perp,S(\Delta)_\lambda,S(\Delta)_V$, the exact ones by $Q\int(\iota_x\AV,c^*)$ and $Q\int(\rho(x)\lambda^*,\lambda^*)$ (Proposition c11-tc0). The two-leg equation is solved by the family of Definition c11-free once $\Ob^{(1)}(\Delta)|_2$ is known; the strict route computes $\Ob^{(1)}(\Delta)|_2$ by \eqref{eq:at-localformula}, regression target
\begin{equation}
\hbar\Ob^{(1)}(\Delta)\big|_2=u\Big[-\tfrac{11}3\int(dA,dA)+\tfrac{47}3\int(\lambda,\dslash\lambda)-\tfrac{14}3\int(dc,A^*)\Big]
\label{eq:at-reg2}
\end{equation}
(Theorem c9-quadratic (ii)), with coefficient $0$ for $\int(d^*\AV,d^*\AV)$ and for every functional with an explicit factor $x$; by Proposition \ref{prop:at-pos} the last statement is $\Ob^{(1)}(P_b)|_2=0$.
\item (ii) $P_b$. The quadratic $Q$-closed functionals of degree $-1$, weight $5$, linear in $b$ and translation invariant are spanned modulo exact ones by $S(P_b)_\perp,S(P_b)_\lambda,S(P_b)_V$ (the position parts of those of (i), Proposition \ref{prop:at-pos}); the exact ones by $Q\int(\iota_b\AV,c^*)$, $Q\int(\rho(b)\lambda^*,\lambda^*)$. Completeness at this weight is task S1-P.
\item (iii) $\mathfrak{so}(V)$, $\mathfrak{so}(W^\perp_{\mathbb R})$. The sector parts of $S(M_\omega)$, of $S(R_r)$ ($\perp$ and $\lambda$ only) are closed; along $R_r$, $F_1^{(1)}(R_r)=0$ solves the equation (Theorem \ref{thm:at-second} (iii)). Along $M_\omega$, in the leg sets $(\AV,\AV)_2$, $(c,A^*|_V)_1$ the closed functionals of degree $0$ are spanned by $Q\int\omega_{\mu\nu}(A_\mu,A^*_\nu)$, so $H^0(Q)=0$ there, and the spin term $\int\omega_{\mu\nu}(A_\mu,A^*_\nu)$ of $F_1^{(1)}(M_\omega)$ is fixed by $\Ob^{(1)}(M_\omega)|_2$. In the leg set $(\lambda_+,\lambda_-)_1$ the $\omega$-linear functionals of degree $0$ are spanned by $Q\int(\tfrac14\omega_{\mu\nu}\rho_\mu\rho_\nu\lambda,\lambda^*)$, the structure $\int(\lambda,\omega_{\mu\nu}\rho_{[\mu}\rho_\nu\rho_{\rho]}\partial_\rho\lambda)$ being a total derivative ($\Gamma$ symmetric), so $H^0(Q)=0$ there as well; with $S(M_\omega)_{\rm spin}:=\int\omega_{\mu\nu}(A_\mu,A^*_\nu)+\int(\tfrac14\omega_{\mu\nu}\rho_\mu\rho_\nu\lambda,\lambda^*)$ both engine routes give $\hbar\Ob^{(1)}(M_\omega)|_2=\tfrac13u\,Q\,S(M_\omega)_{\rm spin}$, uniformly in the two sectors, and $0$ on $(\Ap,\Ap)$ (Section \ref{sec:at-computed}).
\item (iv) $\eta\in S_+\otimes W^*$. The translation-invariant quadratic functionals of degree $-1$ and weight $\tfrac72$ are spanned by $\int(\rho(\AV)\eta,\lambda^*)$ and $\int s_0$; those of degree $0$ by $\int(\eta,\lambda_+)\,d^*\AV$, $\int(\eta,\rho(d\AV)\lambda_+)$, the structure $(\lambda_-,\Ap)_1$ and the structure $(c,\lambda^*_-)_1$. $Q$ is block triangular for the split into position parts (Proposition \ref{prop:at-pos}) and intrinsic parts; its kernel on the full degree $-1$ space is spanned by $S^V(\eta)$ and $S^\perp(\eta)$ (route B of Part Q, task TQ0 along $\Sigma^*$). The cohomology $H^0(Q)$ on the full degree-$0$ two-leg space at weight $\tfrac72$ with position degree at most one (task S1-$\eta$) is one-dimensional, a class in the $(\Ap,\lambda)$ sector (both routes); the computed residual $\mathcal R^{(1)}(\eta)|_2$ has no component along it, so the two-leg equation along $\eta$ is solvable, with the two free directions $S^V(\eta)$, $S^\perp(\eta)$ (Section \ref{sec:at-computed}).
\item (v) $K_b$. The sector parts $S(K_b)_d$ are closed; the intrinsic two-leg target has the leg sets $(c,A^*|_V)_0$, $(\AV,\AV)_1$, $(\Ap,\Ap)_1$, $(\lambda_+,\lambda_-)_0$ (Table \ref{tab:at-census1}); $H^{-1}(Q)$ and $H^0(Q)$ with position degree at most two (task S1-K) have dimensions $7$ and $4$ (both routes): $H^{-1}$ is spanned by the three $[S(K_b)_d]$ and four second-order free-field symmetries; the residual has no component on $H^0$ (Section \ref{sec:at-computed}). $\{\Ob^{(0)}|_2,S(K_b)\}\neq0$ because $\int(d^*dc,d^*\AV)$ is not invariant under the special conformal transformations, and $\{F_1^{(0)},S(K_b)\}\neq0$ for the same reason.
\end{itemize}
\end{proposition}

\begin{proof}
(i) Proposition c11-tc0 and Theorem c9-quadratic (iii); the regression values are stated in Theorem c9-quadratic (ii). (ii) The position part along $b$ commutes with $Q$ (Lemma c11-canonical) and maps the closed and exact directions of (i) onto those listed. (iii) $v.$ commutes with $Q$; the leg-set computation in the vector sector is: $Q\int\omega_{\mu\nu}(\partial_\mu A_\nu,d^*\AV)$ and $Q\int\omega_{\mu\nu}(\partial_\mu c,A^*_\nu)$ are proportional and nonzero, so the closed subspace of these two is one-dimensional and contains $Q\int\omega_{\mu\nu}(A_\mu,A^*_\nu)\ne0$. (iv) The census of Table \ref{tab:at-census1}: $\eta\otimes\lambda^*_-\in V\otimes(\mathbf1\oplus\mathbf{15})$ pairs with $\AV$, $\eta\otimes\lambda^*_+\in(\mathbf1\oplus\Lambda^2_+)\otimes(\mathbf6\oplus\overline{\mathbf{10}})$ with $\Ap$; the degree-$0$ structures likewise. $Q$ commutes with $-\{S(P_b),-\}$, which gives the triangular form. (v) Census; $\{\int(d^*dc,d^*\AV),S(K_b)\}$ is the variation of a functional that is not conformally invariant.
\end{proof}

\subsection{Three and four legs: the sector coefficients and the arity-zero couplings}

\begin{theorem}[The cubic part fixes the sector differences]
\label{thm:at-cubic}
Let $x$ be a generator and write the quadratic part of a solution as $F_1^{(1)}(x)|_2=F^{\rm part}_2(x)+u\sum_d\theta^x_dS(x)_d+(\text{exact})$, with the sectors $d$ of Definition \ref{def:at-names} (ii). The part with three legs of the arity-one equation is
\[
Q\,F_1^{(1)}(x)|_3+\{I_3,F_1^{(1)}(x)|_2\}=\mathcal R^{(1)}(x)|_3,
\]
$I_3$ the cubic part of $I$, and $\mathcal R^{(1)}(x)|_3=\Ob^{(1)}(x)|_3-\{F_1^{(0)}|_2,S_3(x)\}-\{F_1^{(0)}|_3,S_2(x)\}$.
\begin{itemize}
\item (i) Each $\{I_3,S(x)_d\}$ is $Q$-closed, and $\sum_d\{I_3,S(x)_d\}=\{I_3,S_2(x)\}=-QS_3(x)$; the equation depends on the $\theta^x_d$ through $\sum_{d\ne d_0}(\theta^x_d-\theta^x_{d_0})[\{I_3,S(x)_d\}]$ in the cohomology $H^0(Q)_3$ of the cubic functionals of degree $0$, weight $4-\tfrac12\lambda_x$ and position degree at most $p_x$.
\item (ii) The cubic equation is solvable exactly when $[\mathcal R^{(1)}(x)|_3-\{I_3,F^{\rm part}_2(x)\}]=u\sum_{d\ne d_0}(\theta^x_d-\theta^x_{d_0})[\{I_3,S(x)_d\}]$ in $H^0(Q)_3$.
\item (iii) If the classes $[\{I_3,S(x)_d\}]$, $d\ne d_0$, are linearly independent, the cubic equation fixes the differences $\theta^x_d-\theta^x_{d_0}$. They are independent exactly when no combination $\sum_dc_dS(x)_d$ other than the multiples of $S_2(x)$ extends to a $\Lambda^{(0)}$-closed functional through three legs, that is, when the sector parts are not symmetries of the cubic interacting theory.
\item (iv) Along $\Delta$ this is the test GC13 of the composite workflow: the rank of $[\{I_3,S(\Delta)_\perp\}]$, $[\{I_3,S(\Delta)_\lambda\}]$ modulo $Q$ of the cubic functionals of degree $-1$, weight $4$ and position degree at most one; its position part along $b$ is the same test along $P_b$ with $S(P_b)_d$, and $\{F_1^{(0)},S(\Delta)\}=0$ makes $\mathcal R^{(1)}(\Delta)=\Ob^{(1)}(\Delta)$. Along $\xi$ the single class $[\{I_3,S^V_2(\xi)\}]$ decides $\theta^\xi_V-\theta^\xi_\perp$; along $\eta$, $[\{I_3,S^V(\eta)\}]$; along $K_b$, $\mathfrak{so}(V)$, $\mathfrak{so}(W^\perp_{\mathbb R})$, the classes of the brackets with the sector parts.
\end{itemize}
\end{theorem}

\begin{proof}
On three legs, $\Lambda^{(0)}$ applied to a functional with two and three legs gives $Q$ of the cubic part plus the bracket of $I_3$ with the quadratic part; $\{F_1^{(0)},S(x)\}$ has on three legs the two stated terms. (i) $QI_3=0$ (cubic part of the classical master equation, Proposition c9-maps (i)) and $QS(x)_d=0$ give $Q\{I_3,S(x)_d\}=0$ by the derivation rule; the cubic part of $\{S,S(x)\}=0$ is $QS_3(x)+\{I_3,S_2(x)\}=0$. (ii) Definition of $H^0(Q)_3$; the exact quadratic directions contribute $\{I_3,Q\psi\}=\pm Q\{I_3,\psi\}$, exact. (iii) A relation $\sum_dc_d[\{I_3,S(x)_d\}]=0$ means $\{I_3,\sum_dc_dS(x)_d\}=-QY_3$ for a cubic $Y_3$, i.e. $\Lambda^{(0)}(\sum_dc_dS(x)_d+Y_3)=0$ through three legs; the combination $\sum_dS(x)_d$ always extends, by $Y_3=S_3(x)$. (iv) Lemma c9-SDelta (i): $\{F,S(\Delta)\}=-v.F=0$ for $F$ of weight $4$; Proposition \ref{prop:at-pos}.
\end{proof}

\begin{proposition}[The arity-one equations along $\Sigma$ constrain the arity-zero couplings]
\label{prop:at-couplings}
Let $\mathcal H_4:=H^0(\Lambda^{(0)})$ on colour-singlet local functionals of degree $0$ and weight $4$, invariant under $\mathrm{Spin}(4)\times\mathrm{Spin}(6)$ and the reflection: the classical couplings modulo field redefinitions, containing the class of $S$ (the coupling direction of Theorem c8-obstruction (iv)) and the classes $\kappa_1,\dots,\kappa_4$ of Theorem c9-F3 that extend to $\Lambda^{(0)}$-closed functionals. For $\xi\in\Sigma$ let $\delta_\xi:\mathcal H_4\to H^0(\Lambda^{(0)})_{9/2}$, $[G]\mapsto[\{G,S(\xi)\}]$. Then:
\begin{itemize}
\item (i) the arity-one equations along $\Sigma$ are solvable for some arity-zero solution exactly when $[\mathcal R^{(1)}(\xi)]$ lies in the image of $\delta_\xi$ for every $\xi$, the same $[G]$ serving all $\xi$;
\item (ii) the arity-zero freedom left after the arity-one equations along $\Sigma$ is $\bigcap_\xi\ker\delta_\xi$, the supersymmetric classes; it contains $[S]$, since $\{S,S(\xi)\}=0$;
\item (iii) along $\Delta$, $P_b$, $\mathfrak{so}(V)$, $\mathfrak{so}(W^\perp_{\mathbb R})$ and $K_b$ the map $[G]\mapsto[\{G,S(x)\}]$ vanishes on $\mathcal H_4$: the classes have weight $4$, are translation and rotation invariant, and the gauge-invariant functionals of weight $4$ of the fields $A$, $\lambda$ are conformally invariant; these equations impose no condition on the arity-zero couplings.
\end{itemize}
In particular the three- and four-leg parts of the arity-one equations along $\Sigma$, with the cubic finite counterterm of Theorem c9-F3, fix the coefficients $\kappa_k$ and the $\beta$-completions up to the coupling direction and field redefinitions.
\end{proposition}

\begin{proof}
Theorem \ref{thm:at-eq1} (iv), (v). (iii): $\{G,S(x)\}=\mp v.G$ for even $x$ (Lemma c9-SDelta (i), Definition c7-vdot), zero for $G$ of weight $4$ and invariant under translations and rotations; for $K_b$, the functionals $\int(F_A,F_A)$, $\int(d_A\iota_wA,d_A\iota_wA)$, $\int(\lambda,\dslash_A\lambda)$, the Yukawa and quartic terms are conformally invariant with the weights $1$, $\tfrac32$ (Proposition c2-coeff), and the antifield terms of weight $4$ are invariant because $S$ is.
\end{proof}

\begin{hypothesis}[Classical symmetries]
\label{hyp:at-H}
For every generator $x$ (resp. pair $(x,y)$), every $\Lambda^{(0)}$-closed colour-singlet local functional of degree $-1$, of Weyl weight $4-\tfrac12\lambda_x$ (resp. $4-\tfrac12(\lambda_x+\lambda_y)$), of position degree at most $p_x$ (resp. $p_x+p_y$), linear in the Killing data and equivariant as in Theorem \ref{thm:at-bound}, is a linear combination of the functionals $S(z)$, $z\in\mathfrak{psl}(4|4)$ with these quantum numbers, plus a $\Lambda^{(0)}$-exact functional.
\end{hypothesis}

\begin{remark}[Status of the hypothesis]
\label{rem:at-H}
Hypothesis \ref{hyp:at-H} states that the global symmetries of the interacting classical theory at these weights are the superconformal ones. On two legs it fails: $H^{-1}(Q)$ contains the sector parts $S(z)_d$ and, at weight $5$ with position degree one (the weight of the pair $(\Delta,P_b)$), the second-order symmetries of the free fields built from the conformal Killing tensors $b^{(\mu}x^{\nu)}$. The hypothesis is the statement that none of these extends through three legs except the full $S(z)$; Theorem \ref{thm:at-cubic} (iii) is its test at the weights of the arity-one components, and task S1-cubic of the engines computes it, at the weights of the arity-two pairs as well.
\end{remark}

\begin{theorem}[Solution space at arity one]
\label{thm:at-H}
Assume Hypothesis \ref{hyp:at-H} and that the arity-one equations are solvable. Then for every generator $x$ the solutions are $F_1^{(1)}(x)=F^{\rm part}(x)+t_x\,u\,S(x)+\Lambda^{(0)}(\text{local})$ with one number $t_x$, the uniform coefficient: all sector coefficients $\theta^x_d$ are equal to $t_x$ up to the fixed differences of Theorem \ref{thm:at-cubic} (iii) contained in $F^{\rm part}$. In a scheme equivariant for translations, rotations and $\mathfrak{so}(W^\perp_{\mathbb R})$, $t_x$ depends only on the type of $x$: $t_\Delta,t_P,t_M,t_R,t_{\xi_+},t_{\xi_-},t_{\eta_+},t_{\eta_-},t_K$.
\end{theorem}

\begin{proof}
Theorem \ref{thm:at-eq1} (iii) and Hypothesis \ref{hyp:at-H}; the only functional $S(z)$ with the weight, position degree and Killing data of $x$ that is equivariant is $S(x)$ itself. Equivariance under the compact groups makes $t_x$ constant on each irreducible block.
\end{proof}

\begin{remark}[The free directions and the equations that fix them]
\label{rem:at-freetable}
The differences of the sector coefficients are fixed by the cubic part of the arity-one equation (Theorem \ref{thm:at-cubic}) and, along $\Sigma$, $\Sigma^*$, $P_b$, $\Delta$, $\mathfrak{so}(V)$, $\mathfrak{so}(W^\perp_{\mathbb R})$, already by the two-leg parts of the arity-two equations (computed, Section \ref{sec:at-computed}); the uniform parts by the arity-two equations of Section \ref{sec:at-ar2}:
\[
\begin{array}{l|l|l}
x & \text{closed quadratic directions} & \text{uniform part fixed by}\\
\hline
\Delta & S(\Delta)_\perp,S(\Delta)_\lambda,S(\Delta)_V;\ Q\int(\iota_x\AV,c^*),\ Q\int(\rho(x)\lambda^*,\lambda^*) & (\Delta,P_b),\ (\Delta,\xi),\ (\xi,\eta)\\
P_b & S(P_b)_d;\ Q\int(\iota_b\AV,c^*),\ Q\int(\rho(b)\lambda^*,\lambda^*) & \text{up to the inner direction}\\
\mathfrak{so}(V) & S(M_\omega)_d & (K_b,P_{b'}),\ (\xi,\eta)\\
\mathfrak{so}(W^\perp_{\mathbb R}) & S(R_r)_\perp,\ S(R_r)_\lambda & (\xi,\eta)\\
\Sigma & S^V_2(\xi),\ S^\perp_2(\xi);\ \text{difference fixed by }(\xi,\xi'),\ (\xi_\pm,\xi'_\pm) & \text{up to inner and outer}\\
\Sigma^* & S^V(\eta),\ S^\perp(\eta);\ \text{difference fixed by }(\xi,\eta_\mp) & \text{up to inner and outer}\\
K_b & S(K_b)_d;\ \text{exact directions: task S1-K} & \text{up to the inner direction}
\end{array}
\]
\end{remark}

\section{Arity two}
\label{sec:at-ar2}

\subsection{Graphs and the local formula}

\begin{proposition}[Graph classes at arity two]
\label{prop:at-classes2}
Fix generators $x,y$ and $n\ge2$. At order $\hbar$, with $F_1=0$, the part with $n$ legs of the arity-two component along $(x,y)$ of the flow of $D$ consists of: (C2) irreducible cycles with exactly one special edge and either two vertices $S(x)$, $S(y)$ or one vertex $S^{(2)}(x,y)$ on the cycle, the other vertices from $I$; (Rd2) reducible graphs, whose subgraphs are cycles of type (C) along $x$ or $y$, of type (C2) with fewer legs, or of the arity-zero defect, joined by edges $P(\varepsilon,T)$ to trees of $I$, $S(x)$, $S(y)$, $S^{(2)}(x,y)$; (Br2) bridges; (Ct2) the counterterm trees of $D^{(2)}(x,y)$, of $D^{(1)}$ with a separate $S(y)$ or $S(x)$, and of $D^{(0)}$ with two separate vertices. The placements of each vertex are those of Lemma \ref{lem:at-place}; two vertices $S(x)$, $S(y)$ in the line placement are adjacent only through an $\AV\otimes\AV$ edge (both even: the chain $c$--$\AV$--$\AV$--$c$, a double composite line with both ends on ghost legs) or through a $\lambda\otimes\lambda$ edge (both odd, legs $\lambda$ of opposite chiralities), never when one is even and one odd (Lemma o-legs). As $\varepsilon\to0$, (Br2) with the special-edge parts of (Ct2) tends to $0$, (Rd2) to the trees of the components of $\Ob$ with lower arity or fewer legs, the $\log\varepsilon$ terms of (Ct2) cancel (Proposition \ref{prop:at-log}), (Ct2) has no pole (Proposition \ref{prop:at-poles}; $\{\int(A,A),S^{(2)}(x,y)\}=0$), and
\begin{equation}
\sum_{\mathcal C\in(C2),\,n}a_m(\mathcal C)=0\ (m\ge1),\qquad \hbar\,\Ob^{(2)}(x,y)\big|_{n\ \text{legs}}=\sum_{\mathcal C\in(C2),\,n\ \text{legs}}a_0(\mathcal C).
\label{eq:at-local2}
\end{equation}
The sign of a cycle with two vertices $S(x)$, $S(y)$ follows the arity-two component of the product of cochains (Theorem c8-L (iv)); the normalisation of Definition \ref{def:at-names} (i) is fixed by the gate G-sign (Section \ref{sec:at-engines}).
\end{proposition}

\begin{proof}
As Theorem \ref{thm:at-classes}, with Lemma \ref{lem:at-defect} for $D^{(2)}$: its order-$\hbar^0$ marked vertices are the special-edge parts of $(-1)^{|x|}\{S(x),S(y)\}_\varepsilon$ and $\{I,S^{(2)}(x,y)\}_\varepsilon$, and the marked vertices of lower arity enter with the remaining generators on separate vertices. $\{\int(A,A),S^{(2)}(x,y)\}=0$ since $\int(A,A)$ has no antifield and $S^{(2)}$ no leg $A^*$. The adjacency rules are Lemma o-legs.
\end{proof}

\begin{theorem}[The arity-two equation and its cocycle]
\label{thm:at-eq2}
Let $F_1^{(0)}$, $F_1^{(1)}$ solve the equations of arity zero and one.
\begin{itemize}
\item (i) $\mathcal R^{(2)}(x,y)$ is $\Lambda^{(0)}$-closed, of degree $-1$, Weyl weight $4-\tfrac12(\lambda_x+\lambda_y)$, position degree at most $p_x+p_y$; the arity-two equation along $(x,y)$ is solvable exactly when $[\mathcal R^{(2)}(x,y)]=0$ in $H^{-1}(\Lambda^{(0)})$.
\item (ii) Replacing $F_1^{(1)}$ by $F_1^{(1)}+\varphi$ with $\Lambda^{(0)}\varphi(x)=0$ for all $x$ replaces $\mathcal R^{(2)}(x,y)$ by
\[
\mathcal R^{(2)}(x,y)-\varphi([x,y])-(-1)^{|x|}\big(\{\varphi(x),S(y)\}+\{S(x),\varphi(y)\}\big),
\]
minus the Chevalley--Eilenberg coboundary of the one-cochain $[\varphi]$ with values in $H^{-1}(\Lambda^{(0)})$; $\Lambda^{(0)}$-exact $\varphi$ change $\mathcal R^{(2)}$ by exact terms.
\item (iii) The bracket $\{\hbar F_1^{(1)}(x),\hbar F_1^{(1)}(y)\}$ is of order $\hbar^2$ and enters at two loops; at one loop $F_1^{(1)}$ enters the arity-two equation linearly, through $F_1^{(1)}([x,y])$ and its brackets with $S(x)$, $S(y)$. $(x,y)\mapsto[\mathcal R^{(2)}(x,y)]$ is a two-cocycle of $\mathfrak{psl}(4|4)$ with values in $H^{-1}(\Lambda^{(0)})$ (arity-three component of $s(\Ob-sF_1)=0$, with $S^{(3)}=0$), and solvability through arity two for all pairs, for some choice of $F_1^{(1)}$, is the vanishing of its class in $H^2\big(\mathfrak{psl}(4|4);H^{-1}(\Lambda^{(0)})\big)$.
\item (iv) Under Hypothesis \ref{hyp:at-H}, $H^{-1}(\Lambda^{(0)})$ at the weights of the arity-two components is spanned by the classes $[S(z)]$, which form the adjoint module ($\{S(y),S(z)\}\equiv\pm S([y,z])$ modulo exact terms, Lemma c2-components); the uniform directions $\varphi(x)=t_x\,u\,S(x)$ change the class by
\begin{equation}
[\mathcal R^{(2)}(x,y)]\longmapsto[\mathcal R^{(2)}(x,y)]+u\big(t_x+t_y-t_{[x,y]}\big)[S([x,y])],
\label{eq:at-tshift}
\end{equation}
where $t_{[x,y]}S([x,y])$ means $\sum_z t_zS(z)$ over the components $z$ of $[x,y]$ in the blocks of Theorem \ref{thm:at-H}.
\end{itemize}
\end{theorem}

\begin{proof}
(i) Lemma \ref{lem:at-sformulas} (iii). (ii) Definition \ref{def:at-ob}; for exact $\varphi=\Lambda^{(0)}\psi$, $\{\Lambda^{(0)}\psi(x),S(y)\}=-\Lambda^{(0)}\{\psi(x),S(y)\}$ and $\{S(x),\Lambda^{(0)}\psi(y)\}=(-1)^{|x|}\Lambda^{(0)}\{S(x),\psi(y)\}$ by the derivation rule of the proof of Lemma \ref{lem:at-sformulas}, both exact. (iii) The arity-three component of $s$ on a degree-$1$ cochain $G$ with $G^{(0)}=G^{(1)}=0$ is $\Lambda^{(0)}G^{(3)}$ plus a signed Chevalley--Eilenberg differential of $G^{(2)}$ for the action through the brackets with $S(y)$, the signs forced as in Lemma \ref{lem:at-sformulas} (ii) (derivation rule of Definition c1-bracket) by $s^2=0$ on $G=sF^{(2)}$; in $H^{-1}$ the first term vanishes. The arity-three formula is not written out symbol by symbol here; the cocycle property is used only through Theorem \ref{thm:at-fix} (iii). (iv) $(-1)^{|x|}\{S(x),S(y)\}=-S([x,y])-\Lambda^{(0)}S^{(2)}(x,y)$ (Lemma c2-components), so the change in (ii) for $\varphi=tuS$ is $-u\,t_{[x,y]}S([x,y])+u(t_x+t_y)\big(S([x,y])+\Lambda^{(0)}S^{(2)}(x,y)\big)$.
\end{proof}

\begin{theorem}[What the arity-two equations fix]
\label{thm:at-fix}
Assume Hypothesis \ref{hyp:at-H}, the solvability of the arity-one equations, and write $[\mathcal R^{(2)}(x,y)]_{t=0}=u\,[S(\varrho(x,y))]$ with $\varrho(x,y)\in\mathfrak{psl}(4|4)$ of weight $\lambda_x+\lambda_y$, the class at $t_x=0$ for all $x$. The arity-two equations along all pairs are solvable exactly when
\begin{equation}
\varrho(x,y)=\big(t\cdot[x,y]\big)-[t\cdot x,y]-[x,t\cdot y]\qquad\text{for all }x,y,
\label{eq:at-rho}
\end{equation}
with $t\cdot z:=t_zz$ on the blocks, i.e. when the Chevalley--Eilenberg cocycle $\varrho$ with values in the adjoint module is the coboundary of the diagonal one-cochain $t\cdot$. Component by component:
\[
\begin{array}{l|l|l}
\text{pair} & \text{bracket} & \text{equation}\\
\hline
(\Delta,P_b) & -2P_b & \varrho=2t_\Delta P_b\\
(\Delta,\xi) & -\xi & \varrho=t_\Delta\,\xi\\
(\xi_+,\xi_-) & P_{\Gamma(\xi_+,\xi_-)|_V} & \varrho=(t_P-t_{\xi_+}-t_{\xi_-})P_{\Gamma}\\
(\xi_+,\eta_+) & \text{components }\Delta,\ M,\ R & \varrho_z=(t_z-t_{\xi_+}-t_{\eta_+})\,z,\ z\in\{\Delta,M,R\}\\
(\eta_+,\eta_-) & K_{\Gamma(\eta_+,\eta_-)|_V} & \varrho=(t_K-t_{\eta_+}-t_{\eta_-})K_{\Gamma}\\
(K_b,P_{b'}) & (b,b')\Delta+M_{b,b'} & \varrho_\Delta=(t_\Delta-t_K-t_P)(b,b')\Delta\\
 & & \varrho_M=(t_M-t_K-t_P)M_{b,b'}\\
(M_\omega,P_b) & [M_\omega,P_b] & \varrho=-t_M\,[M_\omega,P_b]\\
(M_\omega,\xi),\ (R_r,\xi) & [M_\omega,\xi],\ [R_r,\xi] & \varrho=-t_M\,[M_\omega,\xi],\ -t_R\,[R_r,\xi]\\
(\Delta,\eta),\ (\Delta,K_b) & \eta,\ 2K_b & \varrho=-t_\Delta\eta,\ -2t_\Delta K_b\\
(P_b,\eta),\ (K_b,\xi) & \rho(b)\eta,\ \rho(b)\xi & \varrho=(t_{\xi}-t_P-t_{\eta})\rho(b)\eta,\ (t_\eta-t_K-t_\xi)\rho(b)\xi\\
(\xi_+,\xi'_+),\ (\xi_+,\eta_-),\ (\xi,P_b) & 0 & \varrho=0\\
(\eta_+,\eta'_+),\ (P_b,P_{b'}) & 0 & \varrho=0
\end{array}
\]
(in the rows $(P_b,\eta)$, $(K_b,\xi)$ the chirality of the image is the opposite one, $t_\xi$, $t_\eta$ meaning the coefficients of the image block).
and likewise with the chiralities exchanged. Consequently:
\begin{itemize}
\item (i) $t_\Delta$, $t_M$ and $t_R$ are fixed; $t_\Delta$ is overdetermined by $(\Delta,P_b)$, $(\Delta,\xi_\pm)$ and the $\Delta$-components of $(\xi_\pm,\eta_\pm)$;
\item (ii) the solutions $(t_P,t_{\xi_\pm},t_{\eta_\pm},t_K)$ form an affine plane over the two diagonal derivations of $\mathfrak{psl}(4|4)$: the inner one, $t_z=c\lambda_z$, generated by the field redefinition $\phi^{(0)}=c\,u\,S(\Delta)$ (equivalently the finite part of the operators of Proposition c9-scheme), and the outer one, $t_z=aB_z$, the direction $G_B$ (Definition c8-B);
\item (iii) the remaining relations are consistency conditions: the agreement of the determinations of $t_\Delta$ (from $(\Delta,P_b)$, $(\Delta,\xi)$, $(\Delta,\eta)$, $(\Delta,K_b)$ and the $\Delta$-components of $(\xi,\eta)$), of $t_M$ (from $(M_\omega,P_b)$, $(M_\omega,\xi)$, $(K_b,P_{b'})$ and $(\xi,\eta)$) and of $t_R$ (from $(R_r,\xi)$ and $(\xi,\eta)$), and the linear relations among the values of $\varrho$ on $(\xi_+,\eta_+)$, $(\xi_-,\eta_-)$, $(\xi_+,\xi_-)$, $(\eta_+,\eta_-)$, $(K_b,P_{b'})$, $(P_b,\eta)$, $(K_b,\xi)$, consequences of the cocycle condition of Theorem \ref{thm:at-eq2} (iii);
\item (iii$'$) translation covariance: if $F_1$ is translation covariant (Proposition \ref{prop:at-pos} applied to the solution, as in Definition c11-free), then $t_P=t_\Delta$ and $t_{\eta_\pm}=t_{\xi_\mp}$ up to the outer direction; the inner direction is excluded, since $\phi^{(0)}=c\,u\,S(\Delta)$ is not translation invariant (it shifts $F_1^{(1)}(P_b)$ by $2c\,u\,S(P_b)$ and $F_1^{(1)}(\Delta)$ by $0$). In a translation-covariant scheme the solution set is the outer direction $G_B$ alone;
\item (iv) on the commuting pairs $\varrho=0$ is automatic: $\varrho(x,y)$ depends bilinearly and equivariantly on $(x,y)$, and $\mathfrak{psl}(4|4)$ contains no element of weight $\lambda_x+\lambda_y$ of the required type ($S_+\otimes S_+$ contains no vector for $(\xi_+,\xi'_+)$ and $(\eta_+,\eta'_+)$; $S_+\otimes S_-\otimes W\otimes W$ contains no element of $\mathfrak{so}(V)\oplus\mathfrak{so}(W^\perp_{\mathbb R})\oplus\mathbb C\Delta$ for $(\xi_+,\eta_-)$; no element has weight $-3$ or $-4$ for $(\xi,P_b)$, $(P_b,P_{b'})$). The computed classes on these pairs test Hypothesis \ref{hyp:at-H} (gate G-zero).
\item (v) on two legs the sector coefficients $\theta^x_d$ enter separately, the $[S(z)_d]$ being independent in $H^{-1}(Q)$; the two-leg parts of the commuting-pair equations then impose the sector uniformity: $(\xi_\pm,\xi'_\pm)$ gives $\theta^\xi_V=\theta^\xi_\perp$, $(\xi,P_b)$ gives $\theta^P_V=\theta^P_\perp=\theta^P_\lambda$ (computed, Section \ref{sec:at-computed}).
\end{itemize}
\end{theorem}

\begin{proof}
\eqref{eq:at-tshift}: the equation along $(x,y)$ is $u[S(\varrho(x,y))]+u(t_x+t_y-t_{[x,y]})[S([x,y])]=0$, which is \eqref{eq:at-rho} componentwise. The brackets are those of Proposition c2-special, Remark c2-dilatation and Remark c7-KDelta; $[\xi_+,\eta_-]=0$ because $S_+\otimes S_-\otimes W\otimes W$ contains no element of $\mathfrak{so}(V)\oplus\mathfrak{so}(W^\perp_{\mathbb R})\oplus\mathbb C\Delta$. Solving: $(\Delta,P_b)$ gives $t_\Delta$; $(\xi_\pm,\eta_\pm)$ give $t_M-t_\Delta$, $t_R-t_\Delta$, $t_{\xi_\pm}+t_{\eta_\pm}$; $(\xi_+,\xi_-)$ gives $t_P-t_{\xi_+}-t_{\xi_-}$; $(\eta_+,\eta_-)$ gives $t_K-t_{\eta_+}-t_{\eta_-}$; $(K_b,P_{b'})$ gives $t_K+t_P$. On the six unknowns $t_P,t_{\xi_\pm},t_{\eta_\pm},t_K$ these are five equations of rank four, the sum of the equations from $(\xi_+,\eta_+)$, $(\xi_-,\eta_-)$, $(\xi_+,\xi_-)$ and $(\eta_+,\eta_-)$ being the equation from $(K_b,P_{b'})$ up to the values of $\varrho$: the homogeneous solutions are $t_z=c\lambda_z$ and $t_z=aB_z$, whose diagonal maps are derivations ($\lambda_{[x,y]}=\lambda_x+\lambda_y$; $B$ is a derivation). The inner direction: $(s\phi)^{(1)}(x)=\{\phi^{(0)},S(x)\}=c\,u\{S(\Delta),S(x)\}\equiv c\,\lambda_x\,u\,S(x)$ modulo exact terms; the outer: Proposition c8-laws (c).
\end{proof}

\subsection{The pair $(\Delta,P_b)$ and the coefficient $\theta$}

\begin{theorem}[The sector coefficients along $\Delta$ from the arity-two equation]
\label{thm:at-theta}
Let $\hbar F_1^{(1)}(\Delta)$ and $\hbar F_1^{(1)}(P_b)$ solve the arity-one equations with the quadratic parts of Definition c11-free (the particular part along $P_b$ vanishing on two legs, $\Ob^{(1)}(P_b)|_2=0$). On two legs the arity-two equation along $(\Delta,P_b)$ reads
\begin{equation}
Q\,F_1^{(2)}(\Delta,P_b)\big|_2+2u\sum_{d\in\{\perp,\lambda,V\}}\theta^\Delta_d\,S(P_b)_d=\Ob^{(2)}(\Delta,P_b)\big|_2+Q(\text{local}),
\label{eq:at-theta}
\end{equation}
independently of $\theta^{P_b}_d$. The classes $[S(P_b)_d]$ are independent in $H^{-1}(Q)$ on two legs. Hence, writing $o_d$ for the coefficient of $[S(P_b)_d]$ in $[\Ob^{(2)}(\Delta,P_b)|_2]$, the equation is solvable on two legs exactly when $[\Ob^{(2)}(\Delta,P_b)|_2]$ lies in the span of the $[S(P_b)_d]$, and then
\[
\theta^\Delta_d=\frac{o_d}{2u}\qquad(d=\perp,\lambda,V).
\]
Under Hypothesis \ref{hyp:at-H}, $o_d-2u(\theta^\Delta_d-\theta^\Delta_V)$ is independent of $d$, with the differences of Theorem \ref{thm:at-cubic}, and the components of $[\Ob^{(2)}(\Delta,P_b)|_2]$ along the second-order free symmetries of Remark \ref{rem:at-H} vanish.
\end{theorem}

\begin{proof}
In the third equation of \eqref{eq:at-eqs}, $[\Delta,P_b]=-2P_b$ and $S^{(2)}(\Delta,P_b)=0$. The sector parts of $S(\Delta)$ and of $S(P_b)$ in different sectors have no common legs, and within one sector $\{S(\Delta)_d,S(P_b)_d\}=v.S(P_b)_d=2S(P_b)_d$, the weight of $S(P_b)_d$ being $5$ (Lemma c9-SDelta (i)). The $\theta$-terms are therefore $2u\theta^\Delta_dS(P_b)_d$ from $\{F_1^{(1)}(\Delta),S(P_b)\}$, $2u\theta^{P_b}_dS(P_b)_d$ from $\{S(\Delta),F_1^{(1)}(P_b)\}$ and $-2u\theta^{P_b}_dS(P_b)_d$ from $F_1^{(1)}(-2P_b)$. The remaining quadratic parts of $F_1^{(1)}(\Delta)$ are combinations of field-counting functionals, which commute with $S(P_b)$, and $Q$-exact terms, whose brackets with $S(P_b)$ and $S(\Delta)$ are $Q$-exact. Independence: $Q$ preserves the three sectors $(\Ap,A^*|_\perp)$, $(\lambda,\lambda^*)$ and $(\AV,A^*|_V)\oplus(c,c^*)$ of the quadratic functionals of degree $-1$. In the first there is no functional of degree $-2$: $\int f(x)(A^*|_\perp,A^*|_\perp)=0$ identically, $A^*$ being odd and the fibre pairing symmetric, so $[S(P_b)_\perp]\ne0$ ($S(P_b)_\perp\ne0$). In the second, the image is spanned by $Q\int(\rho(b)\lambda^*,\lambda^*)=2\int(\rho(b)\dslash\lambda,\lambda^*)=2\int(\partial_b\lambda,\lambda^*)+2\int(\rho(b\wedge\partial)\lambda,\lambda^*)$ (and position-dependent images with nonzero position part), which is not a multiple of $S(P_b)_\lambda$. In the third, the translation-invariant image is spanned by $Q\int(\iota_b\AV,c^*)$, which differs from $S(P_b)_V$ by the nonzero $\int(\iota_bd\AV,A^*|_V)$ up to sign. Hence each $[S(P_b)_d]$ is nonzero and they are independent. The last statement: $\mathcal R^{(2)}(\Delta,P_b)$ is $\Lambda^{(0)}$-closed (Theorem \ref{thm:at-eq2} (i)); by Hypothesis \ref{hyp:at-H} its class is a multiple of $[S(P_b)]$, whose two-leg part is $\sum_dS(P_b)_d$.
\end{proof}

\begin{remark}[Cross-checks for $\theta$]
\label{rem:at-thetacheck}
Theorem c9-kappa derives $\theta=-\tfrac13$ from the arity-two component of $L\circ L=0$ on the scalar bilinears, under the hypothesis that the quantisation and the local representatives exist through arity two; Corollary c11-theta does the same family by family. In the strict route this value is the prediction $o_\perp=-\tfrac23u$ for the computed $\Ob^{(2)}(\Delta,P_b)$, a cross-check (X-theta). The same $t_\Delta$ must result from $(\Delta,\xi)$ and from the $\Delta$-components of $(\xi,\eta)$ (gate G-tDelta).
\end{remark}

\subsection{The other pairs}

\begin{proposition}[The pairs along $\Sigma$ and $\Sigma^*$]
\label{prop:at-pairs}
\leavevmode
\begin{itemize}
\item (i) $(\xi,\xi')$, $\xi\in S_+\otimes W$, $\xi'\in S_-\otimes W^*$: $S^{(2)}(\xi,\xi')=\int\big(\tfrac14(\Gamma(\xi,\xi'),\Gamma(\lambda^*,\lambda^*))-\tfrac12(\xi,\lambda^*)(\xi',\lambda^*)-(\iota_{\Gamma(\xi,\xi')}A,c^*)\big)$; the two-leg part of the equation is $QF_1^{(2)}(\xi,\xi')|_2+F_1^{(1)}(P_\Gamma)|_2-\big(\{F_1^{(1)}(\xi),S(\xi')\}+\{S(\xi),F_1^{(1)}(\xi')\}\big)|_2+\{F_1^{(0)},S^{(2)}(\xi,\xi')\}|_2=\Ob^{(2)}(\xi,\xi')|_2$. The sector coefficients enter through $u\theta^{P}_dS(P_\Gamma)_d$ and through the brackets $\{S^d_2(\xi),S_2(\xi')\}$ and $\{S_2(\xi),S^d_2(\xi')\}$, whose sum over $d$ is $-S(P_\Gamma)|_2-QS^{(2)}(\xi,\xi')$ and whose individual classes in $H^{-1}(Q)$ on two legs are combinations of the $[S(P_\Gamma)_{d'}]$ given by the Fierz identities of the arrows $\lambda\to\AV\to\lambda$ and $\lambda\to\Ap\to\lambda$ (task S2-$\xi\xi$). The equation is the one-loop part of $[Q,Q]=P$.
\item (ii) $(\xi_+,\xi'_+)$: the bracket vanishes, $S^{(2)}(\xi_+,\xi'_+)$ is given by (i) with $\Gamma(\xi_+,\xi'_+)|_V=0$ and $\Gamma|_\perp\ne0$; the class of $\mathcal R^{(2)}$ must vanish; for uniform coefficients Hypothesis \ref{hyp:at-H} makes this automatic (Theorem \ref{thm:at-fix} (iv)), and on two legs it is the condition $\theta^\xi_V=\theta^\xi_\perp$ (Theorem \ref{thm:at-fix} (v)); the rank-one case $\xi=\xi'$ enters the twisted obstruction (Remark \ref{rem:at-twist}).
\item (iii) $(\xi,\eta)$: $[\xi,\eta]$ is the vector field $x\mapsto\Gamma(\xi,\rho(x)\eta)|_V$ plus the rotation of $W^\perp_{\mathbb R}$ of (c2-oddodd), with components along $\Delta$, $\mathfrak{so}(V)$ and $\mathfrak{so}(W^\perp_{\mathbb R})$ for $(\xi_\pm,\eta_\pm)$ and zero for $(\xi_\pm,\eta_\mp)$; $S^{(2)}(\xi,\eta)$ is (c2-S2) at the Killing spinor $\rho(x)\eta$. The equation is the one-loop part of $\{Q,S\}=D+M+R$; it determines $t_\Delta$, $t_M$, $t_R$ relative to $t_{\xi}+t_\eta$.
\item (iv) $(\xi,P_b)$: $[\xi,P_b]=0$, $S^{(2)}(\xi,P_b)=0$; the class of $\mathcal R^{(2)}(\xi,P_b)$ must vanish, and under Hypothesis \ref{hyp:at-H} it does, $H^{-1}(\Lambda^{(0)})$ having no class at weight $\tfrac{11}2$; the computed class tests the hypothesis. The cycles with $S(P_b)$ at the special edge or on a leg reproduce the translation of the arity-one cycles along $\xi$; the others contain the composite line along $P_b$ (gate G-xiP).
\item (v) $(\eta,\eta')$: $[\eta_+,\eta_-]=K_{\Gamma(\eta_+,\eta_-)|_V}$; evaluated at finite $T$ (Definition c8-K); the one-loop part of $\{S,S\}=K$; Table \ref{tab:at-census2} has no sector at position degree $0$, the weight $3$ requiring the position for the legs of degree $-1$.
\item (vi) $(\Delta,\xi)$: $[\Delta,\xi]=-\xi$; the one-loop part of $[D,Q]=-Q$; it fixes $t_\Delta$ again (Theorem \ref{thm:at-fix}).
\item (vii) $(K_b,P_{b'})$: at finite $T$; the one-loop part of $[K,P]=D+M$; it fixes $t_K+t_P$ and $t_M-t_\Delta$.
\end{itemize}
\end{proposition}

\begin{proof}
Theorem \ref{thm:at-fix}, Definition c2-S2 and Proposition c2-special. (iv): $S(P_b)$ commutes with every edge, vertex and counterterm; the cycles of (C2) with $S(P_b)$ at the special edge or on a leg reproduce the translation of the arity-one cycles along $\xi$, and the remaining ones carry the composite line along $P_b$.
\end{proof}

\begin{theorem}[One-loop algebra relations]
\label{thm:at-algebra}
Assume Hypothesis \ref{hyp:at-H} and the arity-one equations. For each pair $(x,y)$ the following are equivalent:
\begin{itemize}
\item (a) the arity-two equation along $(x,y)$ is solvable;
\item (b) $[\mathcal R^{(2)}(x,y)]=0$, i.e. \eqref{eq:at-rho} holds at $(x,y)$;
\item (c) with the one-loop corrected generators $S_\hbar(z):=S(z)+\hbar\,t_z\,u\,S(z)$, $z\in\{x,y,[x,y]\}$, the order-$\hbar$ part of $(-1)^{|x|}\{S_\hbar(x),S_\hbar(y)\}+S_\hbar([x,y])$ plus $u\,S(\varrho(x,y))$ is $\Lambda^{(0)}$-exact: the one-loop corrections of the generators close under the bracket of $\mathfrak{psl}(4|4)$ once the arity-two anomaly class $u[S(\varrho(x,y))]$ is included, and the equation is solvable exactly when they compensate it.
\end{itemize}
The relations are: $(\xi,\xi')$ the one-loop correction to $[Q,Q]=P$; $(\xi,\eta)$ to $\{Q,S\}=D+M+R$; $(K_b,P_{b'})$ to $[K,P]=D+M$; $(\Delta,\xi)$ to $[D,Q]=-Q$; $(\Delta,P_b)$ to $[D,P]=-2P$; $(\eta,\eta')$ to $\{S,S\}=K$. On local operators, the map $G\mapsto-\{G,-\}$ turns the action-level identity into the one-body part of the arity-two component of $L\circ L=0$ (Corollary c8-homotopy (ii)) at order $\hbar$: the commutator of the one-body parts of $\hbar L_1^{(1)}(x)$ and $\{S(y),-\}$, plus $\hbar L_1^{(1)}([x,y])$, equals a graded commutator with $\Lambda^{(0)}$ plus $-\{\hbar F_1^{(2)}(x,y),-\}$-type terms.
\end{theorem}

\begin{proof}
(a)$\Leftrightarrow$(b): Theorem \ref{thm:at-eq2} (i). (b)$\Leftrightarrow$(c): with $(-1)^{|x|}\{S(x),S(y)\}=-S([x,y])-\Lambda^{(0)}S^{(2)}(x,y)$, the order-$\hbar$ part of $(-1)^{|x|}\{S_\hbar(x),S_\hbar(y)\}+S_\hbar([x,y])$ for $S_\hbar(z)=S(z)+\hbar t_zuS(z)$ is $u(t_x+t_y-t_{[x,y]})S([x,y])$ modulo exact terms, and \eqref{eq:at-tshift}. The operator statement: Proposition c8-laws (c) and Corollary c8-homotopy (ii), the one-body terms being $-\{\hbar F_1^{(1)}(x),-\}$ (Theorem c11-classes (iii)).
\end{proof}

\begin{remark}[The twist as a cross-check]
\label{rem:at-twist}
For a rank-one $\xi=u_0\otimes z\in S_+\otimes W$, $[\xi,\xi]=0$, and the evaluation $r_\xi$ of Lemma c8-ev maps $\Ob$ to the twisted obstruction $\Ob^{(0)}+\Ob^{(1)}(\xi)-\tfrac12\Ob^{(2)}(\xi,\xi)$ and $F_1$ to $F_1^{(0)}+F_1^{(1)}(\xi)-\tfrac12F_1^{(2)}(\xi,\xi)$, which solves the one-loop twisted master equation of $S^\xi$ exactly when the equations of arities $0,1,2$ along $\xi$ hold (no components of arity $\ge3$, Lemma o-actiontrunc). The comparison with the holomorphic twist of [BK, (3.64)--(3.70)] is a cross-check only (X-twist).
\end{remark}

\section{Closed forms and methods}
\label{sec:at-closed}

\begin{proposition}[Scaled Schwinger form of a defect cycle]
\label{prop:at-schwinger}
For a cycle of Lemma \ref{lem:at-special} with $n_t$ edges, external momenta $p_1,\dots,p_n$ and superficial degree $D$, substitute $t_e=\varepsilon\tau_e$, $\tau_\sigma\in[0,1]$ for the special edge and $\tau_e\in[1,\infty)$ for the others (both segments of a composite line included), and integrate the loop momentum against the Gaussian. With $\tau=\sum_e\tau_e$, $\alpha_e=\tau_e/\tau$, the coefficient of $\varepsilon^{j-D/2}$ of the Taylor expansion in the external momenta is the order-$j$ Taylor coefficient of the Gaussian-integrated integrand integrated over the scaled region; for $j<D$ odd orders vanish by Lorentz parity and even orders give the poles $a_{(D-j)/2}$, the order $j=D$ gives $a_0$, and the orders $j>D$ give $O(\varepsilon^{1/2}\log\varepsilon)$ or smaller. Each $a_0$ is a finite sum of master integrals
\[
J(n_t;\beta;N)=\int_0^1d\tau_\sigma\int_{[1,\infty)^{n_t-1}}\prod_{e\ne\sigma}d\tau_e\;\tau^\beta\,\tau^{-N},
\]
($\tau^\beta$ a monomial in the $\tau_e$), convergent when $N>|\beta|+n_t-1$, with values in $\mathbb Q+\mathbb Q\log2+\mathbb Q\log3$ for $n_t\le3$; on the simplex this is the formula of Lemma \ref{lem:at-special}.
\end{proposition}

\begin{proof}
Theorem c9-sing (i) and its proof, which applies to every cycle with one special edge; the master integrals are those of the Part-C engine of Step 2 (route A), evaluated exactly and confirmed at $50$ digits.
\end{proof}

\begin{proposition}[Small closed forms]
\label{prop:at-small}
With $F_m$ the $m$-th simplex integrand:
\begin{itemize}
\item (i) $n_t=2$: $16\pi^2a_0=\int_0^{1/2}F_0(\alpha)\log\frac{1-\alpha}\alpha\,d\alpha$, $\alpha=\alpha_\sigma$; for constant $F_0$ this is $F_0\log2$;
\item (ii) constant $F_0$ ($D=0$), any $n_t$: $16\pi^2a_0=F_0\,\frac{1}{(n_t-1)!}\log\frac{n_t}{n_t-1}$; constant $F_1$ ($D=2$): $16\pi^2a_1=\frac{F_1}{n_t^2(n_t-1)\,(n_t-1)!}$;
\item (iii) one vertex of $I$ and a composite line ($k=1$, Proposition c11-k1): $16\pi^2a_0=\int_0^{1/3}F_0(\beta)\big[(1-\beta)\log\frac{1-\beta}{2\beta}-(1-3\beta)\big]d\beta$, in $\mathbb Q+\mathbb Q\log2+\mathbb Q\log3$;
\item (iv) triangles of cubic vertices with the special edge: values in $\mathbb Q+\mathbb Q\log\tfrac32$ (the table of Section 9.1).
\end{itemize}
\end{proposition}

\begin{proof}
(i) For $n_t=2$ the region where $\alpha_\sigma$ is the smaller coordinate is $\alpha\le\tfrac12$ and $\min_{f\ne\sigma}\alpha_f=1-\alpha$; $\int_0^{1/2}\log\frac{1-\alpha}\alpha d\alpha=\log2$. (ii) By symmetry the region where $\alpha_\sigma$ is smallest is $1/n_t$ of the simplex, of volume $1/(n_t-1)!$; the logarithmic average over it is $\log\frac{n_t}{n_t-1}$ (Part Q, closed forms B2). (iii) Proposition c11-k1. (iv) Theorem c9-sing and the table of Section 9.1.
\end{proof}

\begin{proposition}[Position factors and Killing data]
\label{prop:at-inject}
For $x\in\{\Delta,\mathfrak{so}(V),\eta,K_b\}$, a vertex $S(x)$ at the point $y$ carries a polynomial of degree $p_x$ in $y$. With a momentum $r$ injected at that vertex and absorbed by the plane wave of the functional, $y^\mu\mapsto-i\partial/\partial r_\mu$ and $y^\mu y^\nu\mapsto-\partial^2/\partial r_\mu\partial r_\nu$ at $r=0$, acting on the integrand and on the momentum delta function. The action on the delta function gives the position parts of Proposition \ref{prop:at-pos}; the action on the integrand gives the intrinsic part for the chosen canonical densities (Definition c11-canonical: all momenta but one held fixed). The $\varepsilon^0$ coefficient is polynomial in $r$, so exact Lagrange interpolation on integer $r$ reproduces the derivatives.
\end{proposition}

\begin{proof}
$\int\mathrm{dvol}\,y^\mu e^{i(r,y)}(\cdots)=-i\partial_{r_\mu}\int\mathrm{dvol}\,e^{i(r,y)}(\cdots)$; Lemma \ref{lem:at-special} gives polynomiality.
\end{proof}

\begin{proposition}[Second route: renormalised amplitudes]
\label{prop:at-routeB}
For each generator and each sector, the renormalised arity-one amplitude $W^{(1)}_{\rm ren}(x)$ of Theorem \ref{thm:at-second} is a sum of one-loop integrals with propagators $e^{-\varepsilon|q|^2}|q|^{-2}$ (and $e^{-2\varepsilon|q|^2}|q|^{-4}$ times the numerator of $(Q^{GF})^2$ on composite lines), minus the trees of $\ICT$. Their finite parts at $\varepsilon\to0$ are Feynman-parameter integrals of rational functions times logarithms of the external invariants; $Q$ acts on the plane-wave value by replacing one leg by its image under the substitutions of Proposition c1-Q, and the nonlocal terms ($\log$ of momenta, ratios $(p,a)(p,b)|p|^{-2}$) cancel in $Q\,W^{(1)}_{\rm ren}(x)$ up to the trees of the components with fewer legs. A third evaluation of each special-edge cycle at finite $\varepsilon$ (numerical, $\varepsilon=2^{-k}$, Richardson extrapolation of $\varepsilon\,\mathcal C$ and $\mathcal C-\sum_m\varepsilon^{-m}a_m$) checks the absence of $\log\varepsilon$ and the value of $a_0$.
\end{proposition}

\begin{proof}
Theorem \ref{thm:at-second} (i); the Feynman-parameter form follows from the Schwinger form by the substitution $t_e=\tau\alpha_e$ and integration over $\tau$; Lemma \ref{lem:at-special} gives the expansion used in the extrapolation.
\end{proof}

\begin{remark}[Theorems of the construction used as cross-checks]
\label{rem:at-cross}
The following identities are consequences of the construction and are used only as cross-checks of the two routes: Theorem \ref{thm:at-second} (v), $\Ob^{(1)}(\Delta)=4\lc^{(0)}+(\text{composite part})$, whose two-leg form is $u[-4,16,-4]+u[\tfrac13,-\tfrac13,-\tfrac23]$ on $\int(dA,dA)$, $\int(\lambda,\dslash\lambda)$, $\int(dc,A^*)$ (X-scal); the value $\theta=-\tfrac13$ of Theorem c9-kappa (X-theta); the twisted obstruction of Remark \ref{rem:at-twist} against the twist literature (X-twist); and $K=P^\dagger$ (Corollary c8-gram) for the special conformal data (X-KP). None of them is used to produce a number.
\end{remark}

\section{Computed two-leg values, solutions and relations}
\label{sec:at-computed}

Values of the two engine routes of Section \ref{sec:at-engines} (step2/strict/action/engine\_a, engine\_b, complete), equal entry by entry in $\mathbb Q(i)+\mathbb Q(i)\log2+\mathbb Q(i)\log3$, with the signs of Lemma \ref{lem:at-sformulas}. Units $u$; $L:=\log\tfrac32$; $T_0,\dots,T_4$, $g_1,\dots,g_4$, $s_0$, $S(\Delta)_d$, $S(P_b)_d$, $N$ as in Proposition \ref{prop:at-xi2} and Remark \ref{rem:at-twocomp}; $u_0:=\int(\rho(\AV)\eta,\lambda^*)$; $S(M_\omega)_{\rm spin}$ as in Proposition \ref{prop:at-two} (iii); $F^{\rm part}(\xi'):=u(L-\tfrac{79}{12})(T_1+T_2)(\xi')$ for $\xi'\in\Sigma$.

\begin{proposition}[Arity-one obstructions on two legs]
\label{prop:at-c-ob1}
\[
\begin{aligned}
&\hbar\Ob^{(0)}|_2=-\tfrac{27}8u\int(d^*dc,d^*\AV),\qquad \hbar\Ob^{(1)}(\Delta)|_2=\eqref{eq:at-reg2},\qquad \hbar\Ob^{(1)}(P_b)|_2=0,\qquad \hbar\Ob^{(1)}(R_r)|_2=0,\\
&\hbar\Ob^{(1)}(M_\omega)|_2=\tfrac13u\,Q\,S(M_\omega)_{\rm spin},\qquad \hbar\Ob^{(1)}(\xi)|_2=u\Big[\big(\tfrac{79}{24}-\tfrac12L\big)(g_1+g_3)+\tfrac{27}8g_2+0\cdot g_4\Big],\\
&\hbar\Ob^{(1)}(\eta)|_2=\sum_\mu x^\mu\,\hbar\Ob^{(1)}(\rho(e_\mu)\eta)|_2+u\Big[-\tfrac52\int(\eta,\lambda)\,d^*\AV-\tfrac52\int(\eta,\rho(d\AV)\lambda)+\big(-\tfrac32+3L\big)\int(\eta,\rho(d\Ap)\lambda)+\tfrac52\int(\rho(dc)\eta,\lambda^*)\Big]+Q(\cdot),
\end{aligned}
\]
the intrinsic part of the last line in the canonical density (Definition c11-canonical); along $K_b$ the obstruction is in the span of the canonical data with position degree $\le2$, with pole sum $\tfrac94\{\int(A,A),S(K_b)\}$. $\lc^{(1)}(\xi)|_2=\{uN,S_2(\xi)\}=2u(T_1+T_2-T_3-T_4)$, so $l_\xi=0$ in Remark \ref{rem:at-logexact} on two legs. The Wess--Zumino identities of Lemma \ref{lem:at-sformulas} (ii) hold on every entry; the printed form before the revision (opposite sign) fails on every entry where it was tested.
\end{proposition}

\begin{proposition}[Two-leg solutions with all free directions]
\label{prop:at-c-sol}
$F_1^{(0)}|_2=\tfrac{27}{16}u\int(d^*\AV,d^*\AV)$ plus the $s$-closed arity-zero directions, and
\[
\begin{aligned}
\hbar F_1^{(1)}(\Delta)|_2&=u\Big[\tfrac{22}3\int(A,A^*)-\tfrac{47}3\int(\lambda^*,\lambda)-\tfrac83\int(c^*,c)\Big]+u\sum_d\theta^\Delta_dS(\Delta)_d\\
&\quad+\kappa_V\,uQ\!\int(\iota_x\AV,c^*)+\kappa_\lambda\,uQ\!\int(\rho(x)\lambda^*,\lambda^*),\\
\hbar F_1^{(1)}(P_b)|_2&=u\sum_d\theta^P_dS(P_b)_d+(\text{the }b\text{-parts of the }\kappa\text{ directions}),\\
\hbar F_1^{(1)}(\xi)|_2&=F^{\rm part}(\xi)+u\,\theta^\xi_VS^V_2(\xi)+u\,\theta^\xi_\perp S^\perp_2(\xi)+\{N_f,S_2(\xi)\}\text{-shifts},\\
\hbar F_1^{(1)}(\eta)|_2&=\sum_\mu x^\mu F^{\rm part}(\rho(e_\mu)\eta)-\big(3L-\tfrac32\big)u\,s_0-\tfrac52u\,u_0+u\,\theta^{\eta}_VS^V(\eta)+u\,\theta^{\eta}_\perp S^\perp(\eta),\\
\hbar F_1^{(1)}(M_\omega)|_2&=\tfrac13u\,S(M_\omega)_{\rm spin}+u\sum_d\theta^M_dS(M_\omega)_d,\qquad \hbar F_1^{(1)}(R_r)|_2=u\sum_d\theta^R_dS(R_r)_d,\\
\hbar F_1^{(1)}(K_b)|_2&=u\int(b,x)\Big\{-\tfrac{22}3(A,A^*)+16(\lambda,\lambda^*)+\tfrac83(c,c^*)\Big\}\\
&\quad+\tfrac23u\int\big\{(\iota_bA)(\iota_xA^*)-(\iota_xA)(\iota_bA^*)\big\}-\tfrac13u\int b_\mu x_\nu(\lambda,\rho_\mu\rho_\nu\lambda^*)+u\sum_d\theta^K_dS(K_b)_d+\sum_{j=1}^4\kappa^K_jG^K_j,
\end{aligned}
\]
with $Q$ of the first bracket equal to $\hbar\Ob^{(1)}(\Delta)|_2$ (gate G-F1Delta), $G^K_1,\dots,G^K_4$ the four second-order free-field symmetries of weight $3$ (closed, not exact, not of the form $S(z)$ on two legs; listed in action/complete/results.json, F1(K\_b)|2, free\_B), and the $x$-linear part of $\hbar F_1^{(1)}(K_b)|_2$ equal to $-(b,x)$ times the particular part along $\Delta$ plus the spin term of the rotation part of $d_{b'}v$ ($16-\tfrac13=\tfrac{47}3$). On two legs $H^0(Q)$ vanishes along $\Delta$, $P_b$, $M_\omega$, $R_r$, $\Sigma$, and is $1$-, $4$-dimensional along $\Sigma^*$ (position degree $\le1$), $K_b$ (position degree $\le2$); in both cases the residual has no component along it.
\end{proposition}

\begin{proposition}[Two-leg arity-two equations]
\label{prop:at-c-ar2}
On two legs, with the solutions of Proposition \ref{prop:at-c-sol}, the classes of the arity-two residuals vanish exactly under the following conditions (both routes):
\[
\begin{array}{l|l}
\text{pair} & \text{condition}\\
\hline
(\Delta,P_b) & \theta^\Delta_\perp=\theta^\Delta_\lambda=\theta^\Delta_V=-\tfrac13\quad(o_d=-\tfrac23u)\\
(\Delta,\xi_\pm) & \theta^\Delta_\perp=\theta^\Delta_\lambda=\theta^\Delta_V=-\tfrac13,\ \theta^\xi\ \text{free}\\
(\xi_+,\xi_-) & \theta^P_d=\theta^{\xi_+}_V+\theta^{\xi_-}_V-\tfrac{27}4-2L\ (\text{all }d),\quad \theta^{\xi_+}_V+\theta^{\xi_-}_V=\theta^{\xi_+}_\perp+\theta^{\xi_-}_\perp\\
(\xi_\pm,\xi'_\pm) & \theta^{\xi_\pm}_V=\theta^{\xi_\pm}_\perp\\
(\xi_\pm,P_b) & \theta^P_V=\theta^P_\perp=\theta^P_\lambda\\
(\xi_+,\eta_-) & \theta^\xi_V-\theta^\xi_\perp=\theta^\eta_V-\theta^\eta_\perp\\
(\xi_+,\eta_+) & \theta^R_\perp=\theta^R_\lambda=:\tau,\ \theta^\Delta_d=\theta^M_d=\tau-\tfrac13,\ \theta^\xi_d+\theta^\eta_d=\tau+\tfrac{77}{12}+2L\ (d=V,\perp)
\end{array}
\]
($\theta^\eta_d$ in the parametrisation of Proposition \ref{prop:at-c-sol}; in the parametrisation of action/complete/notes.txt 2.4, $\theta^\eta_\perp$ there $=\theta^\eta_\perp-3L+\tfrac32$ here.)
The obstructions: $\hbar\Ob^{(2)}(\Delta,\xi)|_2=u\big[8(T_1+T_2)-\tfrac{26}3(T_3+T_4)+\log\tfrac43\,Q(\psi_1+\psi_2)\big]$, $\psi_1=\int x^\mu(\Gamma(\xi,\lambda)^\mu,c^*)$, $\psi_2=\int x^\mu(A^*_\mu,(\xi,\lambda^*))$, solved by $\hbar F_1^{(2)}(\Delta,\xi)|_2=\log\tfrac43\,u(\psi_1+\psi_2)$ plus closed directions; $\hbar\Ob^{(2)}(\xi,P_b)|_2$ is $\tfrac14\log\tfrac43$ times a $Q$-exact functional; $\hbar\Ob^{(2)}(\Delta,P_b)|_2$ has no position part and no pole; every two-leg arity-two pole sum vanishes (along $(\xi,\eta)$ checked on the intrinsic part). With $\theta^\xi_V=\theta^\xi_\perp=:t_\xi$ the rows $(\xi_+,\eta_\pm)$ give $\theta^\eta_V=\theta^\eta_\perp=:t_\eta$ (the position part of $F_1^{(1)}(\eta)$ is then $x^\mu F_1^{(1)}(\rho(e_\mu)\eta)$ with one uniform coefficient) and uniform $t_P$, $t_\Delta$, $t_M$, $t_R$:
\[
t_\Delta=-\tfrac13,\qquad t_M=t_\Delta,\qquad t_R=t_\Delta+\tfrac13,\qquad t_{\xi_+}+t_{\xi_-}-t_P=\tfrac{27}4+2L,\qquad t_{\xi_+}+t_{\eta_+}-t_\Delta=\tfrac{27}4+2L .
\]
\end{proposition}

\begin{proof}
Computed: the cycles of Theorem \ref{thm:at-classes} and Proposition \ref{prop:at-classes2} in the two routes, the cohomology on complete bases fixed by representation counts (for $(\xi,\eta)$ decided on the value vectors directly, the generated $(\lambda,\lambda^*)$ position basis lacking one structure per chirality), and the residual formulas of Definition \ref{def:at-ob}. The arity-two generating function for two odd generators is $I+cS(\xi)+c'S(\xi')-cc'S^{(2)}(\xi,\xi')$, the sign forced by the classical master equation (gate GA-S2/GB-S2: $QS^{(2)}_2+S_2(P_{\Gamma|_V})-\{S_2(\xi),S_2(\xi')\}=0$). The last line: insert $\theta^\Delta=-\tfrac13$ into the $(\xi_+,\eta_+)$ row, giving $\tau=0$.
\end{proof}

\begin{corollary}[Scheme-independent content on the action side]
\label{cor:at-c-inv}
The inner direction ($t_z\mapsto t_z+c\lambda_z$, $\lambda=-1,+1,-2,+2,0$ on $\Sigma,\Sigma^*,P,K$ and the even rotations and $\Delta$) and the outer direction ($t_{\xi_\pm}\mapsto t_{\xi_\pm}\pm a$, $t_{\eta_\pm}\mapsto t_{\eta_\pm}\mp a$) leave invariant exactly the combinations
\[
t_\Delta=-\tfrac13,\qquad t_M-t_\Delta=0,\qquad t_R-t_\Delta=\tfrac13,\qquad t_{\xi_+}+t_{\xi_-}-t_P=\tfrac{27}4+2L,\qquad t_{\xi_+}+t_{\eta_+}-t_\Delta=\tfrac{27}4+2L,
\]
and the sector uniformity of $t_P$, $t_\xi$, $t_\eta$; these are the one-loop parts of $[D,P]=-2P$, $[D,Q]=-Q$, $[Q,P]=0$, $[Q,Q]=P$ and $\{Q,S\}=D+M+R$ in the form of Theorem \ref{thm:at-algebra}. In a translation-covariant scheme $t_P=t_\Delta=-\tfrac13$, hence $t_{\xi_+}+t_{\xi_-}=\tfrac{77}{12}+2L$ and $t_{\eta_+}=t_{\xi_-}$. Scheme artefacts: the individual $t_{\xi_\pm}$, $t_{\eta_\pm}$, $t_P$, $t_K$, $\theta^K_d$, $\kappa^K_j$, $\kappa_V$, $\kappa_\lambda$, the particular parts with $\log\tfrac32$, and every representative above.
\end{corollary}

\begin{remark}[Two-component forms]
\label{rem:at-c-tc}
With Remark \ref{rem:at-twocomp}: $\int(A,A^*)=\int d^4x\{(A_m,A^*_m)+\tfrac14(A_{ij},A^{*ij})\}$, $\int(\lambda^*,\lambda)=\int(\lambda,\lambda^*)=\int d^4x\{(\lambda_\alpha{}^i,\lambda^{*\alpha}{}_i)-(\bar\lambda_{\dot\alpha i},\bar\lambda^{*\dot\alpha i})\}$, $\int(\iota_bA)(\iota_xA^*)=\int d^4x\,(b_mA_m,x_nA^*_n)$, and $S^{(2)}(\xi,\xi')=\int\{\tfrac14(\Gamma(\xi,\xi'),\Gamma(\lambda^*,\lambda^*))-\tfrac12(\xi,\lambda^*)(\xi',\lambda^*)-(\iota_{\Gamma(\xi,\xi')}A,c^*)\}$ with $\Gamma(\xi,\bar\xi')|_V^m=i\,\bar\xi'_{\dot\alpha i}\bar\sigma^{m\dot\alpha\alpha}\xi_\alpha{}^i$, $\Gamma|_\perp=0$ for $\xi\in S_+\otimes W$, $\bar\xi'\in S_-\otimes W^*$, and $\Gamma|_V=0$, $(\Gamma|_\perp)_{kl}=2\varepsilon_{jikl}\,\xi_\alpha{}^j\varepsilon^{\alpha\beta}\xi'_\beta{}^i$ for $\xi,\xi'\in S_+\otimes W$. On the Dirac column the $\rho_\mu\rho_\nu$ of the spin terms is $\gamma^\mu_W\gamma^\nu_W$ (Proposition c8-twocomp), on the undotted block $\sigma^\mu\bar\sigma^\nu$ and on the dotted block $\bar\sigma^\mu\sigma^\nu$.
\end{remark}

\section{Engines and gates}
\label{sec:at-engines}

The specification is \texttt{engine\_spec.json} in this folder. Two independent routes compute every item; neither reads the other's outputs before the comparison.
\begin{itemize}
\item Route A (special-edge cycles in simplex form): enumerates the irreducible defect cycles of Theorems \ref{thm:at-classes} and Proposition \ref{prop:at-classes2} sector by sector, evaluates $a_m$ and $a_0$ exactly by Lemma \ref{lem:at-special} and Proposition \ref{prop:at-schwinger}, decomposes on exact bases of the target sectors by evaluation on integer plane waves in the explicit realisations of $V_{10}$, $\Sigma$ over $\mathbb Q(i)$ and $\mathfrak g=\mathfrak{su}(2)$, $\mathfrak{su}(3)$, $\mathfrak{gl}(3)$, and solves the equations by exact linear algebra. It extends the cycle enumerator of the Part-C action engine (step2/engine\_a, s2\_*) and the defect-chain code of Part Q (step2/composite\_qsk/qsk/engine\_a, qa\_core, qa\_s1) from chains through two letters to closed cycles.
\item Route B (renormalised amplitudes and $Q$): computes $W^{(1)}_{\rm ren}(x)$ and $W^{(2)}_{\rm ren}(x,y)$ by momentum-space numerator algebra with exact Feynman-parameter integrals, applies $Q$, reads the local part, and checks the special-edge cycles numerically at finite $\varepsilon$. It extends step2/engine\_b (ia\_engine, delta\_engine) and step2/composite\_qsk/qsk/engine\_b (qb\_core, qb\_numeps, qb\_numT).
\end{itemize}
Gates (regressions and identities of the construction): G-sign (the component formulas of Lemma \ref{lem:at-sformulas} and $s^2=0$ on random test cochains); G-reg0 (the arity-zero two-leg obstruction $-\tfrac{27}8u\int(d^*dc,d^*\AV)$ and the cubic obstruction \eqref{eq:at-ob3}); G-reg1 (\eqref{eq:at-reg2}); G-pole (\eqref{eq:at-pole} and \eqref{eq:at-local2}); G-log (no $\log\varepsilon$ in any cycle; $s\,\lc=0$); G-logex (Remark \ref{rem:at-logexact}, the closed combination $\lc^{(1)}(x)-\{uN,S(x)\}$); G-R0 ($\Ob^{(1)}(R_r)=0$ cycle sum); G-P2 ($\Ob^{(1)}(P_b)|_2=0$); G-pos (Proposition \ref{prop:at-pos}); G-WZ1, G-WZ2 (Lemma \ref{lem:at-sformulas} (iii) on two and three legs, on the full functional, intrinsic and position parts, with $\Lambda^{(0)}G^{(1)}(x)=-\{G^{(0)},S(x)\}$ and the arity-two form of (ii)); G-T (finite $T$ along $K_b$: $T\partial_T=O(\varepsilon)$); G-AB (route A against route B, entry by entry); G-tDelta (the determinations of $t_\Delta$ agree); G-zero (the classes on the commuting pairs vanish); G-H (Hypothesis \ref{hyp:at-H} at three legs; distinct from the two-leg cohomology gate G-S1xi2 of Proposition \ref{prop:at-xi2}). Cross-checks: X-scal, X-theta, X-twist, X-KP (Remark \ref{rem:at-cross}). Regression values:
\begin{equation}
\hbar\Ob^{(0)}\big|_3=u\Big[\big(\tfrac{11}4+\tfrac{19}4\log\tfrac32\big)(G_1+G_3)+\tfrac{27}8G_2-9\log\tfrac32\,G_5\Big]\quad(\text{Theorem c9-ob3}).
\label{eq:at-ob3}
\end{equation}

\section{Two-component forms}
\label{sec:at-twocomp}

\begin{remark}[Two-component forms of the objects of this note]
\label{rem:at-twocomp}
With Definition c8-twocomp ($\sigma^m=(\sigma^1,\sigma^2,\sigma^3,i)$, $\bar\sigma^m=(-\sigma^1,-\sigma^2,-\sigma^3,i)$, $\varepsilon^{12}=\varepsilon^{\dot1\dot2}=+1$), for $\xi\in S_+\otimes W$ with components $\xi_\alpha{}^i$ ($\bar\xi=0$), $\eta\in S_+\otimes W^*$ with components $\eta^\alpha{}_i$ ($\bar\eta=0$), and the pairing of $\mathfrak g$ written as $(-,-)$ around the components:
\[
\begin{aligned}
&\Gamma_V^m=\Gamma(\xi,\lambda)|_V^m=i\,\bar\lambda_{\dot\alpha i}\bar\sigma^{m\dot\alpha\alpha}\xi_\alpha{}^i,\qquad (\Gamma_\perp)_{kl}=2\varepsilon_{jikl}\,\xi_\alpha{}^j\varepsilon^{\alpha\beta}\lambda_\beta{}^i,\qquad (\xi,\lambda^*)=\xi_\alpha{}^i\lambda^{*\alpha}{}_i,\\
&T_0=-\int d^4x\,\big(\partial_mA_m,\ \xi_\alpha{}^i\lambda^{*\alpha}{}_i\big),\qquad T_1=-i\int d^4x\,\big(\bar\lambda_{\dot\alpha i}\bar\sigma^{m\dot\alpha\alpha}\xi_\alpha{}^i,A^*_m\big),\qquad T_2=-\tfrac12\int d^4x\,\varepsilon_{jikl}\big(\xi_\alpha{}^j\varepsilon^{\alpha\beta}\lambda_\beta{}^i,A^{*kl}\big),\\
&T_3=-\tfrac14\int d^4x\,\big((\partial_mA_n-\partial_nA_m)\,\sigma^m_{\alpha\dot\alpha}\bar\sigma^{n\dot\alpha\beta}\xi_\beta{}^i,\ \lambda^{*\alpha}{}_i\big),\qquad T_4=-\tfrac i2\int d^4x\,\big(\varepsilon^{\alpha\beta}\xi_\beta{}^j\sigma^m_{\alpha\dot\alpha}\partial_mA_{ij},\ \bar\lambda^{*\dot\alpha i}\big),\\
&g_1=\int d^4x\,\big(\Gamma_V^n,\ -\partial^2A_n+\partial_n\partial_mA_m\big),\quad g_2=-\int d^4x\,\big(\Gamma_V^n,\partial_n\partial_mA_m\big),\quad g_3=-\tfrac14\int d^4x\,\big((\Gamma_\perp)_{kl},\partial^2A^{kl}\big),\quad g_4=-\int d^4x\,\big(\partial^2c,\ \xi_\alpha{}^i\lambda^{*\alpha}{}_i\big),\\
&s_0=\int d^4x\,\big(A^{kj}\varepsilon_{\alpha\beta}\eta^\beta{}_j,\ \lambda^{*\alpha}{}_k\big),\qquad \int(\rho(\AV)\eta,\lambda^*)=-i\int d^4x\,\big(A_m\eta^\alpha{}_i\sigma^m_{\alpha\dot\alpha},\ \bar\lambda^{*\dot\alpha i}\big),
\end{aligned}
\]
with $A^{ij}=\tfrac12\varepsilon^{ijkl}A_{kl}$, $d^*=-\iota_\partial$, $d^*d\AV=-\partial^2A_n+\partial_n\partial_mA_m$ and $d^*d$ on the scalars equal to $-\partial^2$. The sector parts are
\[
\begin{aligned}
&S(\Delta)_\perp=\tfrac14\int d^4x\,\big((2x_n\partial_n+2)A_{ij},A^{*ij}\big),\qquad S(\Delta)_V=\int d^4x\,\Big\{\big((2x_n\partial_n+2)A_m,A^*_m\big)-\big(2x_n\partial_nc,c^*\big)\Big\},\\
&S(\Delta)_\lambda=-\int d^4x\,\Big\{\big((2x_n\partial_n+3)\lambda_\alpha{}^i,\lambda^{*\alpha}{}_i\big)-\big((2x_n\partial_n+3)\bar\lambda_{\dot\alpha i},\bar\lambda^{*\dot\alpha i}\big)\Big\},
\end{aligned}
\]
and $S(P_b)_d$ the same with $2x_n\partial_n+2w$ replaced by $b_n\partial_n$. The field-counting functional is $N=2\int(A,A^*)-4\int(\lambda^*,\lambda)-\int(c^*,c)$ with $\int(A,A^*)=\int d^4x\{(A_m,A^*_m)+\tfrac14(A_{ij},A^{*ij})\}$ and $\int(\lambda^*,\lambda)=\int d^4x\{(\lambda_\alpha{}^i,\lambda^{*\alpha}{}_i)-(\bar\lambda_{\dot\alpha i},\bar\lambda^{*\dot\alpha i})\}$. The regression \eqref{eq:at-reg2} reads
\[
\hbar\Ob^{(1)}(\Delta)\big|_2=u\int d^4x\Big\{-\tfrac{11}3\Big[\tfrac12(F^{\rm lin}_{mn},F^{\rm lin}_{mn})+\tfrac14(\partial_mA^{ij},\partial_mA_{ij})\Big]+\tfrac{94}3\,i\big(\bar\lambda_{\dot\alpha i},\bar\sigma^{m\dot\alpha\alpha}\partial_m\lambda_\alpha{}^i\big)-\tfrac{14}3\big(\partial_mc,A^*_m\big)\Big\},
\]
$F^{\rm lin}_{mn}=\partial_mA_n-\partial_nA_m$, and $G_1,\dots,G_5$ of \eqref{eq:at-ob3} are given in Remark c9-twocomp3. On the Dirac column $(\lambda_\alpha{}^i,\bar\lambda^{\dot\alpha}{}_i)$, $\rho(v)=\gamma_W^\mu v_\mu$ gives Weinberg's form of each line (Proposition c8-twocomp).
\end{remark}

\section{Status of the statements}
\label{sec:at-status}

\begin{itemize}
\item Derived from the construction (proofs above): Lemmas \ref{lem:at-defect}, \ref{lem:at-sformulas} (signs revised 2026-11-04 for the bracket of Definition c1-bracket), \ref{lem:at-place}, \ref{lem:at-special}; Propositions \ref{prop:at-poles}, \ref{prop:at-log}, \ref{prop:at-composite}, \ref{prop:at-pos}, \ref{prop:at-T}, \ref{prop:at-xi2}, \ref{prop:at-classes2}; Theorems \ref{thm:at-classes}, \ref{thm:at-second}, \ref{thm:at-bound}, \ref{thm:at-eq1}, \ref{thm:at-cubic}, \ref{thm:at-eq2} ((iii) by the forcing argument, not written out symbol by symbol); Theorems \ref{thm:at-H}, \ref{thm:at-fix}, \ref{thm:at-theta}, \ref{thm:at-algebra} under Hypothesis \ref{hyp:at-H}.
\item Computed by two engine routes and checked by the referee: Propositions \ref{prop:at-c-ob1}, \ref{prop:at-c-sol}, \ref{prop:at-c-ar2} and Corollary \ref{cor:at-c-inv} on two legs; the three-leg regressions \eqref{eq:at-ob3}; $\Ob^{(1)}(R_r)=0$ on three tested three-leg sectors.
\item Open (not computed): the pairs $(K_b,P_{b'})$, $(\eta,\eta')$, $(\xi_-,\eta_\pm)$; all three- and four-leg classes (Theorem \ref{thm:at-cubic}, task S1-cubic, Hypothesis \ref{hyp:at-H} at three legs, Proposition \ref{prop:at-couplings}); the four closed directions $G^K_j$ and the $\mathfrak{su}(4)$-$\mathbf{15}$ free classes; G-log as a finite-$\varepsilon$ extrapolation; G-T as a numerical finite-$T$ run; the five-leg sector along $P_b$.
\item Hypothesis: Hypothesis \ref{hyp:at-H} (classical symmetries at the listed weights), tested at three legs by G-H (not run).
\item Former shortcuts, now theorems from the construction or cross-checks: $\theta=-\tfrac13$ is computed from the construction twice ($(\Delta,P_b)$ and $(\Delta,\xi_\pm)$, Proposition \ref{prop:at-c-ar2}); the operator-level derivation (Remark \ref{rem:at-thetacheck}) is a cross-check only. The scaling identity ``$4$ times the logarithmic coefficient'' along $\Delta$ is Theorem \ref{thm:at-second} (v), used as the cross-check X-scal (it holds placement class by placement class). The twist (Remark \ref{rem:at-twist}) and $K=P^\dagger$ are cross-checks only and enter no number.
\end{itemize}

\end{document}
